% ============================================================
% ヘッジファンドの投資手法とアルゴリズム：基本と研究の整理
% コンパイル：tectonic -X compile ヘッジファンドの投資手法.tex   （XeTeX 系。このフォルダで実行する）
% 章は 章/。第 II 部は Gemini Deep Research の 3 本の調査（../../03_方針と文献調査/ヘッジファンドの手法/）を照合して書いた。
% ============================================================
\documentclass[xelatex,ja=standard,jafont=haranoaji,a4paper,10pt,openany]{bxjsreport}
\setpagelayout*{left=22mm,right=22mm,top=22mm,bottom=24mm,footskip=10mm}

\usepackage{amsmath,amssymb,mathtools,bm}
\usepackage{booktabs,array,multirow,longtable,tabularx}
\usepackage{graphicx}
\usepackage{xcolor}
\usepackage{tikz}
\usetikzlibrary{arrows.meta,positioning,calc,fit}
\usepackage{enumitem}
\usepackage{caption}
\usepackage{float}
\usepackage[most]{tcolorbox}
\usepackage[hidelinks,bookmarksnumbered]{hyperref}
\usepackage{xurl}
\pdfstringdefDisableCommands{\def\Omega{Ω}\def\Sigma{Σ}\def\alpha{α}\def\beta{β}}
\xeCJKDeclareCharClass{CJK}{"0391->"03C9, "2248, "2264->"2267, "2460->"24FF, "25A0->"25FF}

\definecolor{c1}{HTML}{2A78D6}
\definecolor{c2}{HTML}{EB6834}
\definecolor{c3}{HTML}{1BAF7A}
\definecolor{c4}{HTML}{EDA100}
\definecolor{c7}{HTML}{6250D6}
\definecolor{ink2}{HTML}{52514E}
\definecolor{soft1}{HTML}{E4EEF9}
\definecolor{soft2}{HTML}{FBEEDF}
\definecolor{soft3}{HTML}{E3F3E9}
\definecolor{soft4}{HTML}{EEE8FA}
\definecolor{soft5}{HTML}{FDF3DA}
\definecolor{soft6}{HTML}{F1F0EC}

\newcommand{\code}[1]{\texttt{#1}}
\newcommand{\SP}{S\&P500}
\newcommand{\T}{^{\top}}
\newcommand{\E}{\mathrm{E}}
\newcommand{\Var}{\mathrm{Var}}
\newcommand{\Cov}{\mathrm{Cov}}
\newcommand{\IC}{\mathrm{IC}}
\newcommand{\IR}{\mathrm{IR}}
\newcommand{\TE}{\mathrm{TE}}
\newcommand{\SR}{\mathrm{SR}}
\newcommand{\figref}[1]{図~\ref{#1}}
\newcommand{\tabref}[1]{表~\ref{#1}}
\newcommand{\chapref}[1]{第~\ref{#1}~章}
\newcommand{\term}[1]{\textsf{#1}}
\newcolumntype{P}[1]{>{\raggedright\arraybackslash}p{#1}}

\setlist[itemize]{leftmargin=1.6em,itemsep=0.2em,topsep=0.4em}
\setlist[enumerate]{leftmargin=1.8em,itemsep=0.2em,topsep=0.4em}
\captionsetup{font=small,labelfont=bf,skip=6pt}
\renewcommand{\arraystretch}{1.2}
\setlength{\LTpre}{6pt}
\setlength{\LTpost}{6pt}

\makeatletter
\def\@makechapterhead#1{\vspace*{0.3\Cvs}{\parindent\z@\raggedright\normalfont
  \ifnum\c@secnumdepth>\m@ne\if@mainmatter{\Large\headfont\@chapapp\thechapter\@chappos}\par\nobreak\vskip0.4\Cvs\fi\fi
  \interlinepenalty\@M\LARGE\headfont #1\par\nobreak\vskip 1.2\Cvs}}
\def\@makeschapterhead#1{\vspace*{0.3\Cvs}{\parindent\z@\raggedright\normalfont
  \interlinepenalty\@M\LARGE\headfont #1\par\nobreak\vskip 1.2\Cvs}}
\makeatother

\newtcolorbox{keybox}[1][要点]{enhanced,breakable,colback=soft1,colframe=c1,boxrule=0.5pt,arc=2pt,
  left=8pt,right=8pt,top=6pt,bottom=6pt,title={#1},fonttitle=\sffamily\bfseries\small,coltitle=white}
\newtcolorbox{notebox}[1][注意]{enhanced,breakable,colback=soft2,colframe=c2,boxrule=0.5pt,arc=2pt,
  left=8pt,right=8pt,top=6pt,bottom=6pt,title={#1},fonttitle=\sffamily\bfseries\small,coltitle=white}
\newtcolorbox{exbox}[1][例]{enhanced,breakable,colback=soft3,colframe=c3,boxrule=0.5pt,arc=2pt,
  left=8pt,right=8pt,top=6pt,bottom=6pt,title={#1},fonttitle=\sffamily\bfseries\small,coltitle=white}
\newtcolorbox{defbox}[1][定義]{enhanced,breakable,colback=soft6,colframe=ink2,boxrule=0.4pt,arc=2pt,
  left=8pt,right=8pt,top=6pt,bottom=6pt,title={#1},fonttitle=\sffamily\bfseries\small,coltitle=white}
\newtcolorbox{paperbox}[1][論文]{enhanced,breakable,colback=white,colframe=c7,boxrule=0.4pt,arc=2pt,
  left=8pt,right=8pt,top=6pt,bottom=6pt,title={#1},fonttitle=\sffamily\bfseries\small,coltitle=white}
\newtcolorbox{projbox}[1][この研究との関係]{enhanced,breakable,colback=soft5,colframe=c4!85!black,boxrule=0.4pt,arc=2pt,
  left=8pt,right=8pt,top=6pt,bottom=6pt,title={#1},fonttitle=\sffamily\bfseries\small,coltitle=black}

\tikzset{
  box/.style={draw=#1,fill=#1!8,rounded corners=2pt,align=center,font=\footnotesize,inner sep=4pt,
              minimum height=10mm,text width=24mm},
  box/.default=c1,
  arr/.style={-{Stealth[length=2mm]},color=ink2,line width=0.6pt},
  lab/.style={font=\scriptsize,color=ink2},
}

\title{ヘッジファンドの投資手法とアルゴリズム\\[8pt]
  \Large 基本の考え方と、研究で分かっていること}
\author{LLM-BLM の解説}
\date{2026 年 10 月}

\begin{document}
\maketitle
\setcounter{tocdepth}{1}
{\small\tableofcontents}

\chapter*{この報告書について}
\addcontentsline{toc}{chapter}{この報告書について}

\section*{目的}

LLM-BLM の研究では、言語モデルが決算発表文を読んだ点数を Black--Litterman モデルの見通しに入れて、ポートフォリオを作った。読み取りは発表をはさむ 2 日の株価の反応を強く説明したが、その情報は発表の当日にほぼ株価に入りきり、取引費用の後で時価総額加重を上回ることはできなかった。

では、プロのヘッジファンドは何をして収益を上げているのか。どんな戦略があり、どんなアルゴリズムを使い、その根拠になる研究は何を示していて、どこで失敗してきたのか。この報告書は、それを体系的にまとめたものである。読者には、金融の予備知識のない理系の大学生（線形代数・確率・統計・最適化の基本は知っている）を想定した。

\section*{構成}

\begin{description}[leftmargin=2em,style=nextline]
\item[第 I 部　基本（第 1〜11 章）]
  ヘッジファンドの仕組み、戦略の分類、成績の測り方、ファクター、シグナル、予測モデル、ポートフォリオ構築、リスク管理、取引コストと執行、バックテストと過学習、マーケットメイク。教科書や古典的な論文で確立していて、新しく調べなくても書ける標準的な内容である。式と小さな数値例で、仕組みが手で追えるように書いた。
\item[第 II 部　研究で分かっていること（第 12 章以降）]
  Gemini Deep Research に、切り口を変えた 3 つの依頼（戦略の種類から、作り方の工程から、成績の実証と失敗から）で文献を調べさせ、挙がった文献を照合したうえで、戦略ごと・工程ごとに整理した。第 I 部で説明した道具が、実際のデータでどれだけ効いたか、どこで崩れたかを扱う。最後の章で、この研究（LLM-BLM）への含みをまとめる。
\item[付録]
  用語の一覧と、文献の照合の結果。
\end{description}

\section*{文献の扱い}

\begin{itemize}
\item 第 I 部の引用は、分野の古典としてよく知られた論文と教科書に限った。全体の参考になる教科書として、Grinold・Kahn \cite{grinoldkahn2000}、Pedersen \cite{pedersen2015}、Ilmanen \cite{ilmanen2011}、López de Prado \cite{lopezdeprado2018}、Bouchaud ほか \cite{bouchaud2018} がある。
\item 第 II 部の文献は、Deep Research の 3 本の答え（\code{03\_方針と文献調査/ヘッジファンドの手法/}）に挙がったものと、それを補うために自分で選んだものである。報告書で引くすべての文献を、\code{refcheck.py} で題名から arXiv・OpenAlex・Crossref と照合し、実在が確かめられたものだけを使った。Deep Research の答えの誤りや、確かめられなかった記述は、付録~\ref{ch:verify} に一覧にした。
ef{ch:verify} に一覧にした。
\item 論文の数値は、その論文の標本（市場・期間・手数料や費用の前後）のものである。別の期間や市場で同じ値が出るとは限らない。本文では、できるだけ標本の範囲を一緒に書いた。
\item Deep Research の答えをそのまま写してはいない。数値は、照合で原典の要旨や表に当たれたもの、または分野でよく知られたものだけを書き、それ以外は「〜と報告されている」と出典を示して書くか、使わなかった。
\end{itemize}

\begin{notebox}
この報告書は勉強のための整理であり、投資の助言ではない。ヘッジファンドの成績のデータには、第 II 部で述べるように大きな偏りがある。論文で強かった戦略が、これからも同じように儲かるとは限らない。
\end{notebox}


\part{基本}
\chapter{ヘッジファンドとは}
\label{ch:hf}

\section{定義}

ヘッジファンド（hedge fund）には、法律で決まった定義はない。ふつうは、次の特徴をいくつか持つ運用の器（ファンド）をそう呼ぶ。

\begin{defbox}[ヘッジファンドの典型的な特徴]
\begin{enumerate}
\item \textbf{私募}：広く一般に売らず、富裕層・年金・大学基金・保険会社などの限られた投資家から資金を集める。
\item \textbf{手段の自由}：空売り、借入れ（レバレッジ）、先物・オプション・スワップなどのデリバティブを使える。投資信託（ミューチュアルファンド）は、これらを法律や約款で強く制限されている。
\item \textbf{絶対リターンの目標}：市場の指数に勝つこと（相対リターン）ではなく、市場が上がっても下がってもプラスの収益を上げること（絶対リターン）を目標に掲げることが多い。
\item \textbf{成功報酬}：運用の利益の一部（伝統的には 20\%）を運用者が受け取る。
\item \textbf{解約の制限}：毎日は解約できず、月・四半期ごとに、事前の通知をして解約する。最初の 1 年は解約できない（ロックアップ）こともある。
\end{enumerate}
\end{defbox}

「ヘッジ」は、もともと「市場全体の上げ下げのリスクを、空売りで打ち消す」という意味である。最初のヘッジファンドとされる A.~W.~Jones の運用（1949 年）は、割安と考える株を買い、割高と考える株を空売りして、市場の動きに左右されにくくしたうえで、借入れで規模を大きくした。いまの「ヘッジファンド」は、市場のリスクを打ち消さない戦略（たとえば通貨や金利の方向に大きく賭けるグローバル・マクロ）も含む、広い言葉になっている。

\section{歴史の流れ}

\begin{table}[H]
\centering\small
\caption{ヘッジファンドの歴史の主な出来事}
\label{tab:history}
\begin{tabular}{@{}lP{0.74\linewidth}@{}}
\toprule
年 & 出来事 \\
\midrule
1949 & A.~W.~Jones が、ロング・ショートとレバレッジを組み合わせた投資組合を始める。利益の 20\% を報酬とした。 \\
1966 & 経済誌 \textit{Fortune} の記事が Jones の成績を紹介し、同じ形のファンドが増える。 \\
1970 年代 & 1973〜74 年の下げ相場で多くのファンドが損をして閉じる。Bridgewater（1975）などが設立される。 \\
1980 年代 & グローバル・マクロの時代。Soros の Quantum Fund、Robertson の Tiger など。Renaissance Technologies（1982）、D.~E.~Shaw（1988）など、数学・計算機を使う運用会社が生まれる。 \\
1992 & Soros のファンドが英ポンドの売りで大きな利益を上げる（欧州の為替相場メカニズムからの英国の離脱）。 \\
1998 & LTCM（Long-Term Capital Management）が、ロシアの債務不履行をきっかけに巨額の損失を出し、ニューヨーク連銀の仲介で銀行団が救済する。 \\
2000 年代 & 年金・大学基金の資金が流入し、業界が急拡大する。AQR（1998）、Two Sigma（2001）などのクオンツ運用会社が成長する。 \\
2007 & 8 月に株式のクオンツ戦略が数日で急落する（quant quake）。 \\
2008 & 金融危機。Lehman Brothers の破綻でプライムブローカーに預けた資産が凍結され、多くのファンドが解約の制限（ゲート）をかける。 \\
2010 年代 & 手数料の引き下げの圧力、ファクター投資（スマートベータ）の普及、複数の運用チームを抱えるマルチマネージャーの拡大。 \\
2020 年代 & 2020 年 3 月の急落と回復、2021 年の個人投資家による空売りの踏み上げ（GameStop）、代替データと機械学習の普及、LLM の登場。 \\
\bottomrule
\end{tabular}
\end{table}

業界の規模は、データ会社や業界団体の推計で、2020 年代半ばに運用資産 4 兆ドル前後とされる（推計の方法によって幅がある）。ただし、その資産の多くは大手の数十社に集中している。

\section{仕組み}

\subsection{法的な形}

米国のヘッジファンドは、たいてい次の形をとる（\figref{fig:structure}）。

\begin{itemize}
\item \textbf{リミテッド・パートナーシップ}：運用会社（または関連会社）が無限責任のジェネラル・パートナー（GP）、投資家が有限責任のリミテッド・パートナー（LP）になる。
\item \textbf{投資会社法の適用の除外}：米国の 1940 年投資会社法の規制を避けるため、投資家の数を 100 人以下に限るか（3(c)(1)）、一定の資産を持つ「適格購入者」に限る（3(c)(7)）。広告をせずに私募で売る。
\item \textbf{マスター・フィーダー}：米国の課税される投資家のための米国内のファンド（フィーダー）と、米国外の投資家や非課税の投資家のための国外（ケイマン諸島など）のファンド（フィーダー）が、同じマスター・ファンドに資金を入れ、マスターがまとめて売買する。
\item \textbf{運用会社の登録}：2010 年の Dodd--Frank 法の後、一定の規模を超える運用会社は SEC に投資顧問として登録し、ファンドの規模やレバレッジを報告する（Form PF）。先物を主に使う運用会社（CTA）は CFTC の規制も受ける。
\end{itemize}

日本では、金融商品取引法の適格機関投資家等特例業務などの枠組みで、私募のファンドを運用する。

\begin{figure}[H]
\centering
\begin{tikzpicture}[node distance=7mm and 9mm]
\node[box=c4] (inv1) {米国の課税\\投資家};
\node[box=c4,right=of inv1] (inv2) {米国外・非課税\\の投資家};
\node[box=c1,below=of inv1] (f1) {米国内\\フィーダー（LP）};
\node[box=c1,below=of inv2] (f2) {国外フィーダー\\（ケイマン等）};
\node[box=c7,below=11mm of $(f1)!0.5!(f2)$,text width=30mm] (m) {マスター・ファンド\\（売買をまとめる）};
\node[box=c3,left=24mm of m] (gp) {運用会社\\（GP・投資顧問）};
\node[box=c2,right=24mm of m] (pb) {プライム\\ブローカー};
\node[box=ink2,below=8mm of m,text width=30mm] (mk) {市場（株・債券・\\先物・デリバティブ）};
\draw[arr] (inv1)--(f1);
\draw[arr] (inv2)--(f2);
\draw[arr] (f1)--(m);
\draw[arr] (f2)--(m);
\draw[arr] (gp)--node[lab,above]{運用の指示}(m);
\draw[arr,<->] (m)--node[lab,above]{資金・株の貸借}(pb);
\draw[arr,<->] (m)--(mk);
\end{tikzpicture}
\caption{マスター・フィーダーの構造。ほかに、基準価額を計算する管理会社（アドミニストレーター）、監査法人、法律事務所がかかわる。}
\label{fig:structure}
\end{figure}

\subsection{プライムブローカー}

ヘッジファンドの売買を支えるのが、投資銀行のプライムブローカー部門である。

\begin{itemize}
\item \textbf{資金の貸付け}：買う株を担保に資金を貸す（信用買い）。
\item \textbf{株の貸付け}：空売りのための株を探して貸す（ロケート）。借りる株には貸株料がかかり、借り手が多い株ほど高い。貸し手が株を返せと言えば（リコール）、空売りを閉じなければならない。
\item \textbf{保管・決済}：売買の決済と、資産の保管。
\item \textbf{証拠金の計算}：ポジションのリスクから必要な担保（証拠金）を計算し、損失が出れば追加の担保を求める（マージンコール）。
\end{itemize}

2008 年の Lehman Brothers の破綻では、その英国法人にプライムブローカーとして預けていた資産が長く凍結された。この後、多くのファンドが複数のプライムブローカーを使うようになった。

\subsection{解約の条件}

ヘッジファンドの資産には、すぐに売れないもの（小型株、社債、私募の資産など）が含まれる。投資家が一斉に解約すると、安値で投げ売りしなければならない。それを防ぐため、次のような条件をつける。

\begin{itemize}
\item \textbf{解約の頻度と通知の期間}：例えば「四半期末に、45 日前までに通知」。
\item \textbf{ロックアップ}：最初の 1 年などは解約できない。解約に手数料をとる「ソフト・ロックアップ」もある。
\item \textbf{ゲート}：1 回の解約で出せる額を、ファンド全体の一定の割合までに制限する。
\item \textbf{サイドポケット}：売れない資産を別の勘定に分け、売れるまで解約の対象から外す。
\end{itemize}

資産の流動性と、投資家の解約の条件が合っていないと、危機のときに投げ売りが起きる。これは第 II 部で扱う失敗の多くに共通する構造である。

\section{手数料}
\label{sec:fees}

\subsection{2 と 20}

伝統的な手数料は「2 と 20」と呼ばれる。

\begin{itemize}
\item \textbf{管理報酬}（management fee）：運用資産の年 2\%。成績にかかわらずかかる。
\item \textbf{成功報酬}（performance fee、incentive fee）：利益の 20\%。
\end{itemize}

実際の平均はこれより低く、2020 年代には管理報酬が 1\% 台半ば、成功報酬が 10\% 台後半と報じられることが多い。一方で、評判の高い一部のファンドは、もっと高い報酬をとる。複数の運用チームを抱える大手のマルチマネージャーでは、運用チームの報酬・データ・システムなどの費用をそのまま投資家に回す「パススルー」の形が広がり、投資家の負担はむしろ大きくなっている。

\subsection{ハイウォーターマークとハードル}

成功報酬には、次の決まりがつくことが多い。

\begin{itemize}
\item \textbf{ハイウォーターマーク}：基準価額が過去の最高値を超えた分にだけ成功報酬をとる。損をした翌年に取り返しただけでは、成功報酬はかからない。
\item \textbf{ハードル}：短期金利などの一定の収益を超えた分にだけ成功報酬をとる。
\end{itemize}

\begin{exbox}[手数料の計算]
投資家が 100 を預け、手数料の前の収益が年 10\% だったとする（短期金利 0、ハードルなし、ハイウォーターマークの下ではない）。
\begin{enumerate}
\item 管理報酬：$100 \times 2\% = 2$。
\item 成功報酬：手数料の前の利益 10 から管理報酬を引いた 8 の 20\%、つまり $1.6$。
\item 投資家の取り分：$10 - 2 - 1.6 = 6.4$、つまり 6.4\%。
\end{enumerate}
手数料の前の収益の 36\% が運用者に渡る。手数料の前の収益が 5\% なら、$5 - 2 - 0.6 = 2.4$ で、半分以上が運用者に渡る。手数料の前の収益が $-10\%$ なら、管理報酬で $-12\%$ になり、成功報酬はゼロである。
\end{exbox}

成功報酬は、運用者にとって「ファンドの利益に対するコールオプション」のような形をしている。損をしても運用者は損を負わず、利益が出れば一部を受け取る。そのため、リスクを取りすぎる誘因があると指摘されてきた。ハイウォーターマークの下に大きく沈んだファンドは、成功報酬を当分受け取れないので、閉じて別のファンドを立ち上げる誘因も生まれる。

\section{運用会社の型}

\subsection{シングル・マネージャーとマルチマネージャー}

\begin{itemize}
\item \textbf{シングル・マネージャー}：一人（または一つのチーム）の運用者の判断や、一つの戦略の体系で運用する。
\item \textbf{マルチマネージャー}（プラットフォーム型、ポッド型）：数十〜数百の小さな運用チーム（ポッド）に資金を割り当て、会社の中央でリスクを管理する。各チームには損失の上限があり、数\% の損失で割り当てを減らされ、さらに損をするとチームが解散させられる、と報じられる。市場全体やスタイル・ファクター（\chapref{ch:factor}）のリスクは打ち消すよう求められ、チームの「固有の」収益（アルファ）だけを積み上げる設計になっている。
\item \textbf{ファンド・オブ・ファンズ}：複数のヘッジファンドに投資するファンド。分散と選別を売りにするが、手数料が二重にかかる。
\end{itemize}

\subsection{裁量とシステマティック}

\begin{itemize}
\item \textbf{裁量（discretionary）}：運用者が、企業の調査、経済の見通し、経験に基づいて売買を判断する。
\item \textbf{システマティック（systematic、クオンツ）}：売買のルールを、データと統計・最適化のモデルで決め、計算機で実行する。人は、モデルを作り、検証し、監視する。
\end{itemize}

実際には、その中間（クオンツのスクリーニングで候補を絞り、人が判断する「クオンタメンタル」など）も多い。この報告書は、主にシステマティックな方法を扱う。

\subsection{ほかの運用との違い}

\begin{table}[H]
\centering\small
\caption{ヘッジファンドと、ほかの運用の違い（典型的な場合）}
\label{tab:types}
\begin{tabular}{@{}P{0.17\linewidth}P{0.18\linewidth}P{0.18\linewidth}P{0.18\linewidth}P{0.18\linewidth}@{}}
\toprule
 & 投資信託 & ヘッジファンド & 自己勘定の取引会社 & CTA（先物の運用） \\
\midrule
資金 & 一般の投資家 & 富裕層・機関投資家 & 会社自身 & 機関投資家・一般 \\
目標 & 指数に勝つ・連動 & 絶対リターン & 会社の利益 & 絶対リターン \\
空売り・借入れ & 強く制限 & 自由 & 自由 & 先物で自由 \\
報酬 & 資産の 0.1〜1\% 台 & 管理報酬＋成功報酬 & 社員の歩合 & 管理報酬＋成功報酬 \\
解約 & 毎日 & 月・四半期 & なし & 月・毎日 \\
典型的な戦略 & ロングのみの株・債券 & 多様 & 高頻度・マーケットメイク & トレンドフォロー \\
\bottomrule
\end{tabular}
\end{table}

\section{収益はどこから来るか}

ヘッジファンドの収益を、次の三つに分けて考えると見通しがよい（\chapref{ch:perf}、\chapref{ch:factor} で式にする）。

\begin{enumerate}
\item \textbf{市場のベータ}：株式市場全体が上がれば上がる部分。指数のファンドで安く手に入る。
\item \textbf{オルタナティブ・ベータ}（スタイル・プレミアム、リスク・プレミアム）：割安株、モメンタム、キャリー、トレンドフォロー、ボラティリティの売りなど、ルールで再現できる戦略の収益。昔はヘッジファンドの「腕」とされたが、研究が進み、いまは低い手数料の商品でも手に入る。
\item \textbf{アルファ}：上の二つで説明できない部分。情報、分析、速さ、執行の上手さ、ほかの人が取れないリスク（流動性の低い資産を持つ、混雑した取引の反対に立つ）から来る。
\end{enumerate}

研究で繰り返し問われてきたのは、「ヘッジファンドの手数料の後の収益のうち、本当のアルファはどれだけか」である。第 II 部（第 \ref{ch:hfperf} 章）で見るように、多くの研究は、平均的なファンドの収益の大部分が 1 と 2 で説明でき、手数料の後のアルファは小さいか、時代とともに小さくなったと報告している。

\begin{projbox}
LLM-BLM のポートフォリオは、ロングのみで、時価総額加重からの乖離（TE）に上限をつけた。これはヘッジファンドというより、指数に対して運用する伝統的な株式の運用会社（ロングのみのアクティブ運用）の設計に近い。空売りが使えないので、悪いと読んだ銘柄からは「持たない」分しか収益を取れない。この違いは、\chapref{ch:pc} と第 \ref{ch:implications} 章で再び扱う。
\end{projbox}

\chapter{戦略の分類}
\label{ch:strategies}

\section{業界の分類}

ヘッジファンドの成績のデータを集める会社は、ファンドを戦略で分類している。代表的な分け方は、次の四つの大分類である（データ会社 HFR の分類にほぼ沿う）。

\begin{table}[H]
\centering\small
\caption{戦略の大分類と主な小分類}
\label{tab:taxonomy}
\begin{tabular}{@{}P{0.17\linewidth}P{0.37\linewidth}P{0.38\linewidth}@{}}
\toprule
大分類 & 主な小分類 & 一言で \\
\midrule
エクイティ・ヘッジ & ロング・ショート株式、マーケット・ニュートラル、クオンツ株式、セクター特化、ショート・バイアス & 株の銘柄選択。買いと空売りを組み合わせる。 \\
イベント・ドリブン & 合併裁定、ディストレスト（破綻・再建）、アクティビスト、特別な状況（分割・自社株買い・指数の入れ替え） & 企業の出来事の結果に賭ける。 \\
マクロ & グローバル・マクロ（裁量）、システマティック・マクロ、マネージド・フューチャーズ（CTA）、通貨 & 金利・通貨・商品・株価指数の方向に賭ける。 \\
レラティブ・バリュー & 債券の裁定、転換社債の裁定、ボラティリティの裁定、ベーシス取引、統計的裁定 & 似た資産の価格の歪みをとる。方向のリスクは打ち消す。 \\
\bottomrule
\end{tabular}
\end{table}

ほかに、これらを組み合わせるマルチストラテジー、複数のファンドに投資するファンド・オブ・ファンズがある。

\section{各戦略の中身}

ここでは、それぞれの戦略が「何を買い、何を売り、どれだけ持ち、なぜ儲かると考えられているか」を短くまとめる。研究の結果は第 II 部で詳しく扱う。

\subsection{ロング・ショート株式}

割安・好業績と判断した株を買い、割高・業績悪化と判断した株を空売りする。買いと売りの金額の差（ネット）を、資産の 0〜60\% 程度のプラスにすることが多く、市場が上がれば利益が出やすい。判断は、企業の調査（裁量）か、ファクターと統計のモデル（クオンツ）による。

\subsection{マーケット・ニュートラル}

買いと売りの金額をそろえ（ドル・ニュートラル）、さらに市場に対する感応度 $\beta$ もゼロにする（ベータ・ニュートラル）。業種や規模・スタイルのファクターへの感応度も打ち消すことが多い。残るのは、銘柄の相対的な良し悪しの予測からの収益だけである。クオンツ株式の多くはこの形で、数百〜数千の銘柄を持ち、1 銘柄の重みは小さい。

\subsection{統計的裁定}

似た値動きをする銘柄の組（ペア）や、市場全体のファクターからの「ずれ」が、元に戻ることに賭ける。保有期間は数分〜数週間で、売買の回数が多い。ずれは、一時的な需給の偏り（大口の売買）から生まれることが多く、収益は「流動性を供給した対価」と解釈される。

\subsection{トレンドフォロー（マネージド・フューチャーズ）}

株価指数・債券・通貨・商品の先物を、最近上がっていれば買い、下がっていれば売る。数十〜百以上の先物市場に分散し、保有期間は数週間〜数か月。株式の大きな下げ相場（2008 年など）でも利益を上げたことがあり、「危機のときの保険」として持たれる。

\subsection{グローバル・マクロ}

経済の見通し（成長、インフレ、金融政策、政治）から、金利・通貨・株価指数・商品の方向に賭ける。裁量のものが多いが、指標とルールで決めるシステマティック・マクロもある。代表的なルールの一つが \textbf{キャリー}：金利の高い通貨を買い低い通貨を売る、利回り曲線の形から有利な年限を持つ、などである。

\subsection{イベント・ドリブン}

\begin{itemize}
\item \textbf{合併裁定}：買収の発表後、買収される会社の株価は、買収の価格より少し低い。買収が成立すればその差が利益になる。成立しないリスクを引き受ける対価である。
\item \textbf{ディストレスト}：破綻や再建の途中の会社の債券・株を安く買い、再建後の価値に賭ける。
\item \textbf{アクティビスト}：株を大量に買い、経営に提案（自社株買い、事業の売却、経営陣の交代）をして価値を上げる。
\item \textbf{決算・指数のイベント}：決算発表の後の株価の遅れた反応（決算発表後のドリフト）、指数の銘柄の入れ替えに伴う機械的な売買など。
\end{itemize}

\subsection{レラティブ・バリュー}

\begin{itemize}
\item \textbf{債券の裁定}：ほぼ同じ国債（発行されたばかりの銘柄と、少し前に発行された銘柄）の利回りの差、国債と金利スワップの差などをとる。差は小さいので、大きなレバレッジをかける。LTCM の主な戦略だった。
\item \textbf{転換社債の裁定}：株に転換できる社債を買い、その株を空売りして株価のリスクを打ち消す。転換社債が割安に発行されやすいことから収益を得る。
\item \textbf{ベーシス取引}：現物と先物の価格の差（ベーシス）をとる。米国債の現物を買って先物を売る取引が代表で、大きなレバレッジを使う。
\end{itemize}

\subsection{ボラティリティの取引}

オプションの価格に含まれる「予想の変動率」（インプライド・ボラティリティ）は、平均すると、後で実際に起きた変動率（実現ボラティリティ）より高い。そこでオプションを売り続けると、平均的には利益が出る（ボラティリティ・リスク・プレミアム）。ただし、急落のときに大きく損をする。保険を売る商売に似ている。

\subsection{マーケットメイク・高頻度取引}

売りと買いの両方の注文を常に出し、その価格の差（スプレッド）をとる。保有期間は秒以下〜数時間。収益の源は、スプレッドと取引所の報奨金で、損の源は、情報を持った相手との取引（逆選択）と、持った在庫の値下がりである。速さと注文の仕組みの理解が勝負で、多くは自己勘定の取引会社が行う（\chapref{ch:mm}）。

\section{収益の源泉で分ける}

戦略を、「なぜ儲かるのか」で分けると、研究の整理がしやすい。

\begin{table}[H]
\centering\small
\caption{収益の源泉の分類}
\label{tab:sources}
\begin{tabular}{@{}P{0.2\linewidth}P{0.42\linewidth}P{0.3\linewidth}@{}}
\toprule
源泉 & 中身 & 例 \\
\midrule
リスク・プレミアム & 投資家が嫌うリスク（景気の悪いときに損をする、急落する）を引き受ける対価。合理的な市場でも残る。 & 株式のリスク・プレミアム、キャリー、ボラティリティの売り、合併裁定 \\
行動のバイアス & 投資家の判断の偏り（反応の遅れ、過剰な反応、注目の偏り）から生じる価格の歪み。気づく人が増えると小さくなる。 & モメンタム、決算発表後のドリフト、割安株の一部 \\
流動性の供給 & 急いで売買したい人の反対に立つ対価。 & 統計的裁定、短期リバーサル、マーケットメイク \\
構造・制約 & 規制・約款・指数のルールなど、利益以外の理由での売買。 & 指数の入れ替え、レバレッジを使えない投資家の行動（低ベータ） \\
情報・分析 & ほかの人より早く・正しく情報を処理する。 & 企業調査、代替データ、テキストの分析 \\
速さ・技術 & 注文を速く出し、速く取り消す。 & 高頻度取引、遅延の裁定 \\
\bottomrule
\end{tabular}
\end{table}

この区別は、戦略の将来を考えるうえで重要である。リスク・プレミアムは、リスクを嫌う投資家がいる限り残ると期待できるが、急落の時期を耐えなければならない。行動のバイアスによる歪みは、研究が公表され、資金が集まると小さくなりやすい（第 II 部、第 \ref{ch:failures} 章）。

\section{保有期間と容量}

戦略は、保有期間と、運用できる資金の量（容量、キャパシティ）でも大きく違う（\figref{fig:horizon}）。

\begin{figure}[H]
\centering
\begin{tikzpicture}[x=13mm,y=10mm]
\draw[arr] (0,0)--(10.3,0) node[below left,font=\footnotesize]{保有期間};
\draw[arr] (0,0)--(0,5.4) node[above,font=\footnotesize]{容量（運用できる資金）};
\foreach \x/\l in {0.5/ミリ秒,2/秒〜分,3.6/日,5.2/週,6.8/月,8.4/年} \node[lab,below] at (\x,0) {\l};
\node[lab,left] at (0,0.6) {小};
\node[lab,left] at (0,4.8) {大};
\node[box=c2,text width=19mm,minimum height=7mm] at (0.9,0.8) {マーケット\\メイク・HFT};
\node[box=c1,text width=19mm,minimum height=7mm] at (3.1,1.7) {統計的裁定};
\node[box=c4,text width=19mm,minimum height=7mm] at (4.6,2.6) {イベント\\（決算・合併）};
\node[box=c7,text width=19mm,minimum height=7mm] at (6.3,3.3) {ファクター\\株式};
\node[box=c3,text width=19mm,minimum height=7mm] at (7.4,4.6) {トレンド\\フォロー};
\node[box=ink2,text width=19mm,minimum height=7mm] at (9.1,4.0) {マクロ・\\キャリー};
\end{tikzpicture}
\caption{戦略の保有期間と容量の大まかな関係（概念図。位置は目安）。短い戦略ほど 1 回の売買の利益が小さく、取引コストに敏感で、大きな資金を入れると利益が消える。}
\label{fig:horizon}
\end{figure}

短い戦略は、1 年に何千回も独立に近い賭けができるので、シャープレシオ（\chapref{ch:perf}）は高くなりやすい。しかし、1 回の利益が小さく、売買で価格を動かしてしまうので、大きな資金は入れられない。長い戦略は、大きな資金を入れられるが、賭けの回数が少なく、シャープレシオは低くなりやすい。高頻度の取引会社が外部の資金をあまり集めず、マクロや大型のクオンツ運用会社が巨額の資金を集めるのは、この構造による。

\section{クオンツ運用の工程}
\label{sec:pipeline}

システマティックな運用は、どの戦略でも、だいたい同じ工程で作られる（\figref{fig:pipeline}）。第 I 部の残りの章は、この工程の順に並んでいる。

\begin{figure}[H]
\centering
\begin{tikzpicture}[node distance=5mm]
\node[box=ink2,text width=20mm] (d) {データ\\{\scriptsize 価格・財務・テキスト}};
\node[box=c1,text width=20mm,right=of d] (s) {シグナル\\{\scriptsize 第 5 章}};
\node[box=c1,text width=20mm,right=of s] (p) {予測\\{\scriptsize 第 6 章}};
\node[box=c7,text width=20mm,right=of p] (pc) {ポートフォリオ\\{\scriptsize 第 7 章}};
\node[box=c2,text width=20mm,right=of pc] (ex) {執行\\{\scriptsize 第 9 章}};
\node[box=c4,text width=24mm,below=9mm of p] (r) {リスクモデル\\{\scriptsize 第 8 章}};
\node[box=c4,text width=24mm,below=9mm of pc] (c) {取引コストのモデル\\{\scriptsize 第 9 章}};
\node[box=c3,text width=60mm,below=9mm of $(r)!0.5!(c)$] (bt) {検証（バックテスト・過学習の補正）\\{\scriptsize 第 10 章}};
\draw[arr] (d)--(s); \draw[arr] (s)--(p); \draw[arr] (p)--(pc); \draw[arr] (pc)--(ex);
\draw[arr] (r)--(pc); \draw[arr] (c)--(pc);
\draw[arr,dashed] (bt.west) -| (d.south);
\draw[arr,dashed] (ex.south) |- (bt.east);
\end{tikzpicture}
\caption{クオンツ運用の工程。成績の測り方（第 3 章）とファクター（第 4 章）は、すべての工程の共通の言葉になる。}
\label{fig:pipeline}
\end{figure}

\begin{projbox}
LLM-BLM の工程を当てはめると、データ＝決算発表文と株価、シグナル＝LLM・FinBERT・辞書の点数、予測＝リッジ回帰で 2 日の反応に結びつける、ポートフォリオ＝Black--Litterman と TE の上限つきの二次計画、執行＝月末の終値で売買（片道 0.3\% の費用）である。研究で弱かったのは、シグナルの「持続」（情報が当日に株価に入りきる）と、予測の大きさに比べたコストだった。プロの工程との差は、第 II 部の最後（第 \ref{ch:implications} 章）でまとめる。
\end{projbox}

\chapter{成績の測り方}
\label{ch:perf}

戦略の良し悪しを比べるには、共通の物差しが要る。この章では、運用の世界で使われる成績の指標と、その落とし穴をまとめる。

\section{リターン}

\subsection{単純リターンと対数リターン}

時点 $t-1$ から $t$ までの資産の\textbf{単純リターン}は、価格 $P_t$ と配当 $D_t$ から
\begin{equation}
R_t = \frac{P_t + D_t}{P_{t-1}} - 1
\end{equation}
である。\textbf{対数リターン}は $r_t = \ln(1+R_t)$。対数リターンは期間をまたいで足し算できる（$r_{1\to T} = \sum_t r_t$）。一方、ポートフォリオのリターンは単純リターンの加重平均（$R_{p,t} = \sum_i w_i R_{i,t}$）で、対数リターンでは成り立たない。銘柄をまたぐ計算には単純リターン、時間をまたぐ計算には対数リターンが便利である。

\subsection{超過リターンと年率化}

\begin{itemize}
\item \textbf{超過リターン}：安全な資産（短期国債など）の金利 $R_f$ を引いた $R_t - R_f$。ヘッジファンドは現金を担保に先物や空売りを使うので、比べる相手は「現金」であることが多い。
\item \textbf{比べる相手との差}（アクティブ・リターン）：指数に対して運用するなら、$R_{p,t} - R_{b,t}$。
\item \textbf{年率化}：日次の平均 $\bar R$ と標準偏差 $\sigma$ から、年率の平均は $252\,\bar R$、年率の標準偏差は $\sqrt{252}\,\sigma$（月次なら 12 と $\sqrt{12}$）。標準偏差に平方根がつくのは、独立な変動の分散が期間に比例して足されるからである。
\end{itemize}

\subsection{算術平均と幾何平均}

年ごとのリターンの算術平均を $\mu$、標準偏差を $\sigma$ とすると、複利で増える速さ（幾何平均）$g$ はおよそ
\begin{equation}
g \approx \mu - \frac{\sigma^2}{2}
\end{equation}
である。同じ平均でも、変動が大きいほど資産は増えにくい（ボラティリティの引きずり）。例えば $+50\%$ と $-50\%$ を交互に繰り返すと、算術平均は 0 だが、資産は 2 年で $1.5 \times 0.5 = 0.75$ 倍になる。レバレッジを上げると $\mu$ は比例して増えるが、$\sigma^2$ は 2 乗で増えるので、ある倍率を超えると $g$ はかえって下がる（\chapref{ch:pc} のケリー基準）。

\section{シャープレシオ}

\subsection{定義}

最も広く使われる物差しが、Sharpe（1966）\cite{sharpe1966} の\textbf{シャープレシオ}である。
\begin{equation}
\SR = \frac{\E[R - R_f]}{\sigma(R - R_f)}
\end{equation}
リスク 1 単位あたりの超過リターンを表す。レバレッジを $k$ 倍にすると、分子も分母も $k$ 倍になるので、シャープレシオは変わらない。つまり、シャープレシオの高い戦略は、借入れで好きなリスクの水準に調整できる。ヘッジファンドがシャープレシオを重視するのは、このためである。

日次のシャープレシオ $\SR_d$ を年率にするには、リターンが独立なら $\SR = \sqrt{252}\,\SR_d$ とする。

\begin{table}[H]
\centering\small
\caption{シャープレシオの目安（年率、費用の後）}
\label{tab:sr}
\begin{tabular}{@{}lP{0.66\linewidth}@{}}
\toprule
年率の SR & 目安 \\
\midrule
0.3〜0.5 & 米国株式の長期の平均（市場全体をただ持つ） \\
0.5〜1.0 & よく分散したファクター戦略、トレンドフォローの長期の成績 \\
1.0〜2.0 & 優れたヘッジファンド。長く続くことは少ない \\
2 を超える & 容量の小さい短期の戦略、マーケットメイク。あるいはバックテストの過学習 \\
\bottomrule
\end{tabular}
\end{table}

\subsection{推定の誤差}

シャープレシオは推定値であり、誤差が大きい。リターンが独立で正規分布に従うとき、$T$ 期の標本から推定したシャープレシオ $\widehat{\SR}$（同じ期間の単位）の標準誤差は、Lo（2002）\cite{lo2002} によればおよそ
\begin{equation}
\mathrm{SE}(\widehat{\SR}) \approx \sqrt{\frac{1 + \SR^2/2}{T}}
\end{equation}
である。年率で考えると、$\SR$ が小さいとき、$t$ 値はおよそ $\SR \times \sqrt{\text{年数}}$ になる。

\begin{exbox}[腕を見分けるのに何年かかるか]
年率のシャープレシオが $\SR$ の戦略が、偶然でないと言えるには（$t \geq 2$）、およそ $(2/\SR)^2$ 年の記録が要る。
\begin{center}
\begin{tabular}{@{}cccccc@{}}
\toprule
年率の SR & 0.25 & 0.5 & 1.0 & 2.0 & 3.0 \\
\midrule
必要な年数 & 64 & 16 & 4 & 1 & 0.44 \\
\bottomrule
\end{tabular}
\end{center}
年率 0.5 の戦略（よいファクター戦略）が本物かどうかを、成績の記録だけで確かめるには 16 年かかる。しかも、その間に戦略の中身や市場が変わる。短期の戦略が好まれるもう一つの理由は、同じ年数で多くの賭けができ、早く検証できることである。
\end{exbox}

\subsection{シャープレシオの限界}

シャープレシオは、平均と標準偏差しか見ない。次のような戦略は、シャープレシオが実際より良く見える。

\begin{itemize}
\item \textbf{オプションの売り}：普段は小さな利益を積み重ね、まれに大きく損をする（負の歪度）。損の起きない期間を切り取ると、シャープレシオは非常に高い。
\item \textbf{評価の平滑化}：市場の価格がない資産を、運用者の評価や古い価格で値付けすると、リターンの変動が小さく見え、自己相関が出る（\ref{sec:smoothing} 節）。
\item \textbf{レバレッジの解消の連鎖}：同じ取引をしている人が一斉に売ると、普段の標準偏差からは想像できない損失が出る（第 II 部、第 \ref{ch:failures} 章）。
\end{itemize}

そこで、次の指標も一緒に見る。

\section{そのほかの指標}

\begin{description}[leftmargin=1.5em,style=nextline]
\item[ソルティノ・レシオ]
  分母を、下方の変動（目標を下回ったリターンだけの標準偏差）にしたもの。上方の変動を罰しない。
\item[最大ドローダウン]
  資産の過去の最高値からの最大の下落率。
  \begin{equation}
  \mathrm{MDD} = \max_t \frac{\max_{s \le t} V_s - V_t}{\max_{s \le t} V_s}
  \end{equation}
  投資家が実際に耐えられるかどうかを最もよく表す。期間が長いほど大きくなるので、同じ期間で比べる。
\item[カルマー・レシオ]
  年率のリターンを最大ドローダウンで割ったもの。
\item[歪度と尖度]
  リターンの分布の非対称と、裾の厚さ。負の歪度と厚い裾は、まれな大損を示す。
\item[VaR と期待ショートフォール]
  損失の分布の分位点と、その先の平均（\chapref{ch:risk}）。
\item[回転率]
  1 年に売買する額を資産で割ったもの。高いほど取引コストがかかる（\chapref{ch:cost}）。
\end{description}

\section{ベータ・アルファ・情報レシオ}

\subsection{回帰による分解}

戦略の超過リターンを、市場の超過リターンに回帰する。
\begin{equation}
R_{p,t} - R_{f,t} = \alpha + \beta\,(R_{m,t} - R_{f,t}) + \varepsilon_t
\end{equation}
\begin{itemize}
\item $\beta$：市場が 1\% 動いたときに戦略が動く大きさ。指数のファンドで安く手に入る部分。
\item $\alpha$：市場で説明できない平均の超過リターン。「腕」の候補。
\item $\varepsilon_t$：市場と無関係な変動。
\end{itemize}
市場だけでなく、規模・割安・モメンタムなどのファクター（\chapref{ch:factor}）も右辺に入れると、「ファクターでは説明できないアルファ」が測れる。ヘッジファンドの研究では、オプションに似たトレンドフォローのファクターなども使う（第 II 部、第 \ref{ch:hfperf} 章）。

\subsection{トラッキングエラーと情報レシオ}

指数に対して運用するとき、比べる相手との差の標準偏差を\textbf{トラッキングエラー}（TE）と呼ぶ。
\begin{equation}
\TE = \sigma(R_p - R_b), \qquad \IR = \frac{\E[R_p - R_b]}{\TE}
\end{equation}
$\IR$ は\textbf{情報レシオ}で、比べる相手との差についてのシャープレシオである。TE の上限は、「指数からどれだけ離れてよいか」の予算になる。LLM-BLM のポートフォリオで年 3\% の TE の上限を置いたのは、この考え方である。

\section{アクティブ運用の基本法則}
\label{sec:flam}

Grinold（1989）\cite{grinold1989} は、情報レシオを、予測の質と賭けの数に分解した。
\begin{equation}
\IR \approx \IC \times \sqrt{\mathrm{BR}}
\label{eq:flam}
\end{equation}
\begin{itemize}
\item $\IC$（information coefficient、情報係数）：予測と、実際の（リスクで割った）リターンの相関。0.05 なら「よい予測」、0.1 を超えることは珍しい。
\item $\mathrm{BR}$（breadth、幅）：1 年に行う独立な賭けの数。500 銘柄を毎月入れ替えるなら、最大で $500 \times 12 = 6{,}000$。
\end{itemize}

Clarke・de Silva・Thorley（2002）\cite{cdt2002} は、制約（ロングのみ、TE の上限など）で予測を比率に移しきれない分を、\textbf{移転係数} $\mathrm{TC}$（0〜1）として加えた。
\begin{equation}
\IR \approx \mathrm{TC} \times \IC \times \sqrt{\mathrm{BR}}
\end{equation}

\begin{exbox}[基本法則で考える]
\begin{itemize}
\item IC 0.05、500 銘柄、毎月独立に予測：$\IR = 0.05 \times \sqrt{6000} \approx 3.9$。
\item 同じで、ロングのみの制約で $\mathrm{TC} = 0.3$ だったとする（仮の値）：$\IR \approx 1.2$。
\item さらに、銘柄の予測が互いに相関していて、実質的に独立な賭けが 1 年に 600 回しかない：$\IR \approx 0.3 \times 0.05 \times \sqrt{600} \approx 0.37$。
\item 取引コストを引くと、さらに下がる。
\end{itemize}
この式は、弱い予測でも多くの銘柄・多くの時点に広げれば価値があることを示す一方、独立性・制約・コストで、その価値の大半が失われうることも示す。クオンツ運用が数千銘柄を扱い、制約を減らし（空売り、ロング・ショート）、コストを細かくモデル化するのは、この式の各項を大きくするためである。
\end{exbox}

\section{評価の平滑化と自己相関}
\label{sec:smoothing}

流動性の低い資産（社債、私募の資産、小型株）は、毎月の市場の価格が信頼できない。運用者が古い価格や推定で値付けすると、報告されるリターンは、本当のリターンを数か月にわたってならしたものになる。Getmansky・Lo・Makarov（2004）\cite{glm2004} は、報告されるリターン $R^o_t$ を
\begin{equation}
R^o_t = \theta_0 R_t + \theta_1 R_{t-1} + \cdots + \theta_k R_{t-k}, \qquad \textstyle\sum_j \theta_j = 1
\end{equation}
とモデル化した。ならすと、平均は変わらないが、標準偏差は小さくなり、シャープレシオは高く見え、リターンに正の自己相関が出る。

自己相関があるときの年率化は、$\sqrt{12}$ 倍では正しくない。月次の自己相関を $\rho_k$ とすると、$q$ か月の和の分散は
\begin{equation}
\Var\Big(\sum_{j=1}^{q} R_j\Big) = \sigma^2\Big(q + 2\sum_{k=1}^{q-1}(q-k)\rho_k\Big)
\end{equation}
で、正の自己相関があると、年率の変動は $\sqrt{12}\,\sigma$ より大きい。Lo（2002）は、この補正をしたシャープレシオを示した。

\section{手数料の前と後}

ヘッジファンドの成績は、手数料の前（グロス）と後（ネット）で大きく違う（\ref{sec:fees} 節の例）。論文や資料のシャープレシオが、手数料・取引コスト・借入れの費用・貸株料のどれを引いた後の値かを、必ず確かめる必要がある。研究の戦略の成績は、多くが手数料も取引コストも引く前の値である。

\begin{projbox}
LLM-BLM の月次の運用で、年率の TE を 3\% とし、情報レシオが 0.3 だったとすると、比べる相手との差の期待値は年 0.9\%。$t \geq 2$ に必要な年数は $(2/0.3)^2 \approx 44$ 年である。数年の検証期間で「有意に勝った」と言える戦略は、情報レシオが 1 前後以上でなければならない。この研究で時価総額加重との差が有意にならなかったことは、それだけでは「予測が無価値」を意味しない。ただし、2 日の反応で測った IC と、月次の運用で実現した IC の大きな差（情報が当日に株価に入りきる）は、検定力とは別の問題である。
\end{projbox}

\chapter{資産価格とファクター}
\label{ch:factor}

ヘッジファンドの戦略の多くは、「どの資産が、なぜ、平均して高いリターンを持つのか」という資産価格の理論の上に立っている。この章では、その理論と、それを測るための道具（ファクター・モデル）をまとめる。

\section{効率的市場の仮説}

Fama（1970）\cite{fama1970} は、市場の効率性を、価格に織り込まれている情報の範囲で三つに分けた。

\begin{itemize}
\item \textbf{ウィーク型}：過去の価格の情報は、すでに価格に入っている。過去の価格の動きから将来を予測しても儲からない。
\item \textbf{セミストロング型}：公開されたすべての情報（決算、ニュース）は、すでに価格に入っている。
\item \textbf{ストロング型}：公開されていない情報も、価格に入っている。
\end{itemize}

ただし、効率性の検定は、必ず「正しいリスクの調整の方法」の仮定と組みになる（結合仮説の問題）。ある戦略が平均して儲かっても、それが「市場の非効率」なのか「まだ測れていないリスクの対価」なのかは、データだけからは決まらない。

Grossman・Stiglitz（1980）\cite{grossman1980} は、市場が完全に効率的なら、情報を集める費用を回収できないので、誰も情報を集めなくなり、価格は効率的でなくなる、と指摘した。市場は「情報を集める人が費用を回収できる程度に」非効率なはずである。ヘッジファンドの存在は、この均衡の上にある。

\section{CAPM}

Sharpe（1964）\cite{sharpe1964}・Lintner（1965）\cite{lintner1965} の資本資産価格モデル（CAPM）は、すべての投資家が平均分散で最適化し、同じ見通しを持つとき、
\begin{equation}
\E[R_i] - R_f = \beta_i\,(\E[R_m] - R_f), \qquad \beta_i = \frac{\Cov(R_i, R_m)}{\Var(R_m)}
\end{equation}
が成り立つことを示した。個別の銘柄の平均リターンは、市場ポートフォリオとの共分散（ベータ）だけで決まり、銘柄固有の変動は分散で消せるので報われない。

CAPM の重要な帰結は、\textbf{時価総額加重の市場ポートフォリオが最適}ということである。LLM-BLM の Black--Litterman モデルが、時価総額加重の比率から「均衡のリターン」$\pi = \delta \Sigma w_{\text{mkt}}$ を逆算するのは、この考え方による（\chapref{ch:pc}）。

実証では、CAPM はうまく当てはまらない。ベータの高い株の平均リターンは CAPM の予測より低く、ベータの低い株は高い（Black・Jensen・Scholes 1972 以来の観察。第 II 部の低ボラティリティ）。

\section{マルチファクター・モデル}

\subsection{APT}

Ross（1976）\cite{ross1976} の裁定価格理論（APT）は、リターンがいくつかの共通のファクターで動くとき、
\begin{equation}
R_{i,t} = a_i + \sum_{k=1}^{K} b_{ik} f_{k,t} + \varepsilon_{i,t}
\end{equation}
裁定の機会がなければ、平均リターンはファクターへの感応度 $b_{ik}$ の線形関数になることを示した。ファクターが何かは、理論は教えない。

\subsection{Fama--French のモデル}

Fama・French（1993）\cite{ff1993} は、市場に加えて、\textbf{規模}（小型株から大型株を引いた SMB）と\textbf{割安}（簿価時価比率の高い株から低い株を引いた HML）の二つのファクターで、多くの株のポートフォリオの平均リターンが説明できることを示した（3 ファクター・モデル）。
\begin{equation}
R_{i,t} - R_{f,t} = \alpha_i + \beta_i\,\mathrm{MKT}_t + s_i\,\mathrm{SMB}_t + h_i\,\mathrm{HML}_t + \varepsilon_{i,t}
\end{equation}

その後、次のファクターが加えられた。

\begin{table}[H]
\centering\small
\caption{代表的なファクター・モデル}
\label{tab:factormodels}
\begin{tabular}{@{}P{0.33\linewidth}P{0.6\linewidth}@{}}
\toprule
モデル & ファクター \\
\midrule
CAPM & 市場 \\
Fama・French（1993） & 市場、規模（SMB）、割安（HML） \\
Carhart（1997）\cite{carhart1997} & ＋モメンタム（過去 1 年の勝ち組から負け組を引いた UMD） \\
Fama・French（2015）\cite{ff2015} & ＋収益性（RMW：収益性の高い株から低い株を引く）、投資（CMA：資産の伸びの小さい株から大きい株を引く） \\
Hou・Xue・Zhang（2015）\cite{hxz2015} & 市場、規模、投資、自己資本利益率（q ファクター・モデル） \\
\bottomrule
\end{tabular}
\end{table}

\subsection{ファクターの作り方}

ファクターは、ふつう\textbf{ポートフォリオの並べ替え}で作る。

\begin{enumerate}
\item 毎月（または毎年 6 月）、すべての銘柄を特性（簿価時価比率など）で並べる。
\item 上位 30\% と下位 30\%（または上位 10\% と下位 10\%）のポートフォリオを作る（時価総額加重か等加重）。
\item 「上位のリターン − 下位のリターン」を、その月のファクターのリターンとする。
\end{enumerate}

Fama・French は、規模と特性の 2 × 3 の並べ替えで、規模の影響を打ち消している。ファクターのリターンは、ロング・ショートのポートフォリオのリターンそのものなので、「ファクター投資」はこれを実際に売買することである。

\section{ファクターの検定}

\subsection{時系列の回帰}

ある戦略のリターンを、既知のファクターに回帰し、切片 $\alpha$ がゼロかを検定する。切片がゼロなら、戦略は既知のファクターの組み合わせで再現でき、新しい情報はない。多くのポートフォリオの切片が同時にゼロかを検定するのが、Gibbons・Ross・Shanken（1989）\cite{grs1989} の GRS 検定である。

\subsection{横断面の回帰}

Fama・MacBeth（1973）\cite{famamacbeth1973} の方法では、毎期 $t$ に、銘柄のリターンを特性（または推定したベータ）に横断面で回帰する。
\begin{equation}
R_{i,t+1} = \gamma_{0,t} + \gamma_{1,t}\,x_{i,t} + u_{i,t+1}
\end{equation}
係数 $\gamma_{1,t}$ の時系列の平均と標準誤差から、特性 $x$ がリターンを予測するかを検定する。毎期の係数は、特性 $x$ に比例して重みをつけたロング・ショートのポートフォリオのリターンと解釈できる。

\section{リスクか、ミスプライシングか}

ファクターの収益がなぜ生じるかについて、二つの説明がある。

\begin{table}[H]
\centering\small
\caption{ファクターの収益の二つの説明}
\label{tab:riskmisp}
\begin{tabular}{@{}P{0.15\linewidth}P{0.38\linewidth}P{0.38\linewidth}@{}}
\toprule
ファクター & リスクの説明 & 行動・ミスプライシングの説明 \\
\midrule
割安 & 割安な会社は業績が悪く、景気の悪いときにさらに損をする（苦境のリスク）。 & 投資家が過去の業績の悪さを将来に引き伸ばしすぎ、安く売りすぎる。 \\
モメンタム & 時間とともに変わるリスクへの感応度など、説明は難しい。 & 投資家が新しい情報に遅れて反応する（反応の不足）、または後で過剰に反応する。 \\
規模 & 小さい会社は流動性が低く、倒産しやすい。 & 注目されず、情報の反映が遅い。 \\
低ベータ & 借入れの制約のある投資家が、高ベータの株を好むため。 & 宝くじのような株を好む偏り。 \\
\bottomrule
\end{tabular}
\end{table}

どちらの説明でも、裁定には限界がある。Shleifer・Vishny（1997）\cite{shleifer1997} は、裁定をする運用者は他人の資金を使っていて、歪みが一時的に広がって損をすると、投資家が資金を引き上げるため、最も歪みの大きいときに裁定ができなくなることを示した（裁定の限界）。これは、LTCM や 2007 年の quant quake の説明にもなる（第 II 部）。

\section{ファクターの大きさ}

米国株式の長期のデータ（Kenneth French のデータ・ライブラリー、1920 年代後半から）で、代表的なファクターの月次の平均は、おおよそ次のとおりである（期間によって大きく変わる）。

\begin{itemize}
\item 市場の超過リターン：年 7〜8\% 前後、年率の標準偏差 18\% 前後。
\item 割安（HML）：年 3〜4\% 前後。ただし 2007 年ごろから 2020 年まで、長く損をした。
\item モメンタム（UMD）：年 7\% 前後と大きいが、まれに大きく崩れる（2009 年の春など）。
\item 規模（SMB）：年 2〜3\% 前後で、1980 年代以降は小さい。
\end{itemize}

どのファクターも、手数料・取引コストの前の値であり、何年も続けて損をする期間がある。ファクターの研究と、その崩れについては第 II 部（第 \ref{ch:equityfactor} 章）で扱う。

\section{ファクター・ズーと多重検定}

1990 年代以降、平均リターンを予測するとされる特性が数百も報告された。Cochrane（2011）\cite{cochrane2011} はこれを「ファクターの動物園」と呼び、どれが本当に独立な情報を持つのかを問題にした。Harvey・Liu・Zhu（2016）\cite{hlz2016} は、多くの研究者が同じデータで多くの特性を試していることを考えると、新しいファクターの $t$ 値は 2 ではなく 3 を超える必要があると論じた（\chapref{ch:backtest}）。

\section{ファクターとヘッジファンド}

ファクターは、ヘッジファンドの世界で三つの役割を持つ。

\begin{enumerate}
\item \textbf{戦略そのもの}：ファクターのロング・ショートは、クオンツ株式の中心である。
\item \textbf{リスクの物差し}：運用チームの収益のうち、ファクターで説明できる部分（安く手に入る部分）を除き、残りを「本当のアルファ」として評価する。マルチマネージャーは、運用チームにファクターのリスクを打ち消すよう求める。
\item \textbf{リスクモデル}：銘柄の共分散を、少数のファクターの共分散で近似する（\chapref{ch:risk}）。
\end{enumerate}

\begin{projbox}
LLM-BLM では、リッジ回帰の統制に、業種と時期の固定効果、規模（浮動株時価総額の対数）、発表の前のリターン（モメンタムに近い）を入れた。これは、「LLM の点数が、既知のファクターとは別の情報を持つか」を確かめるための、ファクター・モデルの考え方そのものである。
\end{projbox}

\chapter{シグナル}
\label{ch:signal}

\section{シグナルとは}

クオンツ運用で\textbf{シグナル}（アルファ・シグナル、ファクター、特徴量）とは、各時点 $t$・各資産 $i$ について計算される数 $s_{i,t}$ で、将来の（相対的な）リターンを予測するものをいう。

\begin{itemize}
\item 時点 $t$ で\textbf{実際に知られていた}情報だけから計算する。
\item 単独の銘柄の上げ下げより、\textbf{銘柄の間の比較}（どれが相対的に上がるか）を予測するものが多い（横断面のシグナル）。
\item 一つ一つは弱い（IC が 0.01〜0.05 程度）。多くを組み合わせて使う。
\end{itemize}

先物やマクロの戦略では、一つの資産の時系列の上げ下げを予測するシグナル（時系列のシグナル）も多い。

\section{シグナルの種類}

\begin{table}[H]
\centering\small
\caption{代表的なシグナル（株式の場合）}
\label{tab:signals}
\begin{tabular}{@{}P{0.13\linewidth}P{0.3\linewidth}P{0.3\linewidth}P{0.18\linewidth}@{}}
\toprule
分類 & 例 & 計算の例 & 効く期間 \\
\midrule
価格 & モメンタム & 過去 12 か月から直近 1 か月を除いたリターン（12-1） & 数か月〜1 年 \\
 & 短期リバーサル & 直近 1 か月（または 1 週）のリターンの逆 & 数日〜1 か月 \\
 & 低ボラティリティ & 過去 1 年の日次リターンの標準偏差の逆、ベータの逆 & 数か月〜年 \\
\midrule
財務 & 割安 & 簿価時価比率、利益÷株価、キャッシュフロー÷企業価値 & 年 \\
 & 収益性・クオリティ & 売上総利益÷総資産、ROE、利益の安定性、負債の少なさ & 年 \\
 & 投資・発行 & 総資産の伸び、株式の発行（伸びが大きいほど低いリターン） & 年 \\
 & 会計の質 & 発生主義の利益の割合（アクルーアル）の大きさの逆 & 年 \\
\midrule
予想・発表 & 決算の驚き & SUE＝（実際の利益 − 予想）÷ 標準偏差 & 数日〜数か月 \\
 & 予想の修正 & アナリストの利益予想の上方修正の数 − 下方修正の数 & 数週〜数か月 \\
\midrule
需給 & 空売りの残高 & 空売りの株数 ÷ 浮動株（多いほど低いリターン） & 数週〜数か月 \\
 & 資金の流れ & 投資信託の資金の流出入による売買の圧力 & 数週〜数か月 \\
\midrule
代替データ & テキスト & 決算発表文・ニュース・SNS の語調、LLM の読み取り & 数日〜数週 \\
 & その他 & カード決済、衛星画像、ウェブの閲覧、求人、アプリの利用 & 数週〜四半期 \\
\bottomrule
\end{tabular}
\end{table}

先物・マクロでは、次のものが代表的である。

\begin{itemize}
\item \textbf{時系列モメンタム（トレンド）}：その資産自身の過去 1〜12 か月のリターンの符号、移動平均の交差。
\item \textbf{キャリー}：価格が変わらなかったときに得られる収益。通貨なら金利の差、債券なら利回りと年限の効果、商品なら先物の期近と期先の価格差。
\item \textbf{バリュー}：購買力平価からの乖離（通貨）、実質の利回り（債券）、長期の平均からの乖離。
\end{itemize}

\section{シグナルの作り方}

生のデータから使えるシグナルを作るまでに、次の処理をする。

\subsection{時点の正しいデータ}

その時点で知られていた値だけを使う。

\begin{itemize}
\item \textbf{発表の遅れ}：12 月期の決算の数値が公表されるのは、翌年の 2〜3 月である。「12 月末の簿価」を 12 月末の売買に使うと、未来の情報を使ったことになる（先読み、look-ahead）。Fama・French は、前年度の財務データを 6 月末から使う。
\item \textbf{訂正}：財務データは後から訂正される。データベースの「最新の値」は、当時の値とは違う。当時の値を残したデータ（ポイント・イン・タイムのデータ）が要る。
\item \textbf{銘柄の範囲}：今日存在する銘柄だけで過去を検証すると、倒産・上場廃止した銘柄が抜け、成績が良く見える（生存者バイアス）。
\end{itemize}

\subsection{横断面の標準化}

シグナルの水準は、時期によって変わる（例えば、全体の利益の水準）。そこで、毎時点、銘柄の間で標準化する。
\begin{equation}
z_{i,t} = \frac{s_{i,t} - \bar{s}_t}{\mathrm{sd}_t(s)} \qquad \text{または} \qquad
z_{i,t} = \Phi^{-1}\!\left(\frac{\mathrm{rank}_t(s_{i,t}) - 0.5}{N_t}\right)
\end{equation}
後者は順位を正規分布の分位点に直したもの（正規スコア）で、外れ値の影響を受けにくい。LLM-BLM の BLM の見通しでも、点数を正規スコアにした。

\subsection{外れ値の処理}

上下 1\% や平均から 3 標準偏差の値で丸める（ウィンソライズ）。財務の比率は分母が小さいと極端な値になるので、特に必要である。

\subsection{中立化}

シグナルが、意図しないファクターを含んでいることが多い。例えば「割安」は特定の業種（金融、エネルギー）に偏り、「低ボラティリティ」は大型株と公益事業に偏る。毎時点、シグナルを業種のダミーや規模・ベータに回帰し、その残差をシグナルとする（中立化）。
\begin{equation}
z_{i,t} = \sum_{g} c_{g,t}\,\mathbb{1}[i \in g] + d_t\,\ln(\mathrm{size}_{i,t}) + \tilde{z}_{i,t}
\end{equation}
残差 $\tilde{z}_{i,t}$ は、業種と規模の中で比べたシグナルになる。

\section{シグナルの評価}

\subsection{IC}

時点 $t$ のシグナル $z_{i,t}$ と、その後のリターン $r_{i,t+h}$ の銘柄の間の相関を\textbf{情報係数}（IC）と呼ぶ。外れ値に強い順位相関（スピアマン）がよく使われる（rank IC）。
\begin{equation}
\IC_t = \mathrm{corr}_i\big(\mathrm{rank}(z_{i,t}),\ \mathrm{rank}(r_{i,t+h})\big)
\end{equation}
時点ごとの $\IC_t$ の平均と、その安定性（平均÷標準偏差、IC の IR）を見る。

\subsection{分位のポートフォリオ}

シグナルで銘柄を 5 つまたは 10 の組に分け、組ごとの平均リターンを見る。平均リターンが組の順に単調に並ぶか、上位と下位の差（スプレッド）が大きいか、上位と下位のどちらに収益が偏っているかを確かめる。空売りが難しい銘柄（小型・流動性が低い）で下位の収益が大きいなら、実際にはとれない。

\subsection{減衰}

シグナルを計算した後、$h = 1, 2, \dots$ 日のリターンに対する IC を描くと、シグナルの情報がどれだけ速く価格に入るかが分かる（減衰の曲線）。

\begin{figure}[H]
\centering
\begin{tikzpicture}[x=9mm,y=28mm]
\draw[arr] (0,0)--(11,0) node[below left,font=\footnotesize]{シグナルを計算してからの日数};
\draw[arr] (0,0)--(0,1.15) node[above,font=\footnotesize]{累積の IC（相対値）};
\draw[c2,line width=1pt,domain=0:10.5,samples=60] plot (\x,{1-exp(-3*\x)});
\draw[c1,line width=1pt,domain=0:10.5,samples=60] plot (\x,{1-exp(-0.25*\x)});
\node[lab,c2,anchor=west] at (1.2,1.05) {速い減衰（発表の当日に入りきる）};
\node[lab,c1,anchor=west] at (5.5,0.62) {遅い減衰（数週間かけて入る）};
\foreach \x in {0,2,4,6,8,10} \node[lab,below] at (\x,0) {\x};
\end{tikzpicture}
\caption{シグナルの減衰の概念図。速く減衰するシグナルは、すぐに売買しなければ価値がなく、売買の回数も多くなる。}
\label{fig:decay}
\end{figure}

\subsection{回転率と容量}

シグナルの順位が毎期どれだけ入れ替わるか（自己相関）で、売買の量が決まる。シグナルの自己相関が低いほど、回転率が高く、取引コストがかかる。同じ IC でも、ゆっくり変わるシグナルのほうが、実際の運用では価値が高いことが多い。

\subsection{ほかのシグナルとの相関}

新しいシグナルが、既知のシグナル（モメンタム、割安など）と強く相関しているなら、新しい情報は少ない。既知のシグナルに回帰した残差の IC（増分の IC）を見る。

\section{シグナルの組み合わせ}

多くのシグナル $z^{(1)}, \dots, z^{(K)}$ を一つの予測にまとめる方法は、単純なものから順に次のとおり。

\begin{enumerate}
\item \textbf{等しい重み}：標準化したシグナルの平均。重みを推定しないので、過学習しない。意外に強い。
\item \textbf{IC で重み}：過去の IC（あるいは IC の IR）に比例した重み。
\item \textbf{回帰}：過去の将来リターンを目的変数に、シグナルを説明変数にした回帰の係数を重みにする（リッジなどで縮める）。シグナルの間の相関を考慮できる。
\item \textbf{機械学習}：決定木やニューラルネットで、非線形な組み合わせや相互作用を学ぶ（\chapref{ch:model}）。
\end{enumerate}

重みの推定には誤差があり、推定の期間で良かったシグナルが、次の期間も良いとは限らない。重みを等しい重みに向けて縮めること、重みの変化をゆっくりにすることが、実務の定石である。

\section{シグナルの寿命}

シグナルは、時間とともに弱くなることが多い。

\begin{itemize}
\item \textbf{公表}：論文で公表されたシグナルは、公表の後に収益が小さくなる（第 II 部）。
\item \textbf{混雑}：同じシグナルで売買する資金が増えると、収益が小さくなり、一斉に売るときの損が大きくなる。
\item \textbf{構造の変化}：データの入手が容易になる（代替データの普及）、取引の仕組みが変わる（手数料の自由化、ティックの縮小）。
\end{itemize}

運用会社は、シグナルの成績を常に監視し、新しいシグナルを探し続ける。研究の部門の大きな仕事は、「まだ広く使われていないデータ」と「既知のデータの新しい使い方」を探すことである。

\begin{projbox}
LLM の決算発表文の読み取りは、表~\ref{tab:signals} の「テキスト」の一つである。LLM-BLM の実験では、2 日の反応に対する IC は高かったが、減衰の曲線（\figref{fig:decay}）でいえば「速い減衰」で、発表の翌日以降のリターンへの予測力はほとんど残らなかった。この種のシグナルを運用で使うには、発表の直後（数分〜当日）に売買するか、反応の遅い銘柄（注目の低い小型株：実験 S1 の仮説）を探す必要がある。
\end{projbox}

\chapter{予測モデル}
\label{ch:model}

\section{予測の問題}

シグナルの組 $\bm{x}_{i,t} \in \mathbb{R}^K$ から、将来のリターン $r_{i,t+h}$ を予測する関数 $f$ を学ぶ。
\begin{equation}
r_{i,t+h} = f(\bm{x}_{i,t}) + \varepsilon_{i,t+h}
\end{equation}
資産の予測は、画像認識などとは性質がまったく違う。

\begin{itemize}
\item \textbf{信号が極めて弱い}：月次の株のリターンの変動のうち、予測できる部分は 1\% に満たない（決定係数 $R^2 < 0.01$）。99\% 以上は雑音である。
\item \textbf{分布が変わる}：市場の参加者が学習して行動を変えるので、過去の関係が将来も続くとは限らない（非定常）。
\item \textbf{データが少ない}：米国株の月次のデータは、銘柄が数千でも、時点は数百しかない。独立な時点の数が、過学習を防ぐうえでの実質的な標本の大きさになる。
\item \textbf{予測が当たると、それ自体で消える}：多くの人が同じ予測で売買すると、価格がそれを織り込む。
\end{itemize}

このため、単純で頑健なモデルが、複雑なモデルと同じくらい、あるいはそれ以上に良いことが多い。

\section{線形のモデル}

\subsection{最小二乗}

最も基本的なのは、線形回帰
\begin{equation}
\hat{r}_{i,t+h} = \bm{x}_{i,t}\T \bm{b}, \qquad \hat{\bm{b}}_{\mathrm{OLS}} = (X\T X)^{-1} X\T \bm{y}
\end{equation}
である。説明変数が多く、互いに相関していると、係数の推定値が大きく振れ、標本外で悪くなる。

\subsection{正則化回帰}

係数の大きさに罰則をつけて、推定を安定させる。
\begin{align}
\text{リッジ}&: \quad \hat{\bm{b}} = \arg\min_{\bm b} \|\bm y - X \bm b\|^2 + \lambda \|\bm b\|_2^2 = (X\T X + \lambda I)^{-1} X\T \bm y \\
\text{ラッソ}&: \quad \hat{\bm{b}} = \arg\min_{\bm b} \|\bm y - X \bm b\|^2 + \lambda \|\bm b\|_1 \\
\text{エラスティックネット}&: \quad \hat{\bm{b}} = \arg\min_{\bm b} \|\bm y - X \bm b\|^2 + \lambda\big(\alpha \|\bm b\|_1 + (1-\alpha)\|\bm b\|_2^2\big)
\end{align}

\begin{itemize}
\item \textbf{リッジ}は、すべての係数をゼロに向けて縮める。説明変数の主成分のうち、分散の小さい方向の係数を強く縮める。説明変数の数が標本より多くても解ける。LLM-BLM で、言語モデルのベクトル（数百〜数千次元）を株価の反応に結びつけるのに使ったのはこれである。
\item \textbf{ラッソ}は、多くの係数をちょうどゼロにする（変数の選択）。相関した変数の中から一つだけを選びがちで、どれを選ぶかが不安定になる。
\item \textbf{エラスティックネット}は、両者の中間。
\end{itemize}

罰則の強さ $\lambda$ は、交差検証（\ref{sec:cv} 節）で選ぶ。

\subsection{次元の削減}

説明変数が多いとき、まず少数の成分にまとめてから回帰する。

\begin{itemize}
\item \textbf{主成分回帰}（PCR）：説明変数の主成分（分散の大きい方向）の上位を使う。目的変数を見ずに成分を選ぶので、予測に効かない成分を選ぶことがある。
\item \textbf{部分最小二乗}（PLS）：目的変数との共分散が大きい方向を選ぶ。
\end{itemize}

\section{非線形のモデル}

\subsection{決定木とアンサンブル}

\textbf{決定木}は、説明変数の値で標本を分け続け、各葉の平均を予測にする。相互作用（「小型株で、かつ割安なら」）を自然に表せるが、一本の木は不安定で過学習しやすい。そこで多くの木を組み合わせる。

\begin{itemize}
\item \textbf{ランダムフォレスト}：標本と変数を無作為に選んで多くの木を育て、平均する。分散を下げる。
\item \textbf{勾配ブースティング}（XGBoost、LightGBM など）：前の木の誤差（残差）を次の木で当てることを繰り返す。表形式のデータで最も強い方法の一つ。
\end{itemize}

\subsection{ニューラルネット}

多層のパーセプトロン（全結合のネットワーク）は、任意の滑らかな関数を近似できる。資産価格の予測では、層の浅い（2〜3 層）ネットワークが、深いものより良いという報告が多い（第 II 部、第 \ref{ch:ml} 章）。信号が弱く、データが少ないためである。

時系列の構造を使うモデル（LSTM、トランスフォーマー）は、価格の系列から直接予測するのに使われるが、資産の予測で一貫して強いという証拠は、まだ限られる。

\subsection{言語モデル}

LLM や FinBERT のような言語モデルは、文章を数値のベクトル（埋め込み）や点数に変える「特徴量の抽出器」として使う。抽出した特徴量を、上の線形・非線形のモデルで将来のリターンに結びつける。LLM に直接「この株は上がるか」と聞く使い方もあるが、LLM が学習した期間の後の出来事を知っている（先読み）おそれがあり、検証が難しい。

\section{時系列のモデル}

一つの資産の時系列を予測するモデルもある。

\begin{itemize}
\item \textbf{自己回帰}（AR）：$r_t = c + \sum_{j=1}^{p} \phi_j r_{t-j} + \varepsilon_t$。ARMA は誤差の移動平均を加える。
\item \textbf{ベクトル自己回帰}（VAR）：複数の系列（金利、インフレ、株価）を同時に扱う。
\item \textbf{状態空間モデルとカルマン・フィルター}：観測できない状態（真の価値、時間とともに変わるベータ）を、観測から逐次推定する。ペアトレードのヘッジ比率の推定などに使う。
\item \textbf{レジーム・スイッチング}（Hamilton 1989 \cite{hamilton1989}）：経済が「好況」「不況」などの状態を確率的に行き来するとし、状態ごとに別のモデルを当てる。
\item \textbf{ボラティリティのモデル}（GARCH など）：リターンの平均より、分散の予測のほうがずっと容易である（\chapref{ch:risk}）。
\end{itemize}

\section{予測の評価}

\subsection{標本外の決定係数}

Campbell・Thompson（2008）\cite{campbellthompson2008} は、標本外の予測の決定係数を、過去の平均 $\bar r$ を予測に使った場合と比べて定義した。
\begin{equation}
R^2_{\mathrm{OS}} = 1 - \frac{\sum_t (r_t - \hat r_t)^2}{\sum_t (r_t - \bar r_t)^2}
\end{equation}
$R^2_{\mathrm{OS}} > 0$ なら、過去の平均より良い予測である。

\subsection{小さな \texorpdfstring{$R^2$}{R2} の価値}

$R^2$ が小さくても、経済的な価値は大きいことがある。一つの資産の売買のタイミングでいえば、もとのシャープレシオを $\SR_0$ とし、予測の $R^2$ を使って最適に売買すると、シャープレシオは
\begin{equation}
\SR^{*} = \sqrt{\frac{\SR_0^2 + R^2}{1 - R^2}}
\end{equation}
になる（Campbell・Thompson 2008 の考え方）。月次の $R^2 = 0.5\%$ でも、$\SR_0 = 0$ から月次の $\SR^{*} \approx 0.071$、年率で約 0.25 になる。横断面の予測では、多くの銘柄に同時に賭けられるので、基本法則（\ref{sec:flam} 節）により、さらに大きくなりうる。

\subsection{そのほかの物差し}

\begin{itemize}
\item \textbf{IC}：予測とリターンの順位相関（\chapref{ch:signal}）。
\item \textbf{的中率}：符号が当たった割合。50\% を少し超えれば十分なことが多い。
\item \textbf{ポートフォリオの成績}：予測で並べた分位のポートフォリオのリターン。最終的な目的に最も近い。
\item \textbf{予測の比較の検定}：二つの予測の誤差の差の平均がゼロかを検定する（Diebold・Mariano 1995 \cite{dieboldmariano1995}）。
\end{itemize}

\section{検証の方法}
\label{sec:cv}

\subsection{時系列の分割}

ふつうの機械学習の交差検証は、データを無作為に分ける。資産のデータでこれをすると、未来のデータで学習して過去を予測することになり、成績が良く見える。そこで、時間の順に分ける。

\begin{figure}[H]
\centering
\begin{tikzpicture}[x=11mm,y=5.5mm]
\foreach \k in {0,1,2,3} {
  \pgfmathsetmacro{\e}{4+\k*2}
  \fill[c1!35] (0,-\k) rectangle (\e,-\k+0.7);
  \fill[c2!60] (\e+0.3,-\k) rectangle (\e+2.3,-\k+0.7);
  \node[lab,anchor=east] at (-0.1,-\k+0.35) {第 \the\numexpr\k+1\relax 回};
}
\fill[c1!35] (0,-5) rectangle (0.8,-4.3); \node[lab,anchor=west] at (0.9,-4.65) {学習};
\fill[c2!60] (2.6,-5) rectangle (3.4,-4.3); \node[lab,anchor=west] at (3.5,-4.65) {評価};
\draw[c4,line width=0.8pt] (4,-5) -- (4.3,-5) -- (4.3,-4.3) -- (4,-4.3) -- cycle; \node[lab,anchor=west] at (4.4,-4.65) {空ける（エンバーゴ）};
\draw[arr] (0,-3.6)--(11.2,-3.6) node[below left,font=\scriptsize]{時間};
\end{tikzpicture}
\caption{ウォークフォワードの検証（学習の窓を広げていく場合）。学習と評価の間に隙間を空けるのは、正解のリターンの期間が重なるのを防ぐため。}
\label{fig:walkforward}
\end{figure}

\begin{itemize}
\item \textbf{ウォークフォワード}：ある日までのデータで学習し、その後の期間で評価する。評価の期間を進めながら繰り返す。学習の窓は、広げていく（拡大窓）か、一定の長さで動かす（移動窓）。
\item \textbf{パージとエンバーゴ}：正解のリターンが期間を持つ（例：60 日後までのリターン）と、学習の最後の標本の正解の期間が、評価の期間と重なる。重なる標本を学習から除き（パージ）、さらに少し隙間を空ける（エンバーゴ）。López de Prado（2018）\cite{lopezdeprado2018} が体系化した。
\end{itemize}

\subsection{ハイパーパラメータの選び方}

罰則の強さや木の深さなどは、学習の期間の中だけで、さらに時系列の分割をして選ぶ（入れ子の交差検証）。評価の期間を見てハイパーパラメータを選ぶと、評価が楽観的になる。

\begin{projbox}
LLM-BLM のリッジ回帰は、四半期ごとに、その四半期が始まる前に正解の 2 日の窓が閉じた発表だけで学習し、その四半期の発表を予測した。これは、拡大窓のウォークフォワードに、パージを組み合わせた形である。罰則の強さは、学習の期間の中の交差検証（RidgeCV）で選んだ。
\end{projbox}

\chapter{ポートフォリオ構築}
\label{ch:pc}

予測（期待リターン）とリスク（共分散）が手に入ったとして、それをどのような比率の組 $\bm w$ にするか。これがポートフォリオ構築である。予測が弱く誤差が大きい世界では、この工程の作り方で成績が大きく変わる。

\section{平均分散最適化}

\subsection{Markowitz の問題}

Markowitz（1952）\cite{markowitz1952} は、ポートフォリオの期待リターンと分散のトレードオフを最適化の問題にした。$N$ 資産の期待超過リターンを $\bm\mu$、共分散を $\Sigma$ とすると、
\begin{equation}
\max_{\bm w}\ \bm w\T \bm\mu - \frac{\lambda}{2}\,\bm w\T \Sigma \bm w
\end{equation}
制約がなければ、解は
\begin{equation}
\bm w^* = \frac{1}{\lambda}\,\Sigma^{-1}\bm\mu
\label{eq:mvo}
\end{equation}
である。$\lambda$ はリスクを嫌う度合い。すべての $\lambda$ の解は、同じ方向 $\Sigma^{-1}\bm\mu$ に比例し（接点ポートフォリオ）、$\lambda$ はレバレッジだけを決める。

\subsection{推定誤差の問題}

式~\eqref{eq:mvo} は、$\bm\mu$ と $\Sigma$ が正確に分かっていれば最適だが、実際には推定値しかない。

\begin{itemize}
\item \textbf{期待リターンの誤差}：月次のリターンの平均を、標準誤差 1/10 の精度で推定するには、何百年ものデータが要る。期待リターンの推定値は、ほとんど雑音である。
\item \textbf{逆行列による増幅}：$\Sigma^{-1}$ は、共分散の小さな固有値の方向を大きく増幅する。相関の高い 2 資産の小さなリターンの差が、巨大なロングとショートになる。
\end{itemize}
Michaud（1989）\cite{michaud1989} は、平均分散最適化を「誤差の最大化」と呼んだ。推定誤差の最も大きい資産に、最も大きな比率をつけてしまうからである。DeMiguel・Garlappi・Uppal（2009）\cite{dgu2009} は、多くのデータで、推定に基づく最適化が、すべての資産に等しい比率をつける $1/N$ のポートフォリオに、標本外で安定して勝てないことを示した。

\subsection{対策}

\begin{enumerate}
\item \textbf{期待リターンを縮める}：推定値を、全体の平均や均衡のリターンに向けて縮める（ベイズ、James--Stein、Black--Litterman）。
\item \textbf{共分散を安定させる}：縮小推定やファクター・モデル（\chapref{ch:risk}）。
\item \textbf{制約をかける}：比率の上限、ロングのみ、TE の上限。Jagannathan・Ma（2003）\cite{jagannathanma2003} は、空売りの禁止が共分散の縮小と同じ働きをすることを示した。
\item \textbf{期待リターンを使わない}：最小分散、リスクパリティ。
\end{enumerate}

\section{Black--Litterman}
\label{sec:bl}

Black・Litterman（1992）\cite{blacklitterman1992} は、平均分散の不安定さを、ベイズの考え方で解決した。LLM-BLM の中心の道具なので、ここで要点をまとめる。

\subsection{均衡のリターン}

すべての投資家が平均分散で最適化した結果が、時価総額加重の市場ポートフォリオ $\bm w_{\text{mkt}}$ だと考えると、式~\eqref{eq:mvo} を逆に解いて、市場が「暗に」想定している期待リターンが得られる。
\begin{equation}
\bm\pi = \delta\,\Sigma\,\bm w_{\text{mkt}}
\end{equation}
$\delta$ は市場全体のリスク回避度（2.5 前後がよく使われる）。これを事前分布の平均とし、その不確かさを $\tau\Sigma$ とする（$\tau$ は 0.01〜0.05 程度の小さな数）。

\subsection{見通しとの合成}

運用者の見通しを、$P\bm\mu = \bm q + \bm\varepsilon$、$\bm\varepsilon \sim N(0, \Omega)$ の形で書く。$P$ は見通しの対象（どの資産、あるいはどの組み合わせ）、$\bm q$ は見通しの値、$\Omega$ は見通しの不確かさ。ベイズの公式で、事後の期待リターンは
\begin{equation}
\bm\mu_{\mathrm{BL}} = \big[(\tau\Sigma)^{-1} + P\T\Omega^{-1}P\big]^{-1}\big[(\tau\Sigma)^{-1}\bm\pi + P\T\Omega^{-1}\bm q\big]
\end{equation}
となる。見通しのない資産は、均衡のリターン $\bm\pi$ のままで、相関を通じて少しだけ動く。見通しが確かなほど（$\Omega$ が小さいほど）、事後は見通しに近づく。He・Litterman（1999）\cite{helitterman1999} は、$\Omega$ を $\mathrm{diag}(P\,\tau\Sigma\,P\T)$ とする置き方を示し、広く使われている。

\subsection{見通しの大きさ}

見通しの値 $\bm q$ を、予測の点数からどう作るか。Grinold・Kahn の考え方では、標準化した点数 $z$ を、IC と資産の変動の大きさでリターンの単位に直す。
\begin{equation}
\alpha_i = \IC \times \sigma_i \times z_i
\end{equation}
IC が 0.05、年率の変動が 30\%、点数が $+1$ 標準偏差なら、見通しは年 1.5\% の超過リターンである。予測が弱いほど、見通しは小さく、ポートフォリオは市場の比率に近くなる。これは、予測の誤差を自動的に考慮した、合理的な縮め方である。

\begin{projbox}
LLM-BLM の実験 7・9 と S1 の追記で使った BLM は、この形（$\delta = 2.5$、$\tau = 0.025$、He--Litterman の $\Omega$、$\IC = 0.05$）である。見通しの単位を、保有期間 $h$ 日に合わせて日次で $\IC \times \sigma_{\text{日}} \times z / \sqrt{h}$ にするのは、予測の期間と共分散の期間をそろえるためである（単位を取り違えると、見通しが何十倍にもなる）。
\end{projbox}

\section{期待リターンを使わない方法}

\subsection{最小分散}

$\min_{\bm w} \bm w\T\Sigma\bm w$ s.t. $\bm 1\T\bm w = 1$。解は $\bm w \propto \Sigma^{-1}\bm 1$。期待リターンを推定しないので、推定誤差の影響が小さい。結果として、変動の小さい・相関の低い資産に偏り、低ボラティリティのファクターを持つ。

\subsection{逆ボラティリティとリスクパリティ}

\begin{itemize}
\item \textbf{逆ボラティリティ}：$w_i \propto 1/\sigma_i$。相関を無視した簡単な方法。
\item \textbf{リスクパリティ}（等リスク寄与）：各資産がポートフォリオのリスクに寄与する量をそろえる。資産 $i$ のリスク寄与は
\begin{equation}
\mathrm{RC}_i = w_i\,\frac{(\Sigma\bm w)_i}{\sqrt{\bm w\T\Sigma\bm w}}, \qquad \sum_i \mathrm{RC}_i = \sqrt{\bm w\T\Sigma\bm w}
\end{equation}
で、すべての $\mathrm{RC}_i$ が等しくなる $\bm w$ を数値的に求める（Maillard・Roncalli・Teïletche 2010 \cite{mrt2010}）。
\end{itemize}

株式と債券の 60/40 のポートフォリオでは、リスクの 9 割前後が株式から来る。リスクパリティは、債券の比率を大きくし（必要ならレバレッジをかけて）、株と債券のリスクをそろえる。Bridgewater の All Weather などがこの考え方で知られる。

\subsection{階層的リスクパリティ}

López de Prado（2016）\cite{hrp2016} の階層的リスクパリティ（HRP）は、共分散の逆行列を使わない。
\begin{enumerate}
\item 相関から資産の間の距離 $d_{ij} = \sqrt{(1-\rho_{ij})/2}$ を作り、階層的にクラスタリングする。
\item 似た資産が隣り合うように、共分散の行と列を並べ替える。
\item 木の上から、二つの枝にリスクの逆数に比例して資金を分け、再帰的に下へ進む。
\end{enumerate}
逆行列の不安定さを避けるので、資産の数が多く、共分散が不安定なときに頑健だとされる。

\section{レバレッジの決め方}

\subsection{ケリー基準}

繰り返し賭けて資産の対数の期待値（長期の複利の成長率）を最大にする賭けの割合が、ケリー基準である（Kelly 1956 \cite{kelly1956}）。連続時間で、期待超過リターン $\mu$、変動 $\sigma$ の資産一つなら、
\begin{equation}
f^* = \frac{\mu}{\sigma^2}, \qquad \text{多資産なら}\ \bm f^* = \Sigma^{-1}\bm\mu
\end{equation}
このときの成長率は $g = r + \frac{1}{2}\SR^2$（$r$ は安全資産の金利）である。$f^*$ の 2 倍賭けると成長率は $r$ まで下がり、それ以上ではマイナスになる。

$\mu$ の推定の誤差が大きいので、ケリー基準どおりに賭けると、実際には賭けすぎになりやすい。実務では、その半分（ハーフ・ケリー）以下にすることが多い。ハーフ・ケリーは、成長率を 3/4 に保ちながら、変動を半分にする。

\subsection{ボラティリティ・ターゲット}

ポートフォリオの予測した変動 $\hat\sigma_t$ が、目標 $\sigma^*$ になるように、全体の規模を毎日調整する。
\begin{equation}
L_t = \frac{\sigma^*}{\hat\sigma_t}
\end{equation}
変動が大きいときに規模を落とす。トレンドフォローやリスクパリティでは標準的な方法で、急落の時期に変動が高まると自動的にポジションを減らす。ただし、多くの運用者が同時に同じことをすると、変動が上がったときに一斉に売ることになり、下げを加速しうる。

\section{制約}

\subsection{よく使われる制約}

\begin{table}[H]
\centering\small
\caption{ポートフォリオの制約}
\label{tab:constraints}
\begin{tabular}{@{}P{0.25\linewidth}P{0.4\linewidth}P{0.27\linewidth}@{}}
\toprule
制約 & 式の例 & 目的 \\
\midrule
予算 & $\bm 1\T\bm w = 1$（ロング・ショートなら $= 0$） & 資金の全額を使う／ドル・ニュートラル \\
ロングのみ & $w_i \geq 0$ & 規制・約款 \\
比率の上限 & $|w_i - w_{b,i}| \leq c$ & 集中を防ぐ \\
トラッキングエラー & $(\bm w - \bm w_b)\T\Sigma(\bm w - \bm w_b) \leq \TE^2$ & 指数から離れすぎない \\
ベータ・ニュートラル & $\bm\beta\T\bm w = 0$ & 市場の動きを打ち消す \\
ファクター・ニュートラル & $B\T\bm w = \bm 0$ & 業種・スタイルのリスクを打ち消す \\
グロス・レバレッジ & $\sum_i |w_i| \leq G$ & 借入れの上限 \\
回転率 & $\sum_i |w_i - w_{i}^{\text{前}}| \leq u$ & 取引コストを抑える \\
流動性 & $|w_i - w_i^{\text{前}}|\,A \leq k \cdot \mathrm{ADV}_i$ & 売買高に比べて取引を小さく \\
\bottomrule
\end{tabular}
\end{table}

これらの制約は、ほとんどが線形か二次なので、ポートフォリオの問題は\textbf{凸の二次計画}（または二次錐計画）になり、数千銘柄でも速く解ける（CVXPY、MOSEK、Gurobi など）。

\subsection{制約の費用}

制約は、予測を比率に移す効率（移転係数 TC、\ref{sec:flam} 節）を下げる。特に\textbf{ロングのみ}の制約は大きい。時価総額加重の指数では、小さな銘柄の比率は 0.01\% 程度しかないので、「悪い」と予測しても、比率をゼロにする分（$-0.01\%$）しか賭けられない。良いと予測した銘柄を買う側にしか、まともに賭けられないのである。Clarke・de Silva・Thorley（2002）\cite{cdt2002} は、ロングのみの制約で TC が 1 を大きく下回り、予測の価値の多くが失われることを、数値例で示した。

\subsection{ロング・ショートの比率}

ヘッジファンドのロング・ショートでは、次の言葉を使う。
\begin{itemize}
\item \textbf{グロス}：買いと売りの金額の和を、自己資本で割ったもの。マーケット・ニュートラルの株式では 3〜6 倍程度のことが多い。
\item \textbf{ネット}：買い − 売り。ゼロならドル・ニュートラル。
\item \textbf{130/30}：買い 130\%、売り 30\% のポートフォリオ。ロングのみの制約を少しだけ緩め、TC を上げる。
\end{itemize}

\section{多期間とコスト}

ここまでは 1 期間の問題だった。実際には、毎期の取引コストがかかり、予測は時間とともに減衰する。Gârleanu・Pedersen（2013）\cite{garleanupedersen2013} は、コストが売買の量の二乗に比例するとき、最適な売買は「現在の比率と、将来の目標の比率の加重平均へ、少しずつ近づく」ことであり、目標の比率も、減衰の遅い予測を重く、速い予測を軽く見るべきことを示した（「目標に向けて部分的に売買し、目標の先を狙う」）。取引コストの扱いは \chapref{ch:cost} で詳しく扱う。

\begin{exbox}[同じ予測、違う構築]
IC 0.05 の同じ予測を、次の 3 通りで比率にしたとする（概念的な比較）。
\begin{enumerate}
\item \textbf{ロングのみ、TE 3\%}：悪い銘柄を空売りできない。TC は低く、予測の価値の多くが失われる。LLM-BLM の設計。
\item \textbf{130/30、TE 3\%}：悪い銘柄を少し空売りできる。TC が上がる。
\item \textbf{マーケット・ニュートラル、グロス 4 倍}：予測を両側に使い、市場のリスクを打ち消す。TC は高いが、空売りの費用・借入れの費用・レバレッジのリスクを負う。
\end{enumerate}
どれが良いかは、予測の質だけでなく、空売りの費用と、許されるリスクで決まる。
\end{exbox}

\chapter{リスク管理}
\label{ch:risk}

ヘッジファンドが潰れる理由は、予測が外れたことよりも、リスクの取り方を誤ったことのほうが多い。この章では、リスクを測る道具（リスクモデル、VaR）と、リスクを管理する仕組みをまとめる。

\section{リスクの種類}

\begin{table}[H]
\centering\small
\caption{ヘッジファンドが直面するリスク}
\label{tab:risks}
\begin{tabular}{@{}P{0.2\linewidth}P{0.72\linewidth}@{}}
\toprule
リスク & 中身 \\
\midrule
市場リスク & 価格・金利・為替・変動率が不利に動く。 \\
ファクターのリスク & 割安・モメンタムなど、持っているスタイルがまとめて損をする。 \\
流動性のリスク & 売りたいときに売れない、売ると価格が大きく動く（市場の流動性）。 \\
資金繰りのリスク & 借入れの担保の追加を求められる、貸し手が資金を引き上げる、投資家が解約する（資金の流動性）。 \\
信用・取引相手のリスク & 債券の発行体やデリバティブの相手、プライムブローカーが破綻する。 \\
モデルのリスク & モデルの前提（相関、分布、流動性）が崩れる。 \\
混雑のリスク & 同じ取引をしている人が一斉に手じまう。 \\
オペレーションのリスク & システムの誤り、発注の誤り、不正。 \\
\bottomrule
\end{tabular}
\end{table}

市場の流動性と資金の流動性は、互いに強め合う（Brunnermeier・Pedersen 2009 \cite{brunnermeierpedersen2009}）。価格が下がると担保の価値が下がり、担保の追加を求められて売る。売ると価格がさらに下がる。この「流動性のらせん」が、ヘッジファンドの大きな損失の多くの背景にある（第 II 部、第 \ref{ch:failures} 章）。

\section{リスクの物差し}

\subsection{ボラティリティ}

最も基本的な物差しは、リターンの標準偏差（ボラティリティ）である。ポートフォリオのボラティリティは $\sigma_p = \sqrt{\bm w\T\Sigma\bm w}$。

\subsection{VaR}

\textbf{VaR}（バリュー・アット・リスク）は、「ある期間に、ある確率（信頼水準）で超えない損失」である。例えば「1 日の 99\% の VaR が 1 億円」とは、100 日に 1 日程度、1 億円を超える損失が出るという意味である。

\begin{itemize}
\item \textbf{正規分布を仮定}：$\mathrm{VaR}_\alpha = z_\alpha\,\sigma_p\,\sqrt{h} \times A$（$A$ は資産額、$z_{0.99} \approx 2.33$、$h$ は日数）。
\item \textbf{過去のデータ（ヒストリカル）}：過去の各日のリターンを、いまのポートフォリオに当てはめ、損失の分布の分位点をとる。
\item \textbf{モンテカルロ}：リスクの要因を確率モデルで多数生成し、損失の分布を作る。オプションのような非線形な資産に向く。
\end{itemize}

\subsection{期待ショートフォール}

VaR は「どれだけ悪いことが起きうるか」の境目しか見ず、その先の損失の大きさを無視する。\textbf{期待ショートフォール}（ES、CVaR）は、VaR を超えた日の損失の平均である。正規分布なら
\begin{equation}
\mathrm{ES}_\alpha = \sigma_p\,\frac{\phi(z_\alpha)}{1-\alpha}
\end{equation}
（$\phi$ は標準正規の密度）。裾の厚い分布では、ES のほうがリスクをよく表す。銀行の規制も、VaR から ES に移った。

\begin{exbox}[VaR の計算]
資産 100 億円、年率のボラティリティ 10\% のポートフォリオ。1 日のボラティリティは $10\% / \sqrt{252} \approx 0.63\%$。
\begin{itemize}
\item 1 日の 99\% の VaR：$2.33 \times 0.63\% \times 100 \approx 1.47$ 億円。
\item 1 日の 99\% の ES：$0.63\% \times \phi(2.33)/0.01 \approx 0.63\% \times 2.67 \approx 1.68$ 億円。
\item 10 日の 99\% の VaR：$1.47 \times \sqrt{10} \approx 4.6$ 億円（独立を仮定）。
\end{itemize}
実際の株のリターンは裾が厚く、99\% を超える損失は、正規分布の想定よりずっと大きくなりやすい。
\end{exbox}

\subsection{ストレス・テスト}

過去の危機（1987 年のブラックマンデー、1998 年、2008 年、2020 年 3 月）や、仮想の状況（金利が 2\% 上がり、株が 20\% 下がる）を、いまのポートフォリオに当てはめて損失を計算する。統計の物差しでは見えない「相関が崩れたとき」の損失を見るのに欠かせない。

\section{リスクモデル}

\subsection{なぜ必要か}

$N$ 銘柄の共分散 $\Sigma$ には、$N(N+1)/2$ 個の数がある。3,000 銘柄なら約 450 万個。日次のデータが 1 年（252 日）しかなければ、標本の共分散は特異（逆行列がない）になり、推定の誤差が非常に大きい。そこで、構造を仮定して推定する数を減らす。

\subsection{ファクター型のリスクモデル}

銘柄のリターンを、少数（数十）の共通のファクターと、銘柄固有の部分に分ける。
\begin{equation}
\bm r_t = B\,\bm f_t + \bm u_t, \qquad \Sigma = B F B\T + D
\end{equation}
$B$（$N \times K$）は各銘柄のファクターへの感応度、$F$（$K \times K$）はファクターの共分散、$D$ は銘柄固有の分散の対角行列。推定する数は $NK + K(K+1)/2 + N$ に減る。

\begin{itemize}
\item \textbf{ファンダメンタル型}（Barra 型）：感応度 $B$ を、業種のダミー、規模・割安・モメンタム・ボラティリティなどの銘柄の特性で直接与え、毎期の横断面の回帰でファクターのリターン $\bm f_t$ を推定する。商用のリスクモデルの多くはこの形。
\item \textbf{統計型}：主成分分析でファクターを取り出す。解釈は難しいが、データだけで作れる。
\item \textbf{マクロ型}：金利・インフレ・景気などの経済変数をファクターにし、時系列の回帰で感応度を推定する。
\end{itemize}

ファクター型のモデルは、ポートフォリオのリスクを「どのファクターから来ているか」に分解できる。マルチマネージャーが、運用チームにファクターのリスクを打ち消させるときに使うのも、これである。

\subsection{縮小推定}

Ledoit・Wolf（2004）\cite{ledoitwolf2004} は、標本の共分散 $S$ を、構造を持った目標 $F$（すべての相関が等しい行列など）に向けて縮めた
\begin{equation}
\hat\Sigma = \alpha F + (1-\alpha) S
\end{equation}
を提案し、平均二乗誤差を最小にする $\alpha$ をデータから計算する式を示した。計算が簡単で頑健なので、広く使われている。LLM-BLM の BLM でも、直近 252 日の Ledoit--Wolf の推定を使った。

\subsection{ランダム行列による雑音の除去}

$N$ 銘柄、$T$ 日の互いに独立な雑音から作った相関行列の固有値は、$q = N/T$ のとき、Marchenko--Pastur 分布に従い、ほぼ
\begin{equation}
\lambda_\pm = \big(1 \pm \sqrt{q}\big)^2
\end{equation}
の範囲に収まる。実際の株の相関行列の固有値のうち、この範囲を超える少数（市場全体、業種）だけが本当の構造で、範囲の中は雑音と見なせる。範囲の中の固有値を平均値に置き換えることで、相関行列を「掃除」する（Laloux ほか 1999 \cite{laloux1999}、Plerou ほか 1999 \cite{plerou1999} の観察に基づく）。

\section{ボラティリティの予測}

リターンの平均の予測は難しいが、ボラティリティは予測できる。大きな変動の後には大きな変動が続く（ボラティリティのクラスタリング）からである。

\subsection{EWMA}

指数加重の移動平均で分散を更新する。
\begin{equation}
\hat\sigma_t^2 = \lambda\,\hat\sigma_{t-1}^2 + (1-\lambda)\,r_{t-1}^2
\end{equation}
RiskMetrics（1994）は日次で $\lambda = 0.94$ を使った。

\subsection{GARCH}

Engle（1982）\cite{engle1982} の ARCH を一般化した Bollerslev（1986）\cite{bollerslev1986} の GARCH(1,1) は
\begin{equation}
\sigma_t^2 = \omega + a\,r_{t-1}^2 + b\,\sigma_{t-1}^2, \qquad a + b < 1
\end{equation}
で、長期の平均 $\omega/(1-a-b)$ に戻る力を持つ。株では、下げた後に変動が大きくなる非対称（レバレッジ効果）を入れた GJR-GARCH や EGARCH もよく使う。

\subsection{実現ボラティリティ}

日中の 5 分ごとなどの高頻度のリターンの二乗和で、その日の分散を直接測る（実現分散）。Corsi（2009）\cite{corsi2009} の HAR モデルは、日・週・月の実現分散の平均の線形回帰で、翌日の実現分散を予測する。簡単だが強い。

\section{リスクの管理の仕組み}

\subsection{リスクの上限}

運用会社は、次のような上限を置き、超えたら機械的にポジションを減らす。

\begin{itemize}
\item ポートフォリオ全体の VaR・ボラティリティの上限。
\item グロス・ネットのレバレッジの上限。
\item 1 銘柄・1 業種・1 国への集中の上限。
\item ファクターへの感応度の上限。
\item 売買高に対するポジションの大きさの上限（例：1 日の売買高の数日分で手じまえる大きさ）。
\item 損失の上限（ストップロス、ドローダウンの上限）。
\end{itemize}

\subsection{ドローダウンの管理}

運用チームの損失が一定に達したら資金の割り当てを減らし、さらに損をしたらポジションを閉じる、という規則は、マルチマネージャーの標準である。これは、損失が続く運用者を早く止める（腕がないか、市場が変わった可能性が高い）とともに、会社全体の損失を限る。一方で、多くのファンドが同じ規則で同時に損切りすると、損切りの売りが市場を動かす。

\subsection{リスクの予算}

全体のリスク（例：年率 10\% のボラティリティ）を、戦略・運用チーム・ファクターに配分する。期待される情報レシオの高いところに多く配る。リスクパリティ（\chapref{ch:pc}）は、期待を使わずに等しく配る極端な場合である。

\section{リスクモデルの限界}

\begin{itemize}
\item \textbf{危機には相関が上がる}：普段は低い相関の資産が、危機には一斉に同じ方向に動く。分散の効果が、最も必要なときに消える。
\item \textbf{過去の窓の長さ}：過去 1 年の平穏な期間で推定したリスクは、次の危機のリスクを過小評価する。危機の直後は逆に過大評価する。
\item \textbf{混雑は見えない}：ほかの運用者が何を持っているかは、価格のデータからは分かりにくい。
\item \textbf{流動性は見えない}：平常時の売買高は、危機のときに売れる量を表さない。
\end{itemize}

\begin{projbox}
LLM-BLM のポートフォリオは、ロングのみで、事前の TE（リスクモデルで予測した、比べる相手との差のボラティリティ）に上限を置いた。小型株の S1 の追記で、事前の TE の上限を年 3\% のまま「校正しない」としたのは、小型株では事前の TE と実現の TE の比を測った過去がないからである。リスクモデルの予測した TE と、実際に実現した TE を比べることは、それ自体がリスクモデルの検証になる。
\end{projbox}

\chapter{取引コストと執行}
\label{ch:cost}

論文の戦略の多くは、取引コストを引く前の成績で示される。実際の運用では、コストが予測の価値の大半を食うことが珍しくない。クオンツの運用会社の競争力の大きな部分は、予測の質よりも、コストを測り、減らす技術にある。

\section{コストの内訳}

\begin{table}[H]
\centering\small
\caption{取引コストの内訳}
\label{tab:costs}
\begin{tabular}{@{}P{0.24\linewidth}P{0.46\linewidth}P{0.22\linewidth}@{}}
\toprule
コスト & 中身 & 見え方 \\
\midrule
手数料・税 & 証券会社への手数料、取引所の手数料、取引税 & 明示的 \\
スプレッド & 買値と売値の差の半分（すぐに売買すると払う） & 半ば明示的 \\
マーケット・インパクト & 自分の売買が価格を不利に動かす分 & 暗黙 \\
遅延 & 判断してから売買するまでの価格の変化 & 暗黙 \\
機会の費用 & 売買しきれなかった分の、逃した利益 & 暗黙 \\
空売りの費用 & 借りた株の貸株料、配当の支払い、リコールのリスク & 明示的 \\
資金の費用 & 借入れの金利、証拠金に寝る資金 & 明示的 \\
\bottomrule
\end{tabular}
\end{table}

大型株を少し売買するなら、コストの大部分はスプレッドで、片道 0.01〜0.05\% 程度である。小型株や大口の売買では、マーケット・インパクトが支配的になり、片道 0.5\% を超えることも珍しくない。

\section{インプリメンテーション・ショートフォール}

Perold（1988）\cite{perold1988} は、コストを「売買を決めた時点の価格で、すべて売買できたとした紙の上のポートフォリオ」と「実際のポートフォリオ」の成績の差として測ることを提案した（インプリメンテーション・ショートフォール）。
\begin{equation}
\mathrm{IS} = \underbrace{\sum_j q_j\,(p_j - p_0)}_{\text{実際に売買した分のコスト}} + \underbrace{(Q - \textstyle\sum_j q_j)(p_T - p_0)}_{\text{売買しなかった分の機会の費用}} + \text{手数料}
\end{equation}
（買いの場合。$p_0$ は判断の時点の価格、$q_j$, $p_j$ は各約定の数量と価格、$Q$ は目標の数量）。この物差しは、手数料・スプレッド・インパクト・遅延・機会の費用をすべて含む。

\section{マーケット・インパクト}

\subsection{一時的なインパクトと恒久的なインパクト}

大口の注文を出すと、価格は不利に動く。
\begin{itemize}
\item \textbf{一時的なインパクト}：板の流動性を消費したことによる一時的な価格の動き。売買を終えると、ある程度戻る。
\item \textbf{恒久的なインパクト}：売買が情報を含むと市場に解釈され、価格の水準が変わる分。戻らない。
\end{itemize}

Kyle（1985）\cite{kyle1985} のモデルでは、情報を持つ投資家の注文の流れに対して、価格は線形に動く（$\Delta p = \lambda\,Q$、$\lambda$ は「Kyle のラムダ」）。

\subsection{平方根則}

実際の大口の注文（何日かに分けて執行する「メタオーダー」）のデータでは、インパクトは注文の大きさの平方根にほぼ比例することが、多くの市場・時代で観察されている \cite{athl2005,toth2011,bouchaud2018}。
\begin{equation}
I(Q) \approx Y\,\sigma\,\sqrt{\frac{Q}{V}}
\label{eq:sqrt}
\end{equation}
$\sigma$ は日次のボラティリティ、$V$ は 1 日の売買高、$Q$ は注文の量、$Y$ は 1 前後の定数。インパクトは、取引額そのものではなく、\textbf{売買高に対する割合}と、\textbf{ボラティリティ}で決まる。

\begin{exbox}[平方根則で見積もる]
\begin{itemize}
\item 大型株：日次のボラティリティ 1.5\%、1 日の売買高 5 億ドル。100 万ドルの買い（売買高の 0.2\%）：$1.5\% \times \sqrt{0.002} \approx 0.07\%$。
\item 小型株：日次のボラティリティ 3\%、1 日の売買高 500 万ドル。同じ 100 万ドルの買い（売買高の 20\%）：$3\% \times \sqrt{0.2} \approx 1.3\%$。
\end{itemize}
同じ金額でも、小型株ではインパクトが 20 倍近くになる。小型株の戦略が「論文では儲かるのに、大きな資金では儲からない」理由の一つである。
\end{exbox}

\section{最適執行}

\subsection{Almgren--Chriss のモデル}

大口の注文を、どのくらいの速さで執行するか。速く売買するとインパクトが大きく、ゆっくり売買すると、その間の価格の変動のリスクを負う。Almgren・Chriss（2000）\cite{almgrenchriss2000} は、このトレードオフを平均分散の問題にした。

時間 $[0, T]$ で $X$ 株を売る。残りの株数を $x(t)$、売る速さを $v = -\dot x$ とし、一時的なインパクトが速さに比例する（$\eta v$）とすると、コストの期待値とリスクは
\begin{equation}
\E[\text{コスト}] = \int_0^T \eta\,v(t)^2\,dt + (\text{恒久的なインパクトの項}), \qquad
\Var[\text{コスト}] = \sigma^2\int_0^T x(t)^2\,dt
\end{equation}
で、$\E + \lambda\Var$ を最小にする執行の軌跡は
\begin{equation}
x(t) = X\,\frac{\sinh\big(\kappa(T-t)\big)}{\sinh(\kappa T)}, \qquad \kappa = \sqrt{\frac{\lambda\sigma^2}{\eta}}
\end{equation}
となる。リスクを嫌うほど（$\lambda$ が大きいほど）、最初に速く売る。$\lambda \to 0$ なら一定の速さ（TWAP）になる。

\subsection{執行のアルゴリズム}

実務の執行アルゴリズムは、証券会社や運用会社が提供する。

\begin{itemize}
\item \textbf{VWAP}：その日の売買高の時間ごとの分布に合わせて執行し、その日の売買高加重の平均価格に近づける。
\item \textbf{TWAP}：時間で均等に分ける。
\item \textbf{POV}（参加率）：市場の売買高の一定の割合（例：10\%）で執行する。
\item \textbf{IS}（ショートフォール）：判断の時点の価格からの差を最小にする（Almgren--Chriss の考え方）。
\item \textbf{引けの売買}：取引所の終値を決める板寄せ（クロージング・オークション）で売買する。指数のファンドの売買が集中し、流動性が大きい。
\item \textbf{ダークプール・スマートオーダールーティング}：公開されない板や複数の取引所に注文を振り分け、情報の漏れとインパクトを減らす。
\end{itemize}

\section{コストを入れた最適化}

ポートフォリオの最適化（\chapref{ch:pc}）に、コストを直接入れる。
\begin{equation}
\max_{\bm w}\ \bm w\T\bm\alpha - \frac{\lambda}{2}\bm w\T\Sigma\bm w - \sum_i \Big( c_i\,|w_i - w_i^0| + k_i\,|w_i - w_i^0|^{3/2} \Big)
\end{equation}
第 1 の項は比例するコスト（スプレッド・手数料）、第 2 の項は平方根則のインパクト（取引額 $\times$ 平方根 $=$ 1.5 乗）。比例するコストがあると、予測が少し変わっただけでは売買しない「売買しない領域」ができる。これで回転率が大きく下がり、コストの後の成績はむしろ良くなる。

\section{容量}

戦略の\textbf{容量}（キャパシティ）は、コストの後の利益が最大になる、あるいはゼロになる運用資金の大きさである。資金 $A$ の戦略で、取引額は $A$ に比例し、平方根則のインパクトの率は $\sqrt{A}$ に比例するので、
\begin{equation}
\text{コストの後の利益} \approx \alpha A - c\,A^{3/2}
\end{equation}
で、利益は $A^* = (2\alpha / 3c)^2$ で最大になり、$A = (\alpha/c)^2$ でゼロになる。予測の減衰が速い（回転率が高い）戦略、小型株を使う戦略ほど、$c$ が大きく容量は小さい。ヘッジファンドが、成績の良い戦略で新しい資金を断ること（ファンドを閉じる）があるのはこのためである。

\section{空売り}

\begin{itemize}
\item \textbf{ロケート}：空売りの前に、借りられる株を証券会社に確保してもらう。
\item \textbf{貸株料}：大半の銘柄では年 0.3\% 程度以下だが、借りたい人の多い銘柄（小型株、悪材料の銘柄）は年数十\% になることもある。空売りで儲かりそうな銘柄ほど借りるのが高い。
\item \textbf{リコールと踏み上げ}：貸し手が株を返せと言えば、空売りを閉じなければならない。空売りの多い銘柄が急に上がると、損失を限るための買い戻しが殺到し、さらに上がる（ショート・スクイーズ）。2021 年 1 月の GameStop が代表例である。
\item \textbf{規制}：危機の時に、空売りが一時的に禁止されることがある（2008 年の金融株など）。
\end{itemize}

\begin{projbox}
LLM-BLM では、片道 0.3\% のコストを、BLM と比べる相手の両方から引いた。小型株の S1 の運用では、上位 20\% を浮動株時価総額で加重し、毎月入れ替える。平方根則で考えると、小さな資金なら 0.3\% は妥当な水準だが、大きな資金では明らかに足りない。逆に、実際に売買する個人や研究の規模では、容量は問題にならない。研究の結論が「この規模の資金で」のものであることを明記する必要がある。
\end{projbox}

\chapter{バックテストと過学習}
\label{ch:backtest}

\textbf{バックテスト}とは、戦略を過去のデータに当てはめ、もしその戦略で運用していたらどうなったかを計算することである。新しい戦略の検証のほとんどは、バックテストに頼るしかない。しかし、バックテストの成績は、ほぼ必ず、実際の運用の成績より良く出る。この章は、その理由と対策をまとめる。

\section{バックテストの偏り}

\begin{table}[H]
\centering\small
\caption{バックテストの成績を良く見せる偏り}
\label{tab:biases}
\begin{tabular}{@{}P{0.22\linewidth}P{0.42\linewidth}P{0.28\linewidth}@{}}
\toprule
偏り & 中身 & 対策 \\
\midrule
先読み & その時点で知られていなかった情報を使う（決算の公表日の前に決算の値を使う、後で訂正された値を使う、言語モデルが学習した後の出来事を知っている） & ポイント・イン・タイムのデータ、公表の遅れ、学習の期間の止まった言語モデル \\
生存者 & 今日存在する銘柄・ファンドだけで過去を検証する & 上場廃止・倒産した銘柄を含むデータ \\
データ・スヌーピング & 多くの戦略・パラメータを試し、最も良いものを報告する & 試した数を記録し、多重検定の補正（\ref{sec:mt} 節） \\
コストの無視 & 取引コスト・空売りの費用・インパクトを入れない & \chapref{ch:cost} のコストのモデル \\
実行できない売買 & 終値で判断し同じ終値で売買する、売買の止まった銘柄を売買する、借りられない株を空売りする & 判断と売買の時点をずらす、流動性の条件 \\
期間の選び方 & 戦略に都合の良い期間だけを使う & できるだけ長い期間、複数の市場 \\
\bottomrule
\end{tabular}
\end{table}

\section{過学習}

\subsection{多くを試すと、偶然に良いものが見つかる}

本当は予測力がまったくない（真のシャープレシオがゼロの）戦略を $N$ 個試すと、その中の最大のシャープレシオは、$N$ とともに大きくなる。Bailey・López de Prado（2014）\cite{dsr2014} によれば、試した戦略のシャープレシオの推定値のばらつきを $\sigma_{\SR}$ とすると、最大値の期待値はおよそ
\begin{equation}
\E[\max_n \widehat{\SR}_n] \approx \sigma_{\SR}\Big[(1-\gamma)\,\Phi^{-1}\!\Big(1 - \frac{1}{N}\Big) + \gamma\,\Phi^{-1}\!\Big(1 - \frac{1}{N e}\Big)\Big]
\end{equation}
である（$\gamma \approx 0.5772$ はオイラーの定数）。

\begin{exbox}[偶然でどれだけ良く見えるか]
5 年の日次のデータで、真のシャープレシオがゼロの戦略を試す。年率のシャープレシオの推定値の標準誤差は約 $1/\sqrt{5} \approx 0.45$。
\begin{center}
\begin{tabular}{@{}lcccc@{}}
\toprule
試した数 $N$ & 10 & 100 & 1{,}000 & 10{,}000 \\
\midrule
最大の SR の期待値（年率） & 0.70 & 1.13 & 1.46 & 1.73 \\
\bottomrule
\end{tabular}
\end{center}
100 通りのパラメータを試せば、まったく無意味な戦略から、年率 1 を超えるシャープレシオが「見つかる」。Bailey・Borwein・López de Prado・Zhu（2014）\cite{charlatanism2014} は、試した数を示さずにバックテストの成績を示すことを、「疑似数学と金融のいかさま」と呼んで警告した。
\end{exbox}

\subsection{デフレートしたシャープレシオ}

Bailey・López de Prado（2014）の\textbf{デフレートしたシャープレシオ}（DSR）は、観測したシャープレシオ $\widehat\SR$ が、「$N$ 回試したときに偶然出る最大値」$\SR_0$ を超える確率を、リターンの歪度 $\gamma_3$・尖度 $\gamma_4$ と標本の長さ $T$ も考慮して計算する。
\begin{equation}
\mathrm{DSR} = \Phi\!\left(\frac{(\widehat\SR - \SR_0)\sqrt{T-1}}{\sqrt{1 - \gamma_3\,\widehat\SR + \frac{\gamma_4 - 1}{4}\widehat\SR^2}}\right)
\end{equation}
（$\widehat\SR$ と $\SR_0$ は期間あたりの値。$\SR_0$ は上の最大値の期待値。）これが 0.95 を超えなければ、偶然と区別できない。重要なのは、\textbf{試した数 $N$ を記録しておかなければ、この補正ができない}ことである。

\subsection{バックテストの過学習の確率}

Bailey・Borwein・López de Prado・Zhu（2017）\cite{pbo2017} の\textbf{PBO}（probability of backtest overfitting）は、次の手順で計算する（組み合わせ対称の交差検証、CSCV）。

\begin{enumerate}
\item 期間を $S$ 個（例えば 16）の塊に分ける。
\item そのうち半分を学習、残りの半分を評価とする分け方を、すべての組み合わせ（$\binom{16}{8} = 12{,}870$ 通り）で作る。
\item 各分け方で、試したすべての戦略のうち、学習の期間で最も良いものを選び、それが評価の期間で何番目かを見る。
\item 評価の期間で中央値より下になった割合を PBO とする。
\end{enumerate}
PBO が 0.5 に近いなら、「学習の期間で最も良い」ことに、ほとんど意味がない。

\section{多重検定}
\label{sec:mt}

\subsection{補正の方法}

$m$ 個の仮説を同時に検定するとき、それぞれを 5\% で検定すると、少なくとも一つが偶然に有意になる確率は $1 - 0.95^m$（$m = 20$ なら 64\%）になる。

\begin{itemize}
\item \textbf{Bonferroni}：各検定を $\alpha/m$ で行う。最も厳しい。
\item \textbf{Holm}\cite{holm1979}：$p$ 値を小さい順に並べ、$k$ 番目を $\alpha/(m-k+1)$ と比べ、初めて超えたところで止める。Bonferroni より検出力が高く、同じく「一つでも誤る確率」（FWER）を $\alpha$ 以下に抑える。
\item \textbf{Benjamini--Hochberg}\cite{bh1995}：有意とした仮説のうち、誤りの割合の期待値（FDR）を抑える。多くの候補から有望なものを拾うときに向く。
\item \textbf{White の Reality Check}（White 2000 \cite{white2000}）、\textbf{Hansen の SPA}（Hansen 2005 \cite{hansen2005}）：多くの戦略の中で最も良いものが、基準（例：市場）に本当に勝っているかを、ブートストラップで検定する。
\end{itemize}

\subsection{\texorpdfstring{$t > 3$}{t > 3}}

Harvey・Liu・Zhu（2016）\cite{hlz2016} は、学術誌に発表された数百のファクターについて、研究者の全体が試した数を考慮すると、新しいファクターの $t$ 値の基準は 2 ではなく 3 程度にすべきだと論じた。業界でも、新しいシグナルの採用の基準として広く引用される。

\section{検証の設計}

\subsection{手元に残しておく}

最も強い対策は、\textbf{一度も見ていないデータ}で、\textbf{一度だけ}評価することである。

\begin{enumerate}
\item データを、開発の期間と、最後まで見ない期間（ホールドアウト）に分ける。
\item 開発の期間の中で、ウォークフォワードと交差検証（\ref{sec:cv} 節）で戦略を作る。
\item 戦略を固定してから、ホールドアウトで一度だけ評価する。ホールドアウトの結果を見て戦略を直したら、それはもうホールドアウトではない。
\end{enumerate}

\subsection{事前登録}

医学の臨床試験と同じく、検定の仮説・データ・変数・方法・判定の基準を、結果を見る前に文書にして固定する（事前登録）。事前登録にない分析は「探索」として区別する。分析の選択の自由（どの期間、どの統制、どの外れ値の処理）による偶然の発見を防ぐ。

\subsection{紙の上の運用と小さな実運用}

バックテストの後、実際の価格で売買の注文を（出さずに）記録する「紙の上の運用」や、小さな資金での実運用を、数か月〜数年行う。実際の売買のコスト、データの遅れ、システムの誤りは、ここで初めて分かる。

\section{バックテストの確認の一覧}

\begin{keybox}[バックテストを信じる前に確かめること]
\begin{enumerate}
\item 各時点で、その時点で知られていた情報だけを使っているか（公表日、訂正、言語モデルの学習の期間）。
\item 上場廃止・倒産した銘柄が含まれているか。
\item 取引コスト・空売りの費用・インパクトを入れたか。容量はどれくらいか。
\item 何通りの戦略・パラメータを試したか。それを補正したか。
\item 期間を分けても（前半と後半、好況と不況）、市場を変えても、同じ結果か。
\item 結果に経済的な説明があるか（なぜ儲かるのか、誰が損をしているのか）。
\item 少数の銘柄・少数の日が、収益の大半を生んでいないか。
\item 最も良いパラメータの周りで、成績がなだらかに変わるか（鋭い山なら過学習）。
\end{enumerate}
\end{keybox}

\begin{projbox}
LLM-BLM の小型株の実験 S1 は、検定の仮説・変数・判定の基準・コードを、株価の反応を一度も計算する前にコミットして固定し（事前登録）、1 回だけ実行する設計にした。副次の検定には Holm の補正をかけ、事前登録にない分析は「探索」として別に報告する。LLM の点数は、学習の期間が止まった言語モデル（ChronoGPT の 2015 年版）を主に使い、言語モデルが将来を知っている先読みを避けた。これらは、この章の対策をそのまま当てはめたものである。
\end{projbox}

\chapter{マーケットメイクと高頻度取引}
\label{ch:mm}

ここまでの章は、数日〜数か月の予測を前提にしていた。この章は、秒以下の世界、つまり取引所の板の上で何が起きているかを扱う。ここで使う道具（板、スプレッド、逆選択、在庫）は、長い期間の運用の取引コスト（\chapref{ch:cost}）を理解するうえでも基本になる。

\section{板の仕組み}

\subsection{指値の注文と成行の注文}

多くの取引所は、\textbf{指値の注文の板}（リミット・オーダー・ブック）で売買を成立させる。

\begin{itemize}
\item \textbf{指値の注文}：「この価格以下でなら 100 株買う」。すぐには成立せず、板に並んで待つ。流動性を\textbf{供給}する。
\item \textbf{成行の注文}：「いますぐ 100 株買う」。板に並んでいる最も安い売りの指値と成立する。流動性を\textbf{消費}する。
\item \textbf{最良気配}：最も高い買いの指値（ビッド）と、最も安い売りの指値（アスク）。その差が\textbf{スプレッド}、中間が\textbf{仲値}。
\item \textbf{優先の順位}：同じ価格の指値の中では、先に出した注文から成立する（価格優先・時間優先）。良い順番を取るために速さが要る。
\item \textbf{ティック}：価格の最小の刻み。米国株では原則 1 セント。
\end{itemize}

\begin{figure}[H]
\centering
\begin{tikzpicture}[x=7mm,y=5mm]
\foreach \p/\q/\v in {1/3/99.95,2/5/99.96,3/2/99.97,4/6/99.98,5/4/99.99} {
  \fill[c3!55] ({-\q},{\p-0.4}) rectangle (0,{\p+0.4});
  \node[lab,anchor=west] at (0.15,\p) {\v};
}
\foreach \p/\q in {7/2,8/4,9/3,10/6,11/5} {
  \fill[c2!55] (0,{\p-0.4}) rectangle (\q,{\p+0.4});
  \node[lab,anchor=east] at (-0.15,\p) {100.0\the\numexpr \p-6\relax};
}
\draw[ink2,dashed] (-7,6) -- (7,6);
\node[lab,anchor=west] at (2.4,6.35) {スプレッド（100.01 − 99.99 = 0.02）};
\node[font=\footnotesize,c3!70!black] at (-4.5,0) {買いの指値（ビッド）};
\node[font=\footnotesize,c2!80!black] at (4.5,12.2) {売りの指値（アスク）};
\draw[arr] (0,-0.8)--(0,12.8) node[above,font=\scriptsize]{価格};
\end{tikzpicture}
\caption{指値の注文の板の模式図。棒の長さは、その価格に並んでいる注文の量。成行の買いは、最良の売り（100.01）から順に成立する。}
\label{fig:lob}
\end{figure}

\subsection{取引所の手数料の仕組み}

米国の多くの取引所は、流動性を供給した側（メイカー）に報奨金を払い、消費した側（テイカー）から手数料をとる（メイカー・テイカー方式）。マーケットメイカーの収益には、この報奨金も含まれる。

\section{マーケットメイクの経済学}

\subsection{収益と損失}

マーケットメイカーは、ビッドとアスクの両方に指値を出し続ける。買いと売りが同じ量ずつ来れば、1 往復ごとにスプレッドを稼げる。しかし、損失の源が二つある。

\begin{enumerate}
\item \textbf{逆選択}：相手が、自分の知らない情報を持っていることがある。良い知らせを知る相手は買い、マーケットメイカーは売ったあとで価格が上がる。
\item \textbf{在庫のリスク}：買いと売りが偏ると、在庫（ポジション）がたまり、その間の価格の変動のリスクを負う。
\end{enumerate}

\subsection{逆選択のモデル}

Glosten・Milgrom（1985）\cite{glostenmilgrom1985} のモデルでは、注文の一部は情報を持った投資家から、残りは情報を持たない投資家（流動性の必要から売買する）から来る。マーケットメイカーは、相手が誰かは分からないが、「買い注文が来た」こと自体が、価格が上がる可能性を示すと考えて、アスクを高めに、ビッドを低めに置く。スプレッドは、情報を持った投資家に負ける分を、持たない投資家から取り返す大きさに決まる。情報を持った投資家が多い銘柄（小型株、決算の前後）ほど、スプレッドは広い。

Kyle（1985）\cite{kyle1985} のモデルでは、情報を持った投資家は、自分の情報が価格に漏れすぎないように、少しずつ売買する。価格は、注文の流れ（買い − 売り）に比例して動く（$\Delta p = \lambda \times$ 注文の流れ）。$\lambda$ が大きい銘柄ほど、売買の価格への影響（インパクト）が大きい。

\subsection{在庫のモデル}

Ho・Stoll（1981）\cite{hostoll1981} と Avellaneda・Stoikov（2008）\cite{avellanedastoikov2008} は、在庫のリスクを考えて、指値の位置を決めるモデルを作った。Avellaneda・Stoikov のモデルでは、仲値を $s$、在庫を $q$（買い越しなら正）、リスクを嫌う度合いを $\gamma$、価格の変動を $\sigma$、残り時間を $T-t$ として、マーケットメイカーにとっての「本当の価格」（留保価格）は
\begin{equation}
r = s - q\,\gamma\,\sigma^2\,(T - t)
\end{equation}
である。在庫を多く持つほど、留保価格を下げ、ビッドもアスクも下に動かして、売りが成立しやすく、買いが成立しにくくする。こうして、在庫をゼロに戻す。スプレッドの幅は、リスクを嫌う度合いと、注文の来る頻度で決まる。

\section{高頻度取引}

\subsection{種類}

高頻度取引（HFT）は、計算機で秒以下の速さで注文を出し、取り消し、多くの売買をする取引の総称である。主な戦略は次のとおり。

\begin{itemize}
\item \textbf{電子的なマーケットメイク}：上で述べた、流動性の供給。HFT の最大の部分。
\item \textbf{遅延の裁定}：同じ資産（例：株価指数の先物と、その現物の ETF）が別の市場で取引されているとき、一方の価格の動きをもう一方より速く知って売買する。
\item \textbf{短期の統計的裁定}：板の偏り（買いの注文が売りより多い）などから、数秒〜数分の価格の動きを予測する。Cont・Kukanov・Stoikov（2014）\cite{cks2014} は、最良気配の注文の増減の偏り（注文の流れの不均衡）が、短い時間の価格の変化をよく説明することを示した。
\item \textbf{イベントの裁定}：経済指標や企業の発表を、機械で読んで瞬時に売買する。
\end{itemize}

\subsection{速さの競争}

取引所の計算機と同じ建物に自社の計算機を置く（コロケーション）、都市の間をマイクロ波の通信で結ぶ、注文の処理を専用の回路（FPGA）でする、などの投資で、何マイクロ秒を争う。Budish・Cramton・Shim（2015）\cite{bcs2015} は、連続的に売買を成立させる仕組みそのものが、この「速さの軍拡競争」を生むと論じ、一定の間隔（例：0.1 秒）ごとにまとめて売買を成立させる方式（頻繁なバッチ・オークション）を提案した。

\subsection{HFT の功罪}

\begin{itemize}
\item \textbf{良い面}：スプレッドを狭め、流動性を増やし、価格の発見を速くした、という実証が多い。
\item \textbf{悪い面}：2010 年 5 月 6 日の米国株の急落（フラッシュ・クラッシュ）のように、市場が混乱すると、HFT は指値を取り消して退き、流動性が一瞬で消えることがある。遅延の裁定は、遅い参加者から速い参加者へ利益を移すだけで、社会的な価値は小さいとも論じられる。
\end{itemize}

\section{個人や研究で扱えるか}

マーケットメイクと HFT は、ヘッジファンドの戦略の中で、個人や大学の研究が実際に競うことの最も難しい分野である。

\begin{itemize}
\item 実際の板のデータ（全注文の記録）は高価で、量が膨大である。
\item 競争の相手は、速さに巨額の投資をした専門の会社で、遅い参加者は逆選択の相手になる。
\item バックテストでは、自分の注文が板の何番目に並び、いつ成立するかを正確に再現できない。
\end{itemize}

一方で、この分野の考え方（スプレッド、逆選択、インパクト）は、遅い戦略の取引コストを見積もるうえで欠かせない。決算発表の直後の数分は、情報を持った投資家と HFT が最も活発に動く時間帯であり、LLM の読み取りで「当日に」売買しようとすると、この人たちと競うことになる。

\begin{projbox}
LLM-BLM の実験では、決算発表文の情報は発表の当日に株価に入りきった。その「入りきる」過程を担っているのが、まさにこの章の参加者（発表の文章を機械で読んで瞬時に売買する会社と、マーケットメイカー）である。学生の研究が発表の数分後に読み取りを終えても、その時点で情報の大部分は価格に入っている。研究の戦略を、「速さでは競わない」領域（数週〜数か月の、注目の低い銘柄の反応の遅れ）に置く必要があるのは、このためである。
\end{projbox}


\part{研究で分かっていること}
\chapter{調査の方法}
\label{ch:method}

第 II 部は、Gemini の Deep Research（多数のウェブ検索をして長い報告を書くエージェント）に文献を調べさせ、その答えを照合したうえで、自分の知識と合わせて整理したものである。この章では、その手順と、Deep Research の答えの質についての所見をまとめる。

\section{3 つの依頼}

同じテーマ（ヘッジファンドの投資の方法とアルゴリズムの研究）を、切り口を変えて 3 回依頼した。依頼文は \code{03\_方針と文献調査/ヘッジファンドの手法/00\_依頼*.md} にある。

\begin{table}[H]
\centering\small
\caption{Deep Research への 3 つの依頼と実行の記録（2026-10-23、モデル deep-research-preview-04-2026）}
\label{tab:requests}
\begin{tabular}{@{}P{0.2\linewidth}P{0.4\linewidth}rrr@{}}
\toprule
依頼 & 切り口 & 検索の回数 & 時間（秒） & 出力（トークン） \\
\midrule
1 戦略 & 戦略の種類ごとに、中心の論文、収益の源泉、実証の数値、崩れた例、再現の可否 & 38 & 765 & 57,094 \\
2 方法 & クオンツの運用の工程（予測、シグナルの組み合わせ、リスクモデル、構築、執行、検証、テキスト、強化学習）ごとの方法と比較の研究 & 27 & 560 & 37,670 \\
3 実証と失敗 & ヘッジファンドの成績の実証、データの偏り、手数料、混雑、有名な失敗、公表の後の減衰、小規模な研究への教訓 & 17 & 552 & 37,423 \\
\bottomrule
\end{tabular}
\end{table}

どの依頼にも、次の答え方を求めた。

\begin{itemize}
\item 日本語で答える。
\item 文献は、著者・題名・年・掲載先・DOI か URL を書く。存在を確かめられないものは挙げない。数値は、どの論文のどの値かが分かるように書く。
\item 確立した結果と、少数の研究・業界の主張を分けて書く。
\item 最後に、全文献の一覧を付ける。
\end{itemize}

背景として、LLM-BLM の研究の状況（LLM の読み取りを Black--Litterman に入れたが、時価総額加重に勝てなかった）を伝えた。

\section{照合の手順}

\begin{enumerate}
\item 答えの中の、検索ツールの転送用の URL を、元の URL に直した（\code{refcheck.py --resolve}）。
\item 答えに挙がった文献と、自分が本文で引こうとする文献を一つの一覧（\code{照合/文献の一覧.md}、約 190 本）にまとめ、題名で OpenAlex・Crossref・arXiv と照合した（\code{照合/照合の結果.txt}）。
\item 本文に書く数値のうち、確信の持てないものは、論文の要旨（Crossref・arXiv・NBER のページ）や原典の PDF で確かめた。確かめられなかった数値は書かないか、「確かめられなかった」と明記した。
\item Deep Research の答えの誤り・確かめられなかった記述を、付録~\ref{ch:verify} に一覧にした。
\end{enumerate}

\section{答えの質についての所見}

\begin{itemize}
\item \textbf{主な論文の選び方は妥当}だった。3 本の答えが挙げた中心の論文（Fama・French、Moskowitz・Ooi・Pedersen、Gu・Kelly・Xiu、McLean・Pontiff、Khandani・Lo、Ben-David ほか など）は、どれも実在し、分野の標準的な文献である。
\item \textbf{数値の多くは正しかった}が、誤りも混じっていた。例えば、Greenwood・Sammon の指数の効果（7.6\% → 0.8\%）、Ben-David ほかの実効の成功報酬（約 50\%、投資家の取り分 36 セント）、Korajczyk・Sadka の容量（約 50 億ドル）、Yang ほかのメタ分析の平均のアルファ（月 0.36\%）は要旨・原典と一致した。一方、2007 年 8 月の「25 標準偏差」を Khandani・Lo の結果としたこと、Menkveld の論文に要旨にない数値をつけたこと、Jacobs・Müller のアノマリーの数（241 を 231）などの誤りがあった。
\item \textbf{文献の一覧の質は低かった}。答え 2 の文献の一覧の番号は本文の引用の番号と対応しておらず、URL の多くは原典ではなく、解説のブログ、Wikipedia、別の論文の PDF を指していた。答え 3 では、論文の題名を誤ったもの（Sun・Wang・Zheng）、URL が別の論文を指すもの（Jorion 2000 が Amaranth の論文を、Pástor ほかが別の論文を指す）があった。
\item \textbf{学術論文でない出典}（業界の統計のサイト、修士論文、運用会社のブログ）も、「確立された定説」として混じっていた。
\item \textbf{取引できるかの区別が落ちる}ことがあった。Lopez-Lira・Tang の「約 90\% の的中」は、取引できない最初の反応についての値だが、答え 2 はその区別を書いていなかった。
\end{itemize}

これまでの LLM-BLM の文献調査（\code{03\_方針と文献調査/}）でも、GPT・Gemini・Deep Research のどれにも、番号・著者の取り違えや、要旨にない数値があった。今回も同じ傾向で、Deep Research は「何を読むべきかの地図」としては役に立つが、「数値の出典」としては照合なしに使えない。

\section{第 II 部の見取り図}

第 II 部の戦略の章の内容を、先に一つの表にまとめる。

{\footnotesize
\begin{longtable}{@{}P{0.13\linewidth}P{0.16\linewidth}P{0.23\linewidth}P{0.19\linewidth}P{0.17\linewidth}@{}}
\caption{戦略ごとの研究のまとめ}\label{tab:overview}\\
\toprule
戦略 & 収益の源泉 & 中心の証拠 & 崩れた例・弱点 & 個人・研究での再現 \\
\midrule
\endfirsthead
\toprule
戦略 & 収益の源泉 & 中心の証拠 & 崩れた例・弱点 & 個人・研究での再現 \\
\midrule
\endhead
\bottomrule
\endfoot
株式ファクター（第 \ref{ch:equityfactor} 章） & リスクと行動の偏りの両説 & Fama・French、Jegadeesh・Titman、Asness ほか（2013）、Frazzini・Pedersen & モメンタムの急落、割安の長い低迷、公表の後の減衰 & 検証は容易。小型株に依存するものはコストで消える \\
統計的裁定（第 \ref{ch:statarb} 章） & 流動性の供給 & Lehmann、Gatev ほか、Avellaneda・Lee、Nagel & 2007 年 8 月の解消の連鎖。収益の長期の低下 & 日次の検証は容易。実運用は速さとコストで難しい \\
トレンドフォロー（第 \ref{ch:trend} 章） & 反応の不足と行き過ぎ、ヘッジの需要 & Moskowitz・Ooi・Pedersen、Hurst ほか & 方向のない相場、急な反転。量の調整の寄与への批判 & 先物のデータで検証は可能。分散には資金が要る \\
キャリー（第 \ref{ch:macro} 章） & 急落のリスクの対価 & Fama（1984）、Lustig ほか、Koijen ほか & 2008 年の円の急騰などの急落 & 公開のデータで検証は可能 \\
イベント（第 \ref{ch:event} 章） & 注目の不足、裁定の限界、リスクの引き受け & Bernard・Thomas、Womack、Mitchell・Pulvino & PEAD の消滅（Martineau）、指数の効果の消滅 & 決算・推奨のデータで可能。合併裁定は案件の一覧が要る \\
レラティブ・バリュー（第 \ref{ch:rv} 章） & 流動性・資金の制約の対価 & Duarte ほか、Mitchell ほか & LTCM、2008 年、2020 年 3 月 & ほぼ不可能（レバレッジと与信が要る） \\
ボラティリティ（第 \ref{ch:rv} 章） & 変動・急落のリスクの対価 & Coval・Shumway、Bakshi・Kapadia、Carr・Wu & 2018 年 2 月 & オプションのデータがあれば検証は可能 \\
高頻度（第 \ref{ch:hft} 章） & スプレッド、速さ & Menkveld、Brogaard ほか、Budish ほか & フラッシュ・クラッシュ & 不可能 \\
\end{longtable}
}

\section{第 II 部の書き方}

\begin{itemize}
\item 戦略の章（第 \ref{ch:equityfactor}〜\ref{ch:hft} 章）は、主に答え 1 の構成に沿い、答え 1 が挙げなかった中心の論文（Jegadeesh・Titman、Gatev ほか、Avellaneda・Lee、Koijen ほか など）を自分の知識で補った。
\item 方法の章（第 \ref{ch:ml}〜\ref{ch:pcresearch} 章）は、主に答え 2 の構成に沿った。
\item 成績と失敗の章（第 \ref{ch:hfperf}・\ref{ch:failures} 章）は、主に答え 3 の構成に沿った。
\item 論文の要旨は「論文」の枠（\textsf{paperbox}）に、LLM-BLM との関係は「この研究との関係」の枠に分けた。
\end{itemize}

\chapter{株式のロング・ショートとファクター}
\label{ch:equityfactor}

株式のファクター（\chapref{ch:factor}）は、クオンツのヘッジファンドの中心の戦略である。この章では、代表的なファクターの発見の論文、その説明、崩れた例、公表の後の減衰の研究をまとめる。

\section{代表的なファクターの発見}

\begin{table}[H]
\centering\small
\caption{株式の代表的なファクターと、その発見・検証の論文}
\label{tab:factorpapers}
\begin{tabular}{@{}P{0.16\linewidth}P{0.4\linewidth}P{0.36\linewidth}@{}}
\toprule
ファクター & 中身 & 主な論文 \\
\midrule
規模 & 時価総額の小さい株は、ベータで説明できないほど平均リターンが高い。 & Banz（1981）\cite{banz1981}、Fama・French（1993）\cite{ff1993} \\
割安 & 簿価・利益・キャッシュフローに比べて株価の低い株が、高い株を上回る。 & Lakonishok・Shleifer・Vishny（1994）\cite{lsv1994}、Fama・French（1993） \\
モメンタム & 過去 3〜12 か月の勝ち組が、その後 3〜12 か月も負け組を上回る。 & Jegadeesh・Titman（1993）\cite{jt1993}、Carhart（1997）\cite{carhart1997} \\
収益性・クオリティ & 収益性が高く、安定し、負債の少ない株が上回る。 & Novy-Marx（2013）\cite{novymarx2013}、Asness・Frazzini・Pedersen（2019）\cite{qmj2019} \\
投資 & 資産を大きく増やした会社は、その後のリターンが低い。 & Fama・French（2015）\cite{ff2015}、Hou・Xue・Zhang（2015）\cite{hxz2015} \\
低リスク & ベータや変動率の低い株が、リスクに比べて高いリターンを持つ。 & Frazzini・Pedersen（2014）\cite{bab2014}、Ang ほか（2006）\cite{ahxz2006} \\
会計の質 & 利益に占める発生主義の部分（アクルーアル）が大きい会社は、その後のリターンが低い。 & Sloan（1996）\cite{sloan1996} \\
\bottomrule
\end{tabular}
\end{table}

\subsection{モメンタム}

\begin{paperbox}[Jegadeesh・Titman（1993）\cite{jt1993}]
米国株（1965〜1989 年）で、過去 3〜12 か月のリターンで並べた勝ち組を買い、負け組を売る戦略が、その後 3〜12 か月の保有で有意な正のリターンを生むことを示した。過去 6 か月で選び 6 か月持つ戦略は、年率で平均 12\% 程度の超過リターンを上げた。この収益は市場のリスクでは説明できず、保有の 2 年目以降には一部が戻る。
\end{paperbox}

モメンタムは、その後、ほぼすべての国の株式と、債券・通貨・商品でも確認された（Asness・Moskowitz・Pedersen 2013 \cite{amp2013}、\ref{sec:amp} 節）。収益の説明は、主に行動の偏りによる。

\begin{itemize}
\item \textbf{反応の不足}：新しい情報が、少しずつしか価格に入らない。Hong・Stein（1999）\cite{hongstein1999} は、情報が投資家の間にゆっくり広がる市場では、価格が遅れて反応し、それを追う投資家が後で行き過ぎを生むことをモデルにした。
\item \textbf{注目の低い銘柄で強い}：Hong・Lim・Stein（2000）\cite{hls2000} は、規模を調整しても、アナリストの少ない銘柄ほどモメンタムが強く、特に悪い知らせ（負け組）の反映が遅いことを示した（「悪い知らせはゆっくり伝わる」）。
\end{itemize}

\begin{projbox}
Hong・Stein（1999）と Hong・Lim・Stein（2000）は、LLM-BLM の小型株の実験 S1 の仮説（注目の低い小型株では、決算発表文の情報の反映が遅れ、発表後のドリフトが残る）の直接の根拠である。
\end{projbox}

\subsection{割安}

Lakonishok・Shleifer・Vishny（1994）\cite{lsv1994} は、割安株（簿価やキャッシュフローに比べて株価の低い株）が割高株を上回るのは、割安株がリスクの高い株だからではなく、投資家が過去の業績の悪さを将来に引き伸ばしすぎ（外挿）、安く売りすぎるからだと論じた。一方、Fama・French は、割安株は苦境の会社のリスクの対価だと解釈する（\chapref{ch:factor}、表~\ref{tab:riskmisp}）。

\subsection{低リスク}

\begin{paperbox}[Frazzini・Pedersen（2014）「ベータに逆らって賭ける」\cite{bab2014}]
レバレッジを使えない（あるいは制限された）投資家は、高いリターンを求めて高ベータの資産を買う。その結果、高ベータの資産は割高に、低ベータの資産は割安になる。低ベータの資産をレバレッジで買い、高ベータの資産を売る BAB（betting against beta）のファクターが、米国株・国際株・国債・社債・通貨・商品のほとんどで正のリスク調整後のリターンを持つことを示した。また、資金の制約が厳しくなると BAB は損をする。
\end{paperbox}

Ang・Hodrick・Xing・Zhang（2006）\cite{ahxz2006} は、銘柄固有の変動率の高い株の平均リターンが著しく低いことを示した（固有のボラティリティのパズル）。

\subsection{収益性とクオリティ}

Novy-Marx（2013）\cite{novymarx2013} は、売上総利益を総資産で割った値が、簿価時価比率と同じくらい強く将来のリターンを予測し、しかも割安と負の相関を持つので、両方を組み合わせると成績が大きく良くなることを示した。Asness・Frazzini・Pedersen（2019）\cite{qmj2019} は、収益性・成長・安全性・配当などでクオリティを定義し、クオリティの高い株を買い低い株を売る QMJ のファクターが、多くの国で正のリスク調整後のリターンを持つことを示した。クオリティの高い株は、その良さに比べて十分に高い価格がついていない。

\section{あらゆる市場の割安とモメンタム}
\label{sec:amp}

\begin{paperbox}[Asness・Moskowitz・Pedersen（2013）\cite{amp2013}]
米国・英国・欧州・日本の株式、国別の株価指数、国債、通貨、商品の 8 つの市場で、割安とモメンタムの両方が正のリターンを持つことを示した。
\begin{itemize}
\item 割安とモメンタムは、どの市場の中でも負の相関を持つ。両方を組み合わせると、どちらか一方よりずっと安定する。
\item 市場をまたいで、割安どうし・モメンタムどうしは正の相関を持つ。共通の要因があり、その一部は資金の流動性のリスクに関係する。
\end{itemize}
\end{paperbox}

この論文は、ファクター投資が株式だけの現象ではなく、資産クラスをまたぐ一般的な現象であることを示し、多資産のファクター運用（オルタナティブ・リスク・プレミアムの商品）の理論的な根拠になった。

\section{崩れた例}

\subsection{モメンタムの急落}

\begin{paperbox}[Daniel・Moskowitz（2016）\cite{danielmoskowitz2016}]
モメンタムは、大きな下げ相場の後、市場が急に反発する局面で、数か月のうちに大きな損失を出す（1932 年、2009 年など）。下げ相場の後の負け組は、高ベータの株に偏っていて、反発の局面で急騰するからである。こうした急落は、事前の市場の状態と変動率からある程度予測でき、それに応じてモメンタムの量を動的に調整すると、シャープレシオが大きく改善することを示した。
\end{paperbox}

Barroso・Santa-Clara（2015）\cite{barroso2015} は、モメンタムの戦略を、その過去の実現の変動率で割って量を調整する（\chapref{ch:pc} のボラティリティ・ターゲット）だけで、急落がほぼ消え、シャープレシオがほぼ 2 倍になることを示した。

\subsection{割安の長い低迷と、クオンツの冬}

割安（HML）は、2007 年ごろから 2020 年まで長く損をし、特に 2018〜2020 年に大きく下げた。Blitz（2021）\cite{blitz2021} は、この時期を「クオンツの危機」と呼び、割安だけでなく、低ボラティリティ・小型・クオリティなど、多くのファクターが同時に損をしたこと、大型で割高な成長株への資金の集中（ビッグ・グロース）がその原因であったことを分析した。過去には割安が不調のときにモメンタムが補ってきたが、この時期には補いきれなかった、という指摘もある（Deep Research の答え 3 は、これを Blitz の報告として挙げた。原典の数値は確かめていない）。

\subsection{2007 年 8 月}

株式のマーケット・ニュートラルのクオンツが一斉に損をした 2007 年 8 月の出来事は、第 \ref{ch:statarb} 章と第 \ref{ch:failures} 章で扱う。

\section{公表の後の減衰と、再現性}

\subsection{公表の後に小さくなる}

\begin{paperbox}[McLean・Pontiff（2016）\cite{mcleanpontiff2016}]
学術誌に発表された 97 のリターンの予測の変数（アノマリー）について、論文の標本の期間の後、論文の発表の前（標本外）と、発表の後のリターンを比べた。
\begin{itemize}
\item 標本外では、標本の期間の中より平均 26\% 小さい。これは、統計の偏り（データ・マイニング）の上限の目安になる。
\item 発表の後は、標本の期間の中より平均 58\% 小さい。標本外との差（約 32\%）は、発表で情報を知った投資家の裁定による、と解釈される。
\item 標本の期間の中で強かったアノマリーほど、発表の後に大きく減る。裁定の費用が高い（小型、流動性が低い、固有の変動が大きい）銘柄のアノマリーほど、減り方が小さい。
\item 発表の後、アノマリーの銘柄の売買高や空売りが増える。
\end{itemize}
\end{paperbox}

\begin{figure}[H]
\centering
\begin{tikzpicture}[x=1mm,y=0.32mm]
\foreach \y in {0,25,50,75,100} {\draw[ink2!25] (0,\y)--(96,\y); \node[lab,left] at (0,\y) {\y};}
\fill[c1] (8,0) rectangle (26,100);   \node[lab,above] at (17,100) {100};
\fill[c1!60] (38,0) rectangle (56,74); \node[lab,above] at (47,74) {74};
\fill[c2!75] (68,0) rectangle (86,42); \node[lab,above] at (77,42) {42};
\node[lab,align=center,below] at (17,-2) {論文の標本の\\期間の中};
\node[lab,align=center,below] at (47,-2) {標本の後・\\発表の前};
\node[lab,align=center,below] at (77,-2) {発表の後};
\end{tikzpicture}
\caption{McLean・Pontiff（2016）の 97 のアノマリーの平均のリターンの減衰（論文の標本の期間の中を 100 とした値。要旨の 26\%・58\% から描いた）。標本の外での減少（26\%）は統計の偏りの上限の目安、発表の後の追加の減少は裁定による、と解釈される。}
\label{fig:mp}
\end{figure}

Jacobs・Müller（2020）\cite{jacobsmuller2020} は、39 か国の株式市場で 241 のアノマリーを調べ、公表の後の明確な減衰は米国だけで見られ、ほかの国では公表の後もリターンが残ることを示した。裁定の限界（空売りの難しさ、流動性、コスト）が、国による差を生んでいる可能性がある。Chordia・Subrahmanyam・Tong（2014）\cite{cst2014} は、米国で取引のコストが下がり売買が増えた 2000 年代以降（ティックの小数化の後）、アノマリーの多くが小さくなったことを示した。

\subsection{再現できるか}

ファクターの研究には、「そもそも本当にあったのか」という問いもある。

\begin{itemize}
\item \textbf{Hou・Xue・Zhang（2020）}\cite{hxz2020}：452 のアノマリーを、小型株の影響を抑えた同じ方法（NYSE の基準で分け、時価総額で加重）で作り直すと、6 割を超えるものが $t = 1.96$ の基準を超えなかった。多くのアノマリーは、極めて小さな株に依存していた。
\item \textbf{Jensen・Kelly・Pedersen（2023）}\cite{jkp2023}：153 のファクターを 93 か国で作り直し、ベイズの方法でファクターの間の情報を共有して評価すると、大多数のファクターが再現でき、米国の外でも成り立つ。「再現の危機」はないと結論した。
\item \textbf{Chen・Zimmermann（2022）}\cite{chenzimmermann2022}：300 を超える予測の変数のコードとデータを公開し、元の論文の結果の大部分が再現できることを示した。
\end{itemize}

二つの見方は、矛盾するというより、「何を再現とみなすか」の違いである。極めて小さな株を除き、時価総額で加重すると多くのアノマリーは弱くなるが、適切に作れば多くのファクターの方向と大きさは再現できる。

\subsection{取引コストの後}

Novy-Marx・Velikov（2016）\cite{novymarxvelikov2016} は、アノマリーの戦略の取引コストを推定し、回転率の低い戦略（割安、収益性など）はコストの後も残るが、回転率が月 50\% を超えるような戦略（短期の反転など）は、ほとんどがコストの後で有意でなくなることを示した。売買の頻度を下げる工夫（\chapref{ch:cost} の「売買しない領域」）で、コストを大きく減らせる。

Israel・Moskowitz（2013）\cite{israelmoskowitz2013} は、割安とモメンタムの収益が、空売りの側に依存しているわけではなく、大型株でも残り、時代とともに小さくなってもいない（規模を除く）ことを示した。運用会社の研究者による論文で、ファクター投資が現実に運用できることを論じたものである。

\section{確立していること、していないこと}

\begin{keybox}[まとめ]
\begin{description}[leftmargin=1em,style=nextline]
\item[確立している]
  割安・モメンタム・収益性・低リスクのファクターは、長い期間・多くの国・多くの資産クラスで正のリターンを持つ。アノマリーは、公表の後に平均して半分前後に小さくなる（米国）。極めて小さな株に依存するアノマリーは、実際には取れないことが多い。モメンタムはまれに大きく崩れ、変動率に応じた量の調整で和らげられる。
\item[議論が続いている]
  収益の原因がリスクか行動の偏りか。ファクターの収益がこれからも続くか（特に割安）。何百ものファクターのうち、独立な情報を持つものはいくつか（第 \ref{ch:ml} 章）。
\end{description}
\end{keybox}

\chapter{統計的裁定と平均回帰}
\label{ch:statarb}

\section{戦略の中身}

統計的裁定は、似た値動きをする銘柄の間の一時的な「ずれ」が戻ることに賭ける。代表的な形は三つある。

\begin{enumerate}
\item \textbf{ペアトレード}：過去の価格が一緒に動いてきた 2 銘柄の価格差（あるいは比）が、普段より大きく開いたら、上がった方を売り、下がった方を買う。差が戻ったら手じまう。
\item \textbf{ファクターの残差の平均回帰}：各銘柄のリターンを、市場・業種のファクター（主成分や ETF）に回帰し、その残差の累積が平均から離れたら、戻る方向に賭ける。数百〜数千銘柄に同時に小さく賭ける。
\item \textbf{短期の反転}：直近の 1 週間〜1 か月に大きく下げた株を買い、大きく上げた株を売る。
\end{enumerate}

保有期間は数分〜数週間で、回転率が高い。収益の源は、一時的な需給の偏り（大口の売買、指数の資金の流れ、損切り）に対して流動性を供給する対価と解釈されることが多い。

\section{短期の反転}

Lehmann（1990）\cite{lehmann1990} は週次、Jegadeesh（1990）\cite{jegadeesh1990} は月次で、米国株のリターンに強い負の自己相関（反転）があることを示した。直近の負け組を買い勝ち組を売る逆張りの戦略は、取引コストの前では大きな収益を上げた。

\begin{paperbox}[Nagel（2012）「消える流動性」\cite{nagel2012}]
短期の反転の戦略の収益は、流動性を供給する対価と解釈できる。この収益は、市場の変動率（VIX）が高い時期に大きくなる。流動性の供給者（マーケットメイカーやヘッジファンド）自身が資金の制約を受け、流動性の供給が細る時期に、その対価が上がるからである。2008 年の金融危機では、反転の戦略の期待収益が大きく上がった。
\end{paperbox}

ただし、反転の戦略は回転率が非常に高く、取引コストの後にはほとんどの収益が消える（Novy-Marx・Velikov 2016 \cite{novymarxvelikov2016}）。実際にこの収益を取れるのは、コストの低い高頻度の業者と、執行の上手な大手の運用会社に限られる。

\section{ペアトレード}

\begin{paperbox}[Gatev・Goetzmann・Rouwenhorst（2006）\cite{ggr2006}]
米国株（1962〜2002 年）で、12 か月の価格の動きが最も近い（正規化した価格の距離が最小の）組をペアとし、その後 6 か月の間、価格差が過去の標準偏差の 2 倍を超えて開いたら売買する単純な規則を検証した。自己資金の要らないペアのポートフォリオは、年率で最大 11\% 程度の超過リターンを上げ、保守的な取引コストを引いても残った。ただし、収益は期間の後半に小さくなっていた。
\end{paperbox}

Do・Faff（2010）\cite{dofaff2010} は、同じ方法を 2009 年まで延ばし、ペアトレードの収益が 1990 年代以降に低下していること、ただし市場の混乱の時期（2000〜2002 年、2007〜2009 年）には収益が上がることを示した。

ペアを選ぶ方法として、価格の距離のほかに、\textbf{共和分}（Engle・Granger 1987 \cite{englegranger1987}）がよく使われる。二つの価格の系列がそれぞれランダムウォークでも、ある線形結合 $p^A_t - \beta p^B_t$ が定常（平均に戻る）なら、その結合の値を売買の信号にできる。

\section{ファクターの残差の平均回帰}

\begin{paperbox}[Avellaneda・Lee（2010）\cite{avellanedalee2010}]
米国株の各銘柄のリターンを、主成分（相関行列の上位の固有ベクトル）または業種の ETF のリターンに回帰し、残差の累積を平均回帰の過程（Ornstein--Uhlenbeck 過程）としてモデル化した。残差が平均から標準偏差の何倍か離れたら売買する戦略は、1997〜2007 年の検証で高いシャープレシオを示したが、2003 年以降は成績が下がり、2007 年 8 月に大きく損をした。売買高の情報を加えると成績が改善した。
\end{paperbox}

この形は、クオンツのマーケット・ニュートラルの株式の運用の中心に近い。銘柄の数が多いので、基本法則（\ref{sec:flam} 節）の幅が大きく、1 回の賭けの IC が小さくても、全体の情報レシオは高くなりうる。

\section{夜と昼}

Lou・Polk・Skouras（2019）\cite{lps2019} は、米国株のリターンを、夜間（前日の終値から当日の始値）と日中（始値から終値）に分けると、多くのアノマリーが一方に偏ることを示した。モメンタムの収益はほぼすべて夜間に生じ、ほかの多くのアノマリーの収益は日中に生じ、夜間にはむしろ逆に動く。夜間と日中では主な売買の主体（個人と機関投資家）が違い、その「綱引き」が収益の時間の分布を決めている、と解釈している。

\section{2007 年 8 月の quant quake}

\begin{paperbox}[Khandani・Lo（2011）\cite{khandanilo2011}]
2007 年 8 月 6〜9 日、株式のマーケット・ニュートラルのクオンツのファンドの多くが、市場全体が落ち着いている中で、数日のうちに大きな損失を出し、10 日には大きく戻した。著者らは、単純な短期の反転の戦略（Lehmann 型の逆張り）と、ファクターのポートフォリオのリターンで、この出来事を再現した。
\begin{itemize}
\item ある大きなファンドが、別の資産（信用関連など）の損失や資金の必要から、株式のマーケット・ニュートラルのポジションを急いで解消した可能性が高い。
\item 多くのクオンツのファンドが、似たファクター・似たシグナルで、似たポジションを持っていた。一つの解消が、ほかのファンドの損失を生み、それがさらにリスクの上限による解消を呼んだ（解消の連鎖）。
\item 戦略の収益は 1998 年ごろから低下していて、同じ収益を得るために使うレバレッジが大きくなっていた。
\end{itemize}
\end{paperbox}

この出来事の教訓は、「市場全体のリスクがゼロでも、同じ取引をしている人のリスクは消えない」ことである。リスクモデル（\chapref{ch:risk}）は、価格の過去の共分散からリスクを測るが、ほかの運用者のポジションの重なりは、そこに現れない。Goldman Sachs の幹部が当時「25 標準偏差の動きが何日も続いた」と述べたことは、正規分布を前提にしたリスクの物差しが、混雑の解消の前では役に立たないことを示す逸話としてよく引かれる（Deep Research の答え 1 は、この発言を Khandani・Lo の結果として書いていたが、論文の結論ではない）。

\section{個人や研究で再現できるか}

\begin{itemize}
\item 日次の終値によるペアトレードや反転の戦略は、無料の株価のデータで検証できる。ただし、取引コスト（スプレッドとインパクト）と空売りの費用を正しく入れると、ほとんどの収益が消える。
\item 実際の統計的裁定の優位は、日中の細かいデータ、低い取引コスト、速い執行、多数の銘柄の同時の管理にある。個人が実運用で競うのは難しい。
\item 研究として意味があるのは、「どういう時期に収益が上がるか」（変動率が高く流動性が細る時期、Nagel 2012）や、「誰の売買が歪みを生んでいるか」を調べることである。
\end{itemize}

\chapter{トレンドフォロー}
\label{ch:trend}

\section{戦略の中身}

トレンドフォロー（時系列モメンタム、マネージド・フューチャーズ、CTA の戦略）は、株価指数・国債・通貨・商品の先物を、それ自身の過去のリターンが正なら買い、負なら売る。

\begin{equation}
w_{i,t} = \mathrm{sign}\Big(\sum_{k=1}^{12} r_{i,t-k}\Big) \times \frac{\sigma^*}{\hat\sigma_{i,t}}
\end{equation}

\begin{itemize}
\item \textbf{信号}：過去 1〜12 か月のリターンの符号、移動平均の交差（短期の平均が長期の平均を上回る）など。期間の違う信号を組み合わせることが多い。
\item \textbf{量}：各市場の変動率 $\hat\sigma_{i,t}$ の逆数に比例させ、どの市場も同じリスクにする（ボラティリティ・スケーリング）。
\item \textbf{分散}：数十〜百以上の先物市場に分散する。
\end{itemize}

横断面のモメンタム（第 \ref{ch:equityfactor} 章、銘柄を比べて勝ち組を買う）と違い、時系列のモメンタムは、すべての市場が上がっていればすべて買う。

\section{中心の論文}

\begin{paperbox}[Moskowitz・Ooi・Pedersen（2012）「時系列のモメンタム」\cite{mop2012}]
58 の流動性の高い先物・フォワード（株価指数、通貨、商品、国債）で、各資産の過去 12 か月のリターンが、その後 1 か月のリターンを正に予測することを示した。予測の力は 1〜12 か月で続き、より長い期間では一部が戻る。全資産に分散した時系列モメンタムの戦略は、市場・割安・横断面のモメンタムなどの標準的なファクターでは説明できない大きな超過リターンを上げ、市場が極端に上げ下げした時期に最も良い成績を上げた。投機筋（非商業の先物の参加者）がトレンドに乗り、ヘッジをする当業者がその反対側に立っていることも示した。
\end{paperbox}

\begin{paperbox}[Hurst・Ooi・Pedersen（2017）「トレンドフォローの 1 世紀の証拠」\cite{hop2017}]
1880 年まで遡って、株価指数・国債・通貨・商品の時系列モメンタムの戦略を作り直し、ほぼすべての 10 年間で正の平均リターンを上げたこと、景気後退や戦争・金融危機を含む多くの環境で成績を保ったこと、株式や債券の大きな下げの時期に利益を上げることが多かったことを示した。
\end{paperbox}

\section{なぜ儲かるのか}

\begin{itemize}
\item \textbf{反応の不足と、その後の行き過ぎ}：新しい情報への反応が遅れ（最初の反応の不足）、その後、群集の行動で行き過ぎる。どちらの局面でも、トレンドに乗れば利益が出る。
\item \textbf{当業者のヘッジ}：生産者・消費者が価格のリスクを避けるために先物を売買し、その反対に立つ投機筋に対価が支払われる。
\item \textbf{中央銀行の介入}：為替や金利の政策は、ゆっくり一方向に動くことが多い。
\end{itemize}

\section{オプションのような形}

Fung・Hsieh（2001）\cite{funghsieh2001} は、トレンドフォローの運用者のリターンが、「期間の中の最高値で売り、最安値で買う」権利を持つオプション（ルックバック・ストラドル）を買う戦略のリターンに似ていることを示した。トレンドフォローは、大きく一方向に動いた市場で利益を上げ、方向のない市場では小さな損を重ねる。オプションを買い続けているのと同じ形で、リターンの分布は正の歪度を持つ。この「ルックバック・ストラドルのファクター」は、ヘッジファンドの成績の評価に使われる Fung・Hsieh の 7 ファクター・モデル（第 \ref{ch:hfperf} 章）の一部になった。

\section{批判}

\begin{itemize}
\item \textbf{ボラティリティ・スケーリングの寄与}：Kim・Tse・Wald（2016）\cite{kimtsewald2016} は、時系列モメンタムの成績の大部分は、各市場の量を変動率で調整すること（ボラティリティ・スケーリング）から来ていて、量を調整しなければ、単に買い持ちする戦略と大差ないと論じた。
\item \textbf{予測の力は本当にあるか}：Huang・Li・Wang・Zhou（2020）\cite{hlwz2020} は、資産ごとに見ると、過去 12 か月のリターンが将来のリターンを予測する証拠は弱く、全資産をまとめた回帰の有意性は一部の資産と推定の偏りに依存すると論じた。時系列モメンタムの戦略の収益は、各資産の長期の平均リターン（予測の力がなくても得られる部分）でかなり説明できる、という。
\end{itemize}

\section{弱点}

\begin{itemize}
\item \textbf{方向のない相場}：上げ下げを繰り返す相場では、信号が入れ替わるたびに損をする（往復の損）。
\item \textbf{急な反転}：トレンドが急に反転すると、反応が遅れて損をする。2009 年の春、2020 年 3 月の急落と急反発など。
\item \textbf{長い低迷}：2010 年代は、金利が低く変動の小さい環境で、多くの CTA の成績が低迷した。物価と金利が大きく動いた 2022 年には、多くが大きな利益を上げた（業界の指数で広く報じられた事実で、特定の論文の値ではない）。
\end{itemize}

\section{個人や研究で再現できるか}

\begin{itemize}
\item 先物のつなぎ足の価格のデータがあれば、検証は難しくない。規則が単純で、過学習の余地も比較的小さい。
\item 実運用では、数十の先物市場に分散するための証拠金、限月の乗り換え（ロール）のコスト、最小の取引単位が壁になる。個人なら、先物の代わりに ETF で近似することになる。
\item 結果は、信号の期間・ボラティリティの推定の窓・市場の選び方に敏感である。多くの組み合わせを試すと、過学習（\chapref{ch:backtest}）の危険がある。
\end{itemize}

\begin{projbox}
トレンドフォローは、LLM-BLM の研究とは別の種類の情報（テキストではなく価格）を使うが、「予測の力そのものより、量の調整（ボラティリティ・スケーリング）とリスクの配分で成績の多くが決まる」という論争（Kim ほか 2016、Huang ほか 2020）は、LLM の点数を使う運用にもそのまま当てはまる。点数の予測の力を評価するときに、ポートフォリオの作り方の効果と分けて測る必要がある。
\end{projbox}

\chapter{グローバル・マクロとキャリー}
\label{ch:macro}

\section{裁量のマクロとシステマティックなマクロ}

グローバル・マクロは、経済の見通しから、金利・通貨・株価指数・商品の方向に賭ける。有名な運用者の多くは裁量であり、その判断の過程を論文で検証することは難しい。研究で扱えるのは、規則で書ける部分、つまり\textbf{キャリー}、\textbf{割安}（購買力平価などからの乖離）、\textbf{トレンド}（第 \ref{ch:trend} 章）である。この章は、主にキャリーを扱う。

\section{通貨のキャリー}

\subsection{金利平価の破れ}

金利の高い通貨の預金は、金利の低い通貨の預金より多くの利息を生む。カバーなしの金利平価（UIP）が成り立つなら、その差は、高金利の通貨の将来の値下がりでちょうど打ち消され、平均的な利益はゼロになるはずである。

\begin{paperbox}[Fama（1984）\cite{fama1984}]
先渡しの為替レートと、その後の直物の為替レートの関係を検証し、UIP が成り立たないことを示した。金利の高い通貨は、UIP の予測ほど値下がりせず、平均ではむしろ値上がりする傾向がある（フォワード・プレミアムのパズル）。
\end{paperbox}

そこで、金利の低い通貨（円、スイス・フランなど）を借りて売り、金利の高い通貨を買う\textbf{キャリー・トレード}は、平均すると利益を上げてきた。

\subsection{なぜ儲かるのか：リスクの対価}

\begin{paperbox}[Lustig・Roussanov・Verdelhan（2011）\cite{lrv2011}]
通貨を金利の水準で並べたポートフォリオのリターンが、二つの共通のファクターでよく説明できることを示した。一つはドルの全体の水準、もう一つは高金利の通貨を買い低金利の通貨を売るキャリーのファクター（HML\textsubscript{FX}）である。高金利の通貨は、世界の株式市場の変動率が上がる局面（世界的な危機）で損をする。キャリーの収益は、この世界共通のリスクを引き受ける対価だと解釈した。
\end{paperbox}

Menkhoff・Sarno・Schmeling・Schrimpf（2012）\cite{msss2012} は、世界の為替の変動率の予想外の上昇が、キャリーのリターンの横断面を説明することを示した。高金利の通貨は、為替の変動率が急に上がるときに損をする。

\subsection{急落}

\begin{paperbox}[Brunnermeier・Nagel・Pedersen（2008）「キャリー・トレードと通貨の急落」\cite{bnp2008}]
キャリー・トレードのリターンは負の歪度を持つ（ふだんは少しずつ儲け、まれに急に大きく損をする）。金利差が大きいほど、その後の急落のリスクも大きい。市場の変動率（VIX）が上がると、投機筋がキャリーのポジションを解消し、高金利の通貨が急落する。資金の流動性の制約が、急落の引き金になる（\chapref{ch:risk} の流動性のらせん）。
\end{paperbox}

2008 年の秋には、円の急騰で多くのキャリー・トレードが数週間のうちに大きく損をした。「蒸気ローラーの前で小銭を拾う」とたとえられる構造である。

\section{あらゆる資産のキャリー}

\begin{paperbox}[Koijen・Moskowitz・Pedersen・Vrugt（2018）「キャリー」\cite{kmpv2018}]
キャリーを「価格が変わらなかったときに得られる収益」と一般化し、株価指数（配当利回りと先物の価格の関係）、国債（利回り曲線の傾き）、商品（期近と期先の価格差）、通貨（金利差）、オプションなどで定義した。どの資産クラスでも、キャリーの高い資産が低い資産より高いリターンを持つ。資産クラスをまたいだキャリーの戦略の間の相関は低く、分散した「世界のキャリーのファクター」は高いリスク調整後のリターンを持つ。一方で、世界の景気後退や流動性の低下の局面では損をする。
\end{paperbox}

キャリーは、市場の「予想」（先物の価格）と、「何も変わらない」という予測のずれに賭ける戦略と見ることもできる。

\section{危機の局面}

Ang・Bekaert（2002）\cite{angbekaert2002} は、国際的な株式の資産配分を、「平時」と「相関と変動率が高い局面」を確率的に行き来するモデル（レジーム・スイッチング、\chapref{ch:model}）で分析した。下げ相場の局面では国の間の相関が上がり、国際分散の効果が小さくなるが、それでも分散の価値は残ることを示した。マクロの運用者が、局面の変化を見越して資産配分を動かす考え方の、理論的な裏付けの一つである。

\section{裁量のマクロの成績}

裁量のマクロのファンドの成績は、ヘッジファンドのデータベース（第 \ref{ch:hfperf} 章）で調べるしかない。Deep Research の答え 1 は、マクロのファンドのシャープレシオが上げ相場・下げ相場・危機でそれぞれ 1.0・0.8・0.9 だったという数値を挙げたが、出典を確かめられなかったので使わない。データベースの偏り（第 \ref{ch:hfperf} 章）を考えると、裁量のマクロの「腕」の大きさについて、確かなことは言いにくい。

\section{個人や研究で再現できるか}

\begin{itemize}
\item 各国の短期金利・為替レート・先物の価格は、公開のデータで入手しやすい。通貨のキャリーや、資産クラスをまたぐキャリーの戦略の検証は、比較的容易である。
\item ただし、キャリーの収益は、まれな急落の損失で大きく削られる。検証の期間に大きな急落が含まれていないと、成績が過大に見える。
\item 実運用では、個人の為替の取引のスプレッドや、借入れの費用が、機関投資家より不利である。
\end{itemize}

\chapter{イベント・ドリブン}
\label{ch:event}

企業の出来事（決算、アナリストの推奨の変更、合併、指数の入れ替え、破綻、アクティビストの介入）の前後で生じる価格の歪みをとる戦略である。LLM-BLM の研究に最も近いのが、決算発表の後のドリフトなので、そこから始める。

\section{決算発表の後のドリフト}

\subsection{発見}

Ball・Brown（1968）\cite{ballbrown1968} は、会計の利益の情報が株価に反映されることを示した最初の研究の一つで、すでに、利益が予想を上回った会社の株価が、発表の後もしばらく上がり続ける傾向を観察していた。

\begin{paperbox}[Bernard・Thomas（1989, 1990）\cite{bernardthomas1989,bernardthomas1990}]
四半期の利益の驚き（SUE：前年同期からの変化を、その変動の大きさで割ったもの）で銘柄を 10 の組に分けると、驚きの大きい組は発表の後 60 取引日にわたって上がり続け、小さい組は下がり続ける（決算発表の後のドリフト、PEAD）。この遅れた反応は、リスクの対価では説明しにくい。1990 年の論文は、四半期の利益の驚きには自己相関がある（今期が良ければ、来期も前年同期より良いことが多い）のに、株価はそれを十分に織り込まず、次の四半期の発表のときに、予測できる方向の反応が起きることを示した。投資家が、利益の時系列の性質を単純化して捉えている、という解釈である。
\end{paperbox}

\subsection{注目と反応の遅れ}

PEAD は、投資家の注目が足りないときに大きくなる。

\begin{itemize}
\item \textbf{金曜日の発表}：DellaVigna・Pollet（2009）\cite{dellavigna2009} は、金曜日に発表された決算では、発表の直後の反応が小さく、その後の遅れた反応が大きいことを示した。週末の前で、投資家の注目が下がるためと解釈される。
\item \textbf{同じ日の発表の多さ}：Hirshleifer・Lim・Teoh（2009）\cite{hlt2009} は、同じ日にほかの会社の決算発表が多いほど、直後の株価と売買高の反応が小さく、ドリフトが大きいことを示した（注意の分散）。競合するニュースの少ない日の発表では、PEAD を狙う戦略は利益にならなかった。
\item \textbf{情報の処理の遅れ}：Akbas・Markov・Subasi・Weisbrod \cite{akbas2017} は、アナリストの予想のデータベースへの反映が遅れる銘柄ほど、発表のときの反応が小さく、その後のドリフトが大きいことを示した。
\end{itemize}

\subsection{流動性と、消えたドリフト}

\begin{paperbox}[Chordia・Goyal・Sadka・Sadka・Shivakumar（2009）\cite{cgsss2009}]
PEAD は、主に流動性の極めて低い銘柄で起きていることを示した。驚きの大きい株を買い小さい株を売る戦略の時価総額加重のリターンは、最も流動性の高い銘柄で月 0.14\%、最も低い銘柄で月 1.60\% だった。流動性の低い銘柄は売買のコストとインパクトが大きく、さまざまなコストの推定で、紙の上の収益の 63〜100\% がコストで消える。コストそのものがドリフトの残る原因かもしれない、と論じた。
\end{paperbox}

\begin{paperbox}[Martineau（2022）「安らかに眠れ、決算発表の後のドリフト」\cite{martineau2022}]
現代の米国株では、決算の驚きは発表の当日に株価に完全に反映され、PEAD は消えた、と報告した。大型株では 2006 年以降 PEAD は見られず、極めて小さな株（マイクロキャップ）でも最近になって消えた。著者は、PEAD を研究する人に、それが今もあるかを確かめる方法の提案もしている。
\end{paperbox}

\begin{projbox}
LLM-BLM の \SP{} の実験で、LLM の読み取りの情報が発表の当日に株価に入りきり、翌日以降のドリフトを予測しなかったことは、Martineau（2022）の結果とよく一致する。小型株の実験 S1 は、Chordia ほか（2009）がドリフトの集中を示した流動性の低い・注目の低い銘柄に的を絞ったものである。ただし、Martineau（2022）によれば、数値の驚き（SUE）による PEAD は、極めて小さな株でも最近は消えている。S1 が問うのは、数値ではなく\textbf{文章の読み取り}が、小型株でドリフトを予測するかであり（PEAD.txt に近い問い）、先行研究から見て、予測できなくても不思議ではない。また、予測できたとしても、取引コストが高く、紙の上のドリフトが実際に取れるとは限らない（Chordia ほか 2009）。
\end{projbox}

\subsection{テキストによるドリフト}

\begin{paperbox}[Meursault・Liang・Routledge・Scanlon（2023）「PEAD.txt」\cite{peadtxt}]
決算説明会の発言録の文章から、機械学習で「文章による驚き」（SUE.txt）を作り、それで銘柄を並べると、数値の利益の驚き（SUE）によるものより大きく、長く続くドリフトが生じることを示した。数値の驚きのドリフトが近年弱まったのに対し、文章の驚きのドリフトは残っている、という。
\end{paperbox}

文章の情報は、数値の情報より処理が難しく、反映に時間がかかる、という仮説である。LLM による読み取りは、この文章の情報を安く速く処理する道具であり、その普及が、この種のドリフトを小さくしていく可能性がある（第 \ref{ch:text} 章）。

\section{アナリストの推奨の変更}

\begin{paperbox}[Womack（1996）\cite{womack1996}]
大手の証券会社のアナリストが推奨を変更すると、発表の数日に大きな反応がある。その後も同じ方向のドリフトが続き、買いへの変更の後のドリフトは小さく短い（平均 +2.4\%）が、売りへの変更の後のドリフトは大きく、6 か月にわたって続く（平均 $-9.1\%$）。
\end{paperbox}

売りの推奨の情報の反映が遅いのは、空売りの制約と、アナリストが売りの推奨を出したがらないこと（売りの推奨は情報の価値が高い）による、と解釈される。Hong・Lim・Stein（2000）\cite{hls2000} の「悪い知らせはゆっくり伝わる」とも合う。

\section{合併裁定}

\begin{paperbox}[Mitchell・Pulvino（2001）\cite{mitchellpulvino2001}]
1963〜1998 年の約 4,750 件の合併・買収で、買収の発表の後に買収される会社の株を買い（株式の交換による買収なら、買収する会社の株を売り）、成立まで持つ合併裁定の収益を調べた。合併裁定のリターンは、市場が平穏か上げているときは市場とほぼ無関係だが、市場が大きく下げるときには市場と強く連動する。この形は、指数のプット・オプションを売った戦略に似ている。この非線形なリスクを考慮しても、年 4\% 程度の超過リターンが残った。
\end{paperbox}

Baker・Savaşoglu（2002）\cite{bakersavasoglu2002} は、合併裁定の収益が、成立しないリスク（成立の失敗の確率）と、裁定に使える資金の限界（大きな案件ほど裁定者が持ちきれない）の対価であることを示した（裁定の限界）。

合併裁定の損失は、案件が不成立になったとき（独占禁止の当局の反対、資金調達の失敗、相手先の業績の悪化）に集中する。また、2008 年のような危機では、多くの案件が同時に不確かになり、裁定者が資金の制約で一斉に手じまうので、スプレッドが急に広がる（Mitchell・Pulvino 2012 \cite{mitchellpulvino2012}、第 \ref{ch:rv} 章）。

\section{指数の入れ替え}

Harris・Gurel（1986）\cite{harrisgurel1986} と Shleifer（1986）\cite{shleifer1986} は、\SP{} への採用の発表で株価が数\% 上がることを示した。Harris・Gurel は、その上昇の多くが数週間で戻ること（一時的な売買の圧力）を、Shleifer は上昇が残ること（株の需要の曲線が右下がりであること）を強調した。指数のファンドが採用の日に機械的に買うので、その前に買って指数のファンドに売る戦略が成り立つ。

Greenwood・Sammon \cite{greenwoodsammon} は、この「指数の効果」が大きく縮んだことを示した。\SP{} への採用の平均の超過リターンは、1980 年代に 3.4\%、1990 年代に 7.6\% だったが、最近の 10 年は 0.8\% にすぎない。除外も同じで、1980〜90 年代の大きな負の超過リターンが、2010〜2020 年には $-0.6\%$ になった。指数に連動する資産の割合は大きく増えたのに、効果は小さくなった（Deep Research の答え 1 の値は、NBER の要旨と一致した）。

\begin{figure}[H]
\centering
\begin{tikzpicture}[x=1mm,y=4mm]
\foreach \y in {0,2,4,6,8} {\draw[ink2!25] (0,\y)--(80,\y); \node[lab,left] at (0,\y) {\y\%};}
\fill[c1!60] (8,0) rectangle (24,3.4);  \node[lab,above] at (16,3.4) {3.4\%};
\fill[c1] (32,0) rectangle (48,7.6);    \node[lab,above] at (40,7.6) {7.6\%};
\fill[c2!75] (56,0) rectangle (72,0.8); \node[lab,above] at (64,0.8) {0.8\%};
\node[lab,below] at (16,-0.1) {1980 年代};
\node[lab,below] at (40,-0.1) {1990 年代};
\node[lab,below] at (64,-0.1) {最近の 10 年};
\end{tikzpicture}
\caption{\SP{} への採用の平均の超過リターン（Greenwood・Sammon の NBER の要旨の値から描いた）。}
\label{fig:indexeffect}
\end{figure}

\section{ディストレストとアクティビスト}

\begin{itemize}
\item \textbf{ディストレスト}：Jiang・Li・Wang（2012）\cite{jlw2012} は、米国の連邦破産法 11 条の手続きにヘッジファンドが債権者として加わると、再建が成立する確率が上がり、劣後の債権者の回収が増え、経営者の交代が増えることを示した。ヘッジファンドは、破綻した会社の債権を安く買い、手続きに影響力を持つことで収益を上げる。
\item \textbf{アクティビスト}：Brav・Jiang・Partnoy・Thomas（2008）\cite{bjpt2008} は、ヘッジファンドが 5\% 以上の株を持ったと届け出る（13D の提出）と、株価がその前後で大きな超過リターン（約 7\%）を示し、その後に戻らないこと、対象の会社の利益率や配当が改善することを示した。
\end{itemize}

どちらも、大きな資金と、法律・企業の調査の専門性が必要で、規則による再現は難しい。

\section{個人や研究で再現できるか}

\begin{itemize}
\item PEAD や推奨の変更の検証には、決算の日付と予想（アナリストの予想のデータベース）、あるいは決算の数値の時系列が要る。前者は有料、後者は SEC の XBRL などで無料に近い形で作れる。
\item 近年の大型株では PEAD がほぼ消えているので、検証するなら、注目・流動性の低い銘柄や、文章の情報（PEAD.txt）に的を絞ることになる。その場合、取引コストを厳しく見積もる必要がある。
\item 合併裁定は、発表された案件の一覧と価格があれば検証できる。不成立の案件の扱い（生存者バイアス）に注意する。
\end{itemize}

\chapter{レラティブ・バリューとボラティリティ}
\label{ch:rv}

\section{債券の裁定}

債券の裁定は、ほぼ同じ将来のお金の流れを持つ二つの債券（あるいは債券とデリバティブ）の価格の小さな差をとる。差は小さいので、大きなレバレッジをかける。代表的な取引は次のとおり。

\begin{itemize}
\item \textbf{発行されたばかりの銘柄と、少し前の銘柄}：最新の国債は売買が多く、少し割高になる。少し前の銘柄を買い、最新の銘柄を売り、差が縮むのを待つ。
\item \textbf{スワップ・スプレッド}：国債の利回りと、同じ年限の金利スワップの固定金利の差に賭ける。
\item \textbf{利回り曲線}：曲線の形の歪み（ある年限だけが周りより高い）に賭ける。
\item \textbf{現物と先物のベーシス}：国債の現物を買い（レポで資金を借りる）、先物を売る。
\end{itemize}

\begin{paperbox}[Duarte・Longstaff・Yu（2007）「蒸気ローラーの前の小銭か」\cite{dly2007}]
スワップ・スプレッド、利回り曲線、住宅ローン担保証券、金利のボラティリティ、資本構成（同じ会社の株と社債）の裁定の 5 つの債券の裁定の戦略を、規則で作り直して検証した。多くの戦略は有意な超過リターンを上げ、その一部は、「小さな利益を重ねて、まれに大損する」という通説とは違い、正の歪度を持っていた。分析に高度なモデルを要する（「知的な資本」の要る）戦略ほど、超過リターンが大きかった。
\end{paperbox}

\subsection{LTCM}

この種の取引で最も有名な失敗が、1998 年の LTCM である。LTCM は、債券の裁定を中心に、資本の 25 倍前後の資産を持ち（デリバティブを含めるとさらに大きい）、ロシアの債務不履行の後の「質への逃避」で、ほぼすべての取引の差が同時に広がって、数か月で資本の大半を失った。詳しくは第 \ref{ch:failures} 章で扱う。

\subsection{2020 年 3 月の国債の売り}

\begin{paperbox}[Kruttli・Monin・Petrasek・Watugala（2021）\cite{kruttli2021}]
米連邦準備制度の報告（Form PF のデータ）を使い、ヘッジファンドの米国債の取引を分析した。2020 年 3 月の新型コロナの急落で、国債の市場の変動が急に高まると、レポの資金の条件が厳しくなり、国債の取引をしていたファンドは損失を出し、国債のポジションを大きく減らした（論文は、平均的なファンドの国債のポートフォリオの損失を約 7\%、ポジションの縮小を約 2 割と報告している）。売りは、資金の調達の制約の厳しいファンドに集中した。中央銀行が国債の大量の買い入れで介入し、市場は落ち着いた。
\end{paperbox}

現物と先物のベーシスの取引は、その後も拡大し、金融の安定の観点から当局が注目している。

\section{転換社債の裁定}

転換社債（一定の条件で株に転換できる社債）を買い、その株を空売りして、株価の変動のリスクを打ち消す（デルタ・ヘッジ）。転換社債は、発行の時に割安になりやすいことが知られ、その割安さと、株価の変動から生じる利益（ガンマの利益）をとる。

Agarwal・Fung・Loon・Naik（2011）\cite{aflm2011} は、転換社債の裁定のヘッジファンドのリターンが、転換社債の市場の要因（新しい発行の量、割安さ）で説明できることを示した。新しい発行が多く割安な時期には収益が上がり、裁定の資金が集まって割安さが小さくなると、収益も小さくなる。

\subsection{遅い資本}

\begin{paperbox}[Mitchell・Pedersen・Pulvino（2007）「遅い資本」\cite{mpp2007}]
2005 年に転換社債の裁定のファンドから資金が大量に引き上げられたとき、転換社債の価格は、理論的な価値から大きく下がり、その回復には何か月もかかった。歪みを直すはずの裁定の資本は、すぐには入ってこない（遅い資本）。専門の裁定者が損をして資金を失うと、ほかの投資家が代わりに入るまでに時間がかかるからである。
\end{paperbox}

Mitchell・Pulvino（2012）\cite{mitchellpulvino2012} は、2008 年の金融危機で、プライムブローカーが担保の再利用（再担保）による資金の提供を急に止めたことが、転換社債の裁定や合併裁定のスプレッドの急拡大（裁定の崩れ）を引き起こしたことを示した。これは、裁定の収益が「資金の提供者が資金を出し続けること」に依存していることを示す。

\section{ボラティリティのリスク・プレミアム}

\subsection{オプションは平均して割高}

\begin{paperbox}[Coval・Shumway（2001）\cite{covalshumway2001}]
\SP{} の指数のオプションのリターンを調べ、プット・オプションのリターンが CAPM の予測よりずっと低いことを示した。市場のリスクを打ち消した（ベータがゼロの）アット・ザ・マネーのストラドル（コールとプットの買い）は、週に約 3\% の平均の損失を出した。変動率そのもの（あるいは急落）のリスクに、対価が支払われていることを示す。
\end{paperbox}

\begin{paperbox}[Bakshi・Kapadia（2003）\cite{bakshikapadia2003}]
オプションを買って株価の変動のリスクを打ち消したポートフォリオ（デルタ・ヘッジしたオプション）の損益が、平均して負であることを示した。市場の変動率のリスクに負の価格（変動率が上がるリスクを避けるために、投資家が対価を払う）がついていることの証拠である。
\end{paperbox}

\begin{paperbox}[Carr・Wu（2009）「分散のリスク・プレミアム」\cite{carrwu2009}]
オプションの価格から、分散スワップ（将来の実現の分散と、事前に決めた固定の分散を交換する取引）の公正な固定の値を計算し、実際に実現した分散と比べた。\SP{} などの株価指数では、実現の分散は平均してオプションが示す分散よりずっと小さく、分散スワップの買いは平均して大きな損失を出した（分散のリスク・プレミアムが大きく負）。このプレミアムは、CAPM のベータや Fama・French のファクターでは説明できない。
\end{paperbox}

\subsection{相関のリスク}

Driessen・Maenhout・Vilkov（2009）\cite{dmv2009} は、株価指数のオプションが、個別の株のオプションに比べて割高であることを示し、それが「銘柄の間の相関が上がるリスク」への対価であると解釈した。指数のオプションを売り、個別の株のオプションを買う取引（分散の取引、ディスパージョン）は、この相関のリスク・プレミアムをとる戦略である。

\subsection{ボラティリティの売りの崩れ：2018 年 2 月}

\begin{paperbox}[Augustin・Cheng・Van den Bergen（2021）「ボルマゲドン」\cite{volmageddon2021}]
2018 年 2 月 5 日、VIX（\SP{} のオプションが示す予想の変動率の指数）が 1 日で 2 倍以上（17 台から 37 台）に上がり、VIX の先物を売る形の上場商品（XIV など）は、価値の 9 割以上を失って償還された。これらの商品は、日々の値動きに応じて、引けにかけて VIX の先物を機械的に買い戻す仕組みだった。VIX が上がると、その買い戻しの量が膨らみ、VIX の先物をさらに押し上げた。著者らは、この仕組みの欠陥と、その売買が予測できたことを分析した。
\end{paperbox}

ボラティリティの売りは、保険を売る商売に似ている。平時のシャープレシオは非常に高く見えるが、それはまれな大損がまだ起きていないからである。成績を評価するには、シャープレシオではなく、最悪の損失と、その起きる確率を見る必要がある（\chapref{ch:perf}）。

\section{個人や研究で再現できるか}

\begin{itemize}
\item 債券の裁定、転換社債の裁定、ベーシス取引は、価格のデータ（債券や転換社債の日々の価格、レポの金利）、大きなレバレッジ、プライムブローカーの与信が必要で、個人や研究での再現はほぼできない。
\item ボラティリティのリスク・プレミアムの検証は、オプションの価格のデータ（有料）があれば可能である。ただし、まれな大損を含む長い期間で評価しないと、成績を過大に見る。
\item これらの戦略の研究の教訓は、収益の多くが「急落と資金の制約のリスクを引き受ける対価」であり、平時の成績からは、そのリスクが見えないことである。
\end{itemize}

\chapter{高頻度取引とマーケットメイクの研究}
\label{ch:hft}

第 \ref{ch:mm} 章で、板の仕組み、逆選択と在庫のモデルを説明した。この章では、高頻度の取引会社が実際に何をし、どれだけ儲け、市場にどんな影響を与えているかの実証研究をまとめる。

\section{新しいマーケットメイカー}

\begin{paperbox}[Menkveld（2013）「高頻度取引と新しいマーケットメイカー」\cite{menkveld2013}]
欧州の株式で、既存の取引所と、急成長していた新しい取引所（Chi-X）の両方で売買していた一つの大手の高頻度の取引会社の全取引を分析した。この会社は、既存の取引所の取引の 8.1\%、新しい取引所の取引の 64.4\% に参加し、どちらでも取引の約 5 分の 4 は、自分の指値が相手の注文で成立する「受け身の」取引（流動性の供給）だった。在庫の値動きでは損をし、スプレッドで稼いでいた。シャープレシオの計算は、資本のコストをどう仮定するかに強く依存する、と要旨は述べている。
\end{paperbox}

Deep Research の答え 1 は、この論文の値として「1 銘柄・1 日に 1,397 回の取引、1 回 0.88 ユーロの利益、シャープレシオ 9.35」を挙げたが、要旨には出てこず、本文で確かめられなかったので、この報告書では使わない。

この研究は、高頻度取引の大きな部分が、昔の取引所の立会場のマーケットメイカーの役割を、計算機で引き継いだものであることを示した。

\section{流動性と価格の発見}

\begin{itemize}
\item \textbf{Hendershott・Jones・Menkveld（2011）}\cite{hjm2011}：ニューヨーク証券取引所が 2003 年に気配の自動配信を導入し、アルゴリズム取引が増えた変化を使って、アルゴリズム取引がスプレッドを狭め、気配の情報を増やし、流動性を改善したことを示した。
\item \textbf{Brogaard・Hendershott・Riordan（2014）}\cite{bhr2014}：米国株の高頻度の取引会社の売買を分析し、高頻度の取引会社は、全体として、価格の恒久的な変化の方向に売買し、一時的な価格の歪みの反対に売買することで、価格の発見を助けていることを示した。一方で、高頻度の取引会社が流動性を供給する取引は、逆選択を受ける（相手に情報がある）傾向があった。
\end{itemize}

\section{フラッシュ・クラッシュ}

\begin{paperbox}[Kirilenko・Kyle・Samadi・Tuzun（2017）\cite{kkst2017}]
2010 年 5 月 6 日、米国の株価指数の先物（E-mini \SP{}）が数分で急落し、すぐに戻った。先物の取引の口座ごとのデータを分析し、高頻度の取引会社は急落の引き金ではなかったが、大口の売りの注文が市場に入ったとき、在庫を素早く互いに売り渡し（「熱いじゃがいも」）、売りの圧力を増幅したことを示した。高頻度の取引会社の売買のやり方は、急落の日も普段と変わらなかった。
\end{paperbox}

\section{速さの競争}

\begin{paperbox}[Budish・Cramton・Shim（2015）\cite{bcs2015}]
株価指数の先物と、その ETF の価格は、秒の単位ではほぼ完全に連動するが、ミリ秒の単位では連動が崩れ、機械的な裁定の機会が生じる。この機会は速い者だけが取れるので、速さへの投資の競争が起きるが、機会の大きさは速さの競争で小さくならない。連続的に売買を成立させる市場の仕組みの欠陥であり、一定の間隔でまとめて成立させる方式（頻繁なバッチ・オークション）が解決策だと論じた。
\end{paperbox}

\begin{paperbox}[Aquilina・Budish・O'Neill（2022）\cite{abo2022}]
ロンドン証券取引所のすべての注文の記録（成立しなかった注文や、失敗した注文を含む）から、同じ価格の変化に向けて複数の会社が同時に注文を送る「速さの競争」を数えた。FTSE100 の銘柄では、こうした競争が頻繁に起き、取引量のかなりの部分を占める。遅延の裁定は、投資家にとっての取引のコストを押し上げる「税」のように働き、それをなくすと流動性のコストが下がると推計した。
\end{paperbox}

\section{利益の集中}

Baron・Brogaard・Hagströmer・Kirilenko（2019）\cite{bbhk2019} は、スウェーデンの株式市場の取引会社のデータで、高頻度取引の利益が少数の会社に集中し、速さの順位が利益を強く左右し、新しく参入した会社は利益が小さく退出しやすいことを示した。高頻度取引は、規模と速さの経済が強く働く、勝者が多くを取る市場である。

\section{板の情報による短期の予測}

Cont・Kukanov・Stoikov（2014）\cite{cks2014} は、最良気配の注文の量の増減の偏り（注文の流れの不均衡）が、数秒〜数分の価格の変化を線形によく説明し、その傾きは板の厚さに反比例することを示した。短期の価格の予測とインパクトのモデルの基本になっている。

\section{確立していること}

\begin{keybox}[まとめ]
\begin{itemize}
\item 高頻度取引の大きな部分はマーケットメイクであり、平常時には、スプレッドを狭め、価格の発見を助けている。
\item 収益は、極めて多くの小さな取引の積み重ねで、シャープレシオは非常に高いが、容量は小さく、少数の速い会社に集中する。
\item 市場が混乱すると、流動性の供給を引っ込め、急落を増幅することがある。
\item 速さの競争は、市場の仕組み（連続的な売買の成立）が生んでいて、社会的な価値の小さい投資を招いている、という議論がある。
\end{itemize}
\end{keybox}

個人や大学の研究がこの分野で競うことは事実上できないが、決算発表の直後の数秒〜数分は、この人たちが情報を価格に入れる時間であることを知っておく価値はある（第 \ref{ch:mm} 章の「この研究との関係」）。

\chapter{機械学習による予測}
\label{ch:ml}

\chapref{ch:model} で予測のモデルの種類を説明した。この章では、それらを資産の予測に当てはめた主な研究が、何を示したかをまとめる。

\section{比較の基準になった研究}

\begin{paperbox}[Gu・Kelly・Xiu（2020）「機械学習による実証的資産価格」\cite{gkx2020}]
米国株（1957〜2016 年、約 3 万銘柄）の月次のリターンを、94 の銘柄の特性、8 つのマクロ変数、業種、それらの相互作用（合わせて約 900 の説明変数）から予測し、線形回帰、正則化回帰（エラスティックネット）、主成分回帰・PLS、一般化線形モデル、ランダムフォレスト、勾配ブースティング、1〜5 層のニューラルネットを比べた。
\begin{itemize}
\item 標本外の月次の決定係数は、どの方法でも 1\% に満たない。最も良いもの（3 層程度のニューラルネット）で 0.4\% 程度。すべての変数をそのまま入れた最小二乗は、大きく負になる（過学習）。
\item 木のモデルとニューラルネットが最も良く、予測で並べたロング・ショートのポートフォリオは、文献の代表的な回帰の戦略の成績を、場合によっては 2 倍にする。
\item どの方法も、重要な説明変数として、ほぼ同じもの（モメンタムの変種、流動性、変動率）を選ぶ。機械学習の利点は、新しい変数の発見ではなく、既知の変数の非線形な関係と相互作用を捉えることにある。
\item 層の深いネットワークは、浅いものより良くはならない。信号が弱く、データが少ないためである。
\end{itemize}
\end{paperbox}

この論文は、資産の予測の機械学習の研究の基準になった。同時に、「予測できる部分は極めて小さい」ことを、改めて数字で示した。

\section{たくさんの特性をどうまとめるか}

\begin{itemize}
\item \textbf{Kozak・Nagel・Santosh（2020）}\cite{kns2020}：多数の特性のポートフォリオから確率的割引ファクター（すべての資産の価格を説明する一つの量）を推定するとき、少数の特性だけを選ぶ（ラッソのような）疎なモデルではうまくいかず、多くの特性に少しずつ重みをつけ、主成分の分散の小さい方向を強く縮める（リッジに近い）推定が、標本外で最も良いことを示した。ただし、特性を主成分に回転させれば、少数の主成分で十分に説明できる。
\item \textbf{Freyberger・Neuhierl・Weber（2020）}\cite{fnw2020}：特性とリターンの非線形な関係を、ノンパラメトリックな方法（適応的なグループ・ラッソ）で推定し、62 の特性のうち、独立な情報を持つのは十数個にすぎないことを示した。
\item \textbf{Kelly・Pruitt・Su（2019）}\cite{kps2019}：特性を、銘柄のファクターへの感応度（ベータ）を決める変数として使う IPCA を提案した。少数（5 つ程度）の潜在的なファクターで、多くのアノマリーのアルファが説明できる。特性がリターンを予測するのは、それがリスクへの感応度を表すからだ、という解釈である（「特性は共分散である」）。
\item \textbf{Chen・Pelger・Zhu（2024）}\cite{cpz2024}：裁定がない条件を課した深層学習（敵対的な学習）で確率的割引ファクターを推定し、標本外で高いシャープレシオを得た。
\end{itemize}

\section{機械学習の成績は取れるのか}

\begin{paperbox}[Avramov・Cheng・Metzker（2023）\cite{acm2023}]
機械学習の予測による戦略の収益が、裁定の難しい銘柄（極めて小さな株、格付けの低い会社、苦境の会社、流動性の低い株）に集中していることを示した。そうした銘柄を除いたり、経済的な制約（売買のコスト、空売りの難しさ）を課したりすると、収益は大きく下がる。また、機械学習の戦略は回転率が高く、コストの影響が大きい。
\end{paperbox}

Leippold・Wang・Zhou（2022）\cite{lwz2022} は、中国の株式市場で同じ比較をし、機械学習の予測力は米国より強いこと、それが個人投資家の多い市場の性質（小型株、短期の反転）と関係することを示した。どこで予測できるかは、市場の参加者の構成に依存する。

\section{複雑さは良いことか}

Kelly・Malamud・Zhou（2024）\cite{kmz2024} は、説明変数の数が標本の数より多い「過剰にパラメータを持つ」モデル（ランダムな特徴量を極めて多く作り、ほぼ罰則なしで当てはめる）が、標本外で、単純なモデルより良い市場のタイミングの予測をすることを、理論と実証で示し、「複雑さの美徳」と呼んだ。古典的な統計の直感（パラメータは少ないほど良い）に反する結果で、議論を呼んでいる。

\section{機械学習は金融を学べるか}

Israel・Kelly・Moskowitz（2020）\cite{ikm2020} は、運用会社と研究者の立場から、金融の予測に機械学習を使う難しさを整理した。

\begin{enumerate}
\item \textbf{データが小さい}：月次のデータは数百か月しかない。画像や言語のデータの何万分の一である。
\item \textbf{信号が弱い}：予測できる部分は雑音のごく一部。
\item \textbf{市場が変わる}：予測が使われると、関係そのものが消える。
\item \textbf{解釈が要る}：顧客とリスク管理のために、なぜその売買をするのかを説明する必要がある。
\end{enumerate}

結論は、機械学習は役に立つが、その効果は、ほかの分野で見られるような劇的なものではなく、既存のクオンツの方法の「改良」の程度だ、というものである。

\section{価格の図を読む CNN}

\begin{paperbox}[Jiang・Kelly・Xiu（2023）「価格のトレンドを（再）画像化する」\cite{jkx2023}]
過去 5〜60 日の株価と売買高を、ローソク足の図のような画像にし、画像認識に使う畳み込みニューラルネット（CNN）に、その後のリターンの上げ下げを予測させた。CNN の予測は、既知のモメンタムや反転の変数を超える予測力を持ち、学習しなかった期間・国・時間の長さでも同じように機能した。CNN が見つけたパターンは、人間のテクニカル分析の図形とは必ずしも一致しない。
\end{paperbox}

\section{確立していること}

\begin{keybox}[まとめ]
\begin{description}[leftmargin=1em,style=nextline]
\item[確立している]
  標本外の予測の決定係数は月次で 1\% に満たない。非線形のモデル（木、浅いニューラルネット）は線形のモデルより少し良い。重要な変数は、モメンタム・流動性・変動率などの既知のもの。予測の収益は、取引の難しい小型株に偏り、コストで大きく減る。
\item[議論が続いている]
  深い・大きなモデルが本当に良いか（複雑さの美徳）。機械学習の戦略の収益が、公表と普及の後も残るか。
\end{description}
\end{keybox}

\begin{projbox}
LLM-BLM で、ChronoGPT のベクトル（数千次元）をリッジで反応に結びつけたのは、Kozak ほか（2020）の「多くの説明変数に少しずつ重みをつけて縮める」考え方に近い。また、Avramov ほか（2023）の結果は、小型株の S1 で予測の力が見つかっても、それが取れるかどうかは別の問題であることを示している。
\end{projbox}

\chapter{テキストと LLM}
\label{ch:text}

LLM-BLM の研究に最も近い分野である。ニュース、決算発表、決算説明会、提出書類の文章から、株価の予測に使える情報を取り出す研究は、辞書から機械学習、そして LLM へと進んできた。

\section{辞書の時代}

\begin{paperbox}[Tetlock（2007）\cite{tetlock2007}]
Wall Street Journal の株式市場のコラムの文章の悲観の度合い（一般の心理の辞書で数えた否定的な語の割合）が高い日の翌日は、株価が下がり、その後数日で戻ることを示した。悲観の高さは、売買高の多さも予測した。メディアの語調が、投資家の心理を通じて、一時的に価格を動かすという解釈である。
\end{paperbox}

Tetlock・Saar-Tsechansky・Macskassy（2008）\cite{tsm2008} は、個別の会社についてのニュースの否定的な語の割合が、その会社の将来の利益とリターンの低さを予測し、市場はその情報にわずかに遅れて反応することを示した。

\begin{paperbox}[Loughran・McDonald（2011）\cite{lm2011}]
一般の心理の辞書（Harvard の辞書）で「否定的」とされる語の約 4 分の 3 は、10-K（米国の年次報告書）の文脈では否定的な意味を持たない（「負債」「税」「費用」など）ことを示し、金融の文章のための辞書を作った。この辞書による否定的な語の割合は、提出の前後の株価の反応などと関係する。
\end{paperbox}

LLM-BLM の \code{dictionary} の点数は、この Loughran--McDonald の辞書による語調である。

\subsection{文章の変化}

\begin{paperbox}[Cohen・Malloy・Nguyen（2020）「怠惰な価格」\cite{cmn2020}]
10-K と 10-Q の文章が、前年の同じ書類から変わった度合いを測り、文章を大きく変えた会社（変更者）の株が、その後大きく下がることを示した。変更者を売り、変えなかった会社を買うポートフォリオは、月に最大 188 ベーシスポイント（年 22\% 超）のアルファを上げた。投資家は長い書類の変化を読まず、情報は数か月かけて価格に入る。
\end{paperbox}

この研究は、「文章の情報は、数値の情報よりゆっくり価格に入る」ことの強い証拠の一つで、PEAD.txt（第 \ref{ch:event} 章）と同じ方向の結果である。

\section{機械学習の時代}

\begin{itemize}
\item \textbf{Ke・Kelly・Xiu（2019）}\cite{kkx2019}：ニュースの文章から、リターンの符号と一緒に現れやすい語を教師ありで学び、語調の点数を作る方法（SESTM）を提案した。この点数は、既存の商用の語調の指標より強くリターンを予測し、その情報は、小さく変動の大きい株ほどゆっくり価格に入った。
\item \textbf{Frankel・Jennings・Lee（2022）}\cite{fjl2022}：決算説明会や 10-K の文章で、機械学習による語調の測り方が、辞書による測り方より、開示に対する株価の反応をよく説明することを示した。
\item \textbf{FinBERT}：言語モデル BERT を金融の文章で追加学習したもの。Araci（2019）\cite{araci2019} は金融のニュースの語調の分類で、Huang・Wang・Yang（2023）\cite{hwy2023} はアナリストの報告書などで、辞書や従来の機械学習より高い精度を示した。LLM-BLM の \code{finbert} と \code{finbert\_ft} はこの系統である。
\end{itemize}

\section{LLM の時代}

\subsection{LLM はリターンを予測するか}

\begin{paperbox}[Lopez-Lira・Tang（2023）\cite{lopezliratang2023}]
ChatGPT に、ニュースの見出しがその会社の株価に良いか悪いかを答えさせた点数が、その後のリターンを予測することを示した。LLM の言語の理解の能力が、金融の予測の能力として現れたことを示した最初期の研究の一つである。改訂を重ねた arXiv の最新の版の要旨によれば、
\begin{itemize}
\item GPT-4 は、学習の知識の期限より後の見出しで、取引できない最初の反応（見出しが出た時点の反応）の方向を、ポートフォリオの日の約 90\% で当てた。
\item GPT-4 の点数は、その後のドリフトも有意に予測し、特に小さな株と悪いニュースで強かった。
\item 予測の力はモデルが大きいほど高い。
\item LLM の利用が広がるにつれて、戦略の収益は下がった（価格の効率が上がったことと整合的）。
\end{itemize}
「約 90\%」は取引できない最初の反応についての値で、取引できる収益の値ではないことに注意する（Deep Research の答え 2 は、この区別を落としていた）。
\end{paperbox}

\begin{paperbox}[Chen・Kelly・Xiu（2022）「期待リターンと LLM」\cite{ckx2022}]
LLM の点数（語調）ではなく、LLM が文章から作るベクトル（埋め込み）全体を説明変数にして、ニュースからリターンを予測した。世界の多くの市場で、埋め込みによる予測は、語調の点数や従来の言語の方法より強く、大きなモデルほど予測が良かった。
\end{paperbox}

\begin{paperbox}[Kim・Muhn・Nikolaev（2024）\cite{kmn2024}]
GPT-4 に、会社名を伏せた財務諸表（貸借対照表と損益計算書）だけを与えて、来期の利益が増えるか減るかを予測させた。GPT-4 の予測の正しさは、アナリストの予想を上回り、専門の機械学習のモデルに匹敵した。その予測による売買の戦略は、高いシャープレシオとアルファを示した。
\end{paperbox}

\subsection{先読みの問題}

LLM の研究の最大の問題は、\textbf{先読み}（\chapref{ch:backtest}）である。LLM は、学習したデータの期間の出来事を知っている。2015 年の決算発表文を、2023 年までの文章で学習した LLM に読ませると、LLM はその会社のその後の株価や業績を「知っている」可能性がある。

\begin{itemize}
\item \textbf{Glasserman・Lin（2023）}\cite{glasserman2023}：GPT の語調の点数による売買の成績を、元の見出しと、会社名などを伏せた見出しで比べた。LLM の学習の期間の中では、意外にも伏せた見出しのほうが成績が良く、特に大きな会社で強かった。会社についての一般の知識が語調の判断を乱す効果（気をそらす効果）のほうが、先読みより大きかったと解釈した。学習の期間の外では、先読みは問題にならないが、気をそらす効果は残りうる。
\item \textbf{ChronoGPT}（He・Lv・Manela・Wu 2025）\cite{chronogpt2025}：各時点までに手に入った文章だけで学習した言語モデルの系列（ChronoBERT・ChronoGPT）を作り、先読みのない検証を可能にした。ニュースから翌日のリターンを予測する応用では、その場その場の出力によるシャープレシオが、ずっと大きな Llama のモデルと同程度で、先読みの偏りは小さいと報告した。LLM-BLM の主な検定は、2015 年までの文章で学習した版を使った。
\end{itemize}

\subsection{金融に特化した LLM}

\begin{itemize}
\item \textbf{BloombergGPT}（Wu ほか 2023）\cite{bloomberggpt2023}：Bloomberg が、自社の金融のデータ（3,630 億トークン）と一般のデータ（3,450 億トークン）で、500 億のパラメータの LLM を一から学習した。金融の課題では、既存のモデルより大きく良く、一般の課題でも劣らなかった、と報告した。
\item \textbf{FinGPT}（Yang・Liu・Wang 2023）\cite{fingpt2023}：公開の LLM を、金融のデータで安く追加学習する（LoRA などの軽い方法）枠組みを公開した。一から学習するより、はるかに安い。
\end{itemize}

\section{効果は残っているか}

文章の情報の予測力については、次のことが分かっている。

\begin{enumerate}
\item \textbf{発表の直後の反応}は、文章の点数で強く説明できる（LLM-BLM の \SP{} の実験もこれを確認した）。
\item \textbf{翌日以降のドリフト}は、大型株では小さいか、消えている（第 \ref{ch:event} 章の Martineau 2022）。残っているのは、小さな株、悪い知らせ、長く複雑な文章（10-K の変化、決算説明会の発言）である。
\item \textbf{LLM の普及}は、文章の情報の処理を安く速くし、こうした遅れを小さくしていくと予想される。過去のデータで見つかった LLM の戦略の成績が、これからも続くとは限らない。
\item \textbf{先読み}を除くと、学習の期間の中での LLM の成績は、小さくなりうる。
\end{enumerate}

\begin{projbox}
LLM-BLM の結果（2 日の反応は強く説明、ドリフトは有意でない）は、この分野の研究と整合的である。研究として新しい情報を加えられる余地は、(1) 先読みのない LLM（ChronoGPT）での検証、(2) 注目の低い小型株でのドリフト（実験 S1）、(3) 数値でなく文章でしか分からない情報（ガイダンス、リスクの記述、文章の変化）、(4) 発表の当日の反応の大きさの予測（取引ではなく、リスク管理やイベントの前の量の調整に使う）、などにある。
\end{projbox}

\chapter{ポートフォリオ構築・リスク・執行の研究}
\label{ch:pcresearch}

\chapref{ch:pc}〜\chapref{ch:cost} で説明した道具が、実際のデータでどれだけ効いたかの研究をまとめる。

\section{推定誤差と単純な方法}

\begin{paperbox}[DeMiguel・Garlappi・Uppal（2009）\cite{dgu2009}]
14 の平均分散型の最適化の方法（ベイズの縮小、制約つき、最小分散など）を、7 つのデータで、すべての資産に等しく投資する $1/N$ のポートフォリオと比べた。標本外のシャープレシオ・確実性等価・回転率のどれでも、一貫して $1/N$ に勝つ方法はなかった。平均分散の最適化が $1/N$ に勝つには、25 資産なら約 3,000 か月、50 資産なら約 6,000 か月の推定の期間が要る、と計算した。
\end{paperbox}

この結果は、期待リターンの推定の誤差がいかに大きいかを示す。ただし、$1/N$ は「予測を持たない」場合の比較であり、予測の情報を持つ運用者には、Black--Litterman のように予測を控えめに混ぜる方法が向く。

\begin{itemize}
\item \textbf{Jagannathan・Ma（2003）}\cite{jagannathanma2003}：最小分散のポートフォリオに空売りの禁止の制約をかけると、標本の共分散を縮めたのと同じ効果が生じ、標本外の成績が良くなる。「間違った制約が役に立つ」。
\item \textbf{López de Prado（2016）}\cite{hrp2016}：階層的リスクパリティ（\chapref{ch:pc}）が、モンテカルロの実験で、二次計画による最小分散より標本外の分散が小さいことを示した。ただし、実際のデータでの優位は、比較の条件に依存するという報告もあり、確立した結果とは言いにくい。
\end{itemize}

\section{共分散の推定}

\begin{itemize}
\item \textbf{Ledoit・Wolf（2004）}\cite{ledoitwolf2004}：標本の共分散を、すべての相関を等しいとした行列に向けて縮めると、標本外の最小分散のポートフォリオの変動率が下がる。縮める強さは、データから解析的に計算できる。
\item \textbf{Ledoit・Wolf（2020）}\cite{ledoitwolf2020}：固有値ごとに違う強さで縮める非線形の縮小を、解析的な式で速く計算する方法を示した。資産の数が多いときに、線形の縮小より良い。
\item \textbf{ランダム行列}：Laloux ほか（1999）\cite{laloux1999} と Plerou ほか（1999）\cite{plerou1999} は、株の相関行列の固有値の大部分が、ランダムな行列の固有値の分布（Marchenko--Pastur）の範囲に収まり、範囲の外の少数の固有値（市場全体、業種）だけが安定した構造を持つことを示した。
\end{itemize}

\section{ボラティリティの予測と、量の調整}

Hansen・Huang・Shek（2012）\cite{hhs2012} の Realized GARCH は、日中の高頻度のデータで測った実現の変動率を GARCH に組み込み、日次のリターンだけの GARCH より良い変動率の予測をした。Corsi（2009）\cite{corsi2009} の HAR モデルとともに、実現の変動率による予測は、日次のデータだけのモデルより一般に良い。

\begin{paperbox}[Moreira・Muir（2017）と、その批判]
Moreira・Muir（2017）\cite{moreiramuir2017} は、市場やファクターのポートフォリオの量を、前月の実現の分散の逆数に比例して調整する（変動が大きいときに減らす）と、大きなアルファが生じ、シャープレシオが上がることを示した。変動率は予測できるが、平均リターンは変動率に比例して上がらないためである。

Cederburg・O'Doherty・Wang・Yan（2020）\cite{cederburg2020} は、100 を超える株式の戦略で同じ調整を検証し、その場その場で実行できる形（調整の大きさを事前に決める）では、調整した戦略が調整しない戦略に一貫して勝つわけではないことを示した。標本の全期間を使って調整の大きさを決めると、成績が良く見える。
\end{paperbox}

モメンタムのような急落する戦略では、量の調整の効果が大きい（Barroso・Santa-Clara 2015 \cite{barroso2015}、第 \ref{ch:equityfactor} 章）。一方、市場全体の量の調整の効果は、検証の仕方に大きく依存する。

\section{取引コストと容量}

\subsection{インパクトのモデル}

\begin{itemize}
\item \textbf{Almgren・Thum・Hauptmann・Li（2005）}\cite{athl2005}：大手の証券会社の米国株の大口の注文のデータから、一時的なインパクトが売買の速さの約 0.6 乗に比例する（平方根に近い）ことなどを推定した。
\item \textbf{Tóth ほか（2011）}\cite{toth2011}：大口の注文のインパクトが、注文の大きさの平方根に比例することを、多くの市場で確かめ、板に表れない潜在的な流動性が、注文の大きさに比例して薄くなるというモデルで説明した。
\item 平方根則の多くの実証は、Bouchaud ほかの教科書 \cite{bouchaud2018} にまとめられている。
\end{itemize}

\subsection{コストの後の収益と容量}

\begin{itemize}
\item \textbf{Korajczyk・Sadka（2004）}\cite{korajczyksadka2004}：インパクトを含む取引コストを入れると、モメンタムの戦略の収益は、運用の資金が大きくなるにつれて小さくなり、ある資金の大きさ（損益分岐の規模）でゼロになることを示した。取引の量を流動性に合わせる工夫をした一部のモメンタムの戦略では、見かけの収益の機会が消えるまでに最大で約 50 億ドル（1999 年 12 月の時価総額の水準で）を投じられるが、よく研究されてきた従来のモメンタムの戦略の多くは、大きな規模では実行できないと推計した。
\item \textbf{Novy-Marx・Velikov（2016）}\cite{novymarxvelikov2016}：回転率の低いアノマリーはコストの後も残るが、高いものは多くが消える（第 \ref{ch:equityfactor} 章）。
\item \textbf{Frazzini・Israel・Moskowitz（2018）}\cite{fim2018}：大手の運用会社の実際の大量の売買の記録から取引コストを推定し、学術研究で使われてきたコストの推定より、実際のコストはかなり小さいこと、そのため割安・モメンタムなどの戦略の容量は、これまでの推計より大きいことを論じた。運用会社自身の研究であり、執行の上手な大手だから実現できるコストである点に注意が要る。
\item \textbf{DeMiguel・Martín-Utrera・Nogales・Uppal（2020）}\cite{dmnr2020}：多くの特性を同時に使うと、それぞれの売買が一部打ち消し合い（ある特性で買い、別の特性で売る）、コストが下がる。そのため、取引コストを考えると、むしろ多くの特性を使うことが有利になる。
\end{itemize}

\section{強化学習と深層学習}

\subsection{執行}

最適執行（\chapref{ch:cost}）は、状態（残りの量、時間、板の状態）に応じて行動（どれだけ売買するか）を決める逐次の意思決定の問題で、強化学習に向く。

\begin{itemize}
\item \textbf{Nevmyvaka・Feng・Kearns（2006）}\cite{nfk2006}：米国株の板のデータを使い、強化学習で指値の位置と量を決める執行の方策を学習し、単純な方策よりコストを下げた。
\item \textbf{Ning・Lin・Jaimungal（2021）}\cite{nlj2021}：深層の Q 学習（二重 DQN）で最適執行を学習し、TWAP などより良い成績を示した。
\item \textbf{Hafsi・Vittori}\cite{hafsivittori}：多数の売買の主体を模した市場のシミュレータ（ABIDES）の中で、板の状態を入力にした DQN の執行の方策を学習し、TWAP などの単純な方策より執行のコストを下げた（Deep Research の答え 2 が挙げた研究で、内容はその要約による。題名・著者は照合で確かめ、2026 年の国際会議 ACDSA の会議録に載っている。シミュレータの中の結果であり、実際の市場での検証ではない）。
\end{itemize}

執行は、データが豊富で（毎日、多くの注文）、結果がすぐに分かり、目的（コスト）がはっきりしているので、強化学習が実用に近い分野とされる。

\subsection{ヘッジ}

Buehler・Gonon・Teichmann・Wood（2019）\cite{deephedging2019} の「深層ヘッジ」は、取引コストなどの摩擦がある市場で、デリバティブのヘッジの方策をニューラルネットで直接学習する枠組みを示した。実際の銀行のデリバティブの業務にも使われている、と報じられている。

\subsection{ポートフォリオの配分}

Jiang・Xu・Liang（2017）\cite{jxl2017} や、オープンソースの FinRL \cite{finrl2020} など、深層の強化学習でポートフォリオの比率を直接決める研究は多い。しかし、

\begin{itemize}
\item 報酬（リターン）の信号が極めて弱く、学習が不安定で過学習しやすい。
\item 市場が変わるので、過去の環境で学んだ方策が将来も良いとは限らない。
\item 取引コストや市場への影響を、シミュレーションで正しく再現しにくい。
\end{itemize}
という理由で、標本外の頑健な成績を示した研究は少なく、実務での利用も限られる（Deep Research の答え 2 も、研究の段階と評価していた）。

\begin{keybox}[まとめ]
\begin{description}[leftmargin=1em,style=nextline]
\item[確立している]
  期待リターンの推定の誤差は非常に大きく、単純な方法（$1/N$、縮小、制約）が強い。共分散の縮小推定と実現の変動率による予測は有効。インパクトは平方根則に近い。コストは戦略の容量を決め、回転率の高い戦略の収益の多くを消す。強化学習は執行では実用に近い。
\item[議論が続いている]
  変動率による量の調整の効果（検証の仕方に依存）。階層的リスクパリティの実データでの優位。実際の取引コストの大きさ（学術の推定と運用会社の推定の差）。強化学習による配分。
\end{description}
\end{keybox}

\chapter{ヘッジファンドは儲かっているのか}
\label{ch:hfperf}

第 \ref{ch:equityfactor}〜\ref{ch:hft} 章は「戦略」の研究だった。この章は、実際のヘッジファンドの「成績」の研究をまとめる。結論を先に書くと、平均的なヘッジファンドの手数料の後のアルファは小さく、データの偏りを補正するとさらに小さくなる。少数のファンドには本物の腕があるらしいが、それを事前に見分けるのは難しい。

\section{ヘッジファンドのリターンを説明するファクター}

ヘッジファンドは、空売り・レバレッジ・デリバティブを使い、ポジションを頻繁に変えるので、株式の投資信託の評価に使うファクター（市場、規模、割安、モメンタム）では、そのリスクを捉えきれない。

\begin{itemize}
\item \textbf{Fung・Hsieh（1997）}\cite{funghsieh1997}：ヘッジファンドのリターンは、伝統的な資産クラスとの相関が低く、少数の「スタイル」（売買の仕方）の要因で説明できることを示した。
\item \textbf{Fung・Hsieh（2001）}\cite{funghsieh2001}：トレンドフォローの運用者のリターンが、ルックバック・ストラドル（オプション）のリターンに似ることを示した（第 \ref{ch:trend} 章）。
\item \textbf{Agarwal・Naik（2004）}\cite{agarwalnaik2004}：多くの株式のヘッジファンドの戦略のリターンが、株価指数のアウト・オブ・ザ・マネーのプット・オプションを売った形に似る（左の裾のリスクを持つ）ことを示し、オプションの形のファクターを提案した。平均と分散だけの評価は、この裾のリスクを過小評価する。
\item \textbf{Fung・Hsieh（2004）}\cite{funghsieh2004}：株式市場、規模の差、国債の利回りの変化、信用のスプレッドの変化、債券・通貨・商品のトレンドフォロー（ルックバック・ストラドル）の 7 つのファクターのモデルを提案した。このモデルは、ヘッジファンドの指数のリターンの変動の大部分（約 8 割）を説明する。ヘッジファンドの研究の標準のベンチマークになった。
\end{itemize}

\begin{projbox}
「ヘッジファンドの収益の多くは、ルールで再現できるファクター（オルタナティブ・ベータ）で説明できる」という結果（第 \ref{ch:hf} 章の 3 分類の 2）は、この一連の研究に基づく。
\end{projbox}

\section{手数料の後のアルファ}

\subsection{メタ分析}

\begin{paperbox}[Yang・Havranek・Irsova・Novak（2022）\cite{yhin2022}]
2001〜2021 年に発表された 74 の研究から、ヘッジファンドのリターンをリスク・ファクターに回帰した切片（アルファ）の推定値 1,019 個を集めたメタ分析。
\begin{itemize}
\item 報告されたアルファの単純な平均は月 0.36\%（年 4.32\%）。発表の偏り（有意な結果だけが発表されやすい偏り）は小さく、それを補正した推定の多くは月 0.30〜0.40\% の範囲にある。
\item データの偏り（遡って成績が入る偏り、生存者バイアス）を明示的に補正した研究ほど、低いアルファを報告している。
\end{itemize}
\end{paperbox}

Deep Research の答え 3 は、この論文が「データの期間の中央が 2015 年以降の研究では、平均のアルファが負になる」ことを示したと書いたが、照合した 2022 年 7 月の版の要旨と本文の検索では、その記述を見つけられなかった。この報告書では使わない。

\subsection{時期による変化}

\begin{paperbox}[Fung・Hsieh・Naik・Ramadorai（2008）\cite{fhnr2008}]
1995〜2004 年のファンド・オブ・ファンズを分析し、平均的なファンド・オブ・ファンズが有意なアルファを上げたのは、1998 年 10 月〜2000 年 3 月の短い期間だけだったことを示した。アルファを上げたファンドには資金が流入し、流入の後にアルファは下がった。資金の流入が、アルファの機会を食いつぶしていく（資本の形成）。
\end{paperbox}

\subsection{ラグのある市場の感応度}

Asness・Krail・Liew（2001）\cite{akl2001} は、ヘッジファンドの月次のリターンを、同じ月の市場のリターンだけでなく、前の月の市場のリターンにも回帰すると、市場への感応度（ベータ）がずっと大きくなり、アルファの多くが消えることを示した。流動性の低い資産の評価が遅れて市場の動きを反映するため（\ref{sec:smoothing} 節の平滑化）である。Hasanhodzic・Lo（2007）\cite{hasanhodziclo2007} は、いくつかの戦略の種類で、ヘッジファンドのリターンのかなりの部分が、株・債券・通貨・商品・信用などの共通のファクターへの線形の感応度で再現できる（「クローン」が作れる）ことを示した。

\subsection{投資家が実際に得た収益}

\begin{paperbox}[Dichev・Yu（2011）\cite{dichevyu2011}]
ファンドの成績（買って持ち続けた場合の収益）ではなく、投資家の資金の出し入れを考慮した収益（金額で加重した収益）を計算した。投資家は成績の良い後に資金を入れ、悪い後に引き出すので、金額で加重した年率の収益は、買って持ち続けた収益より数\% 低く（SSRN に出た版の要旨では約 4\%）、幅広い株価指数の収益も下回り、変動も大きかった。ヘッジファンドの投資家の実際の経験は、それまで考えられていたよりずっと悪い。

\end{paperbox}

\section{データベースの偏り}

ヘッジファンドには成績の公開の義務がなく、データベースへの報告は任意である。このため、データベースの成績は、実際の平均より良く見える。

\begin{table}[H]
\centering\small
\caption{ヘッジファンドのデータベースの偏り}
\label{tab:hfbias}
\begin{tabular}{@{}P{0.2\linewidth}P{0.42\linewidth}P{0.3\linewidth}@{}}
\toprule
偏り & 仕組み & 主な研究 \\
\midrule
生存者バイアス & 閉じた・報告をやめたファンドがデータから消える。 & Liang（2000）\cite{liang2000}、Brown・Goetzmann・Ibbotson（1999）\cite{bgi1999}、Fung・Hsieh（2000）\cite{funghsieh2000} \\
遡って入る偏り & 新しいファンドは、成績が良かったときだけデータベースに加わり、過去の成績も一緒に入る（インスタント・ヒストリー）。 & Fung・Hsieh（2000）、Malkiel・Saha（2005）\cite{malkielsaha2005} \\
自己選択 & 報告を始めるのは成績が良い後、やめるのは悪い後。 & Agarwal・Fos・Jiang（2013）\cite{afj2013}、Aiken・Clifford・Ellis（2013）\cite{ace2013} \\
評価の平滑化・操作 & 流動性の低い資産の評価をならす。小さな損失を報告したがらない。 & Getmansky・Lo・Makarov（2004）\cite{glm2004}、Bollen・Pool（2009）\cite{bollenpool2009} \\
\bottomrule
\end{tabular}
\end{table}

偏りの大きさの推定は、データベース・期間・方法で違う。Yang ほか（2022）のまとめによれば、生存者バイアスの推定は年 2\% 前後から 3\% 台（例えば、ある研究は年 1.8〜2.4\%、別の研究は年 2.0〜3.6\%）で、小さく若くレバレッジの高いファンドではさらに大きい。Liang（2000）\cite{liang2000} は年 2\% を超える生存者バイアスを報告した。Malkiel・Saha（2005）\cite{malkielsaha2005} は、遡って入る偏りと生存者バイアスがともに大きく、ヘッジファンドのリターンは一般に思われているより低く、リスクは高いと結論した。

\begin{itemize}
\item \textbf{報告をやめたファンド}：Aiken・Clifford・Ellis（2013）\cite{ace2013} は、ファンド・オブ・ファンズの保有の開示から、データベースに報告していないファンドの成績を調べ、データベースから抜けたファンドのその後の成績が非常に悪いことを示した（成績の分布の左の裾がデータから消えている）。
\item \textbf{ゼロの前後の不自然さ}：Bollen・Pool（2009）\cite{bollenpool2009} は、ヘッジファンドの月次のリターンの分布で、小さな損失が小さな利益より不自然に少ないこと（ゼロでの不連続）を示し、損失の報告を避ける評価の操作の可能性を指摘した。
\end{itemize}

\begin{projbox}
LLM-BLM の小型株のデータの作り方（生き残った会社だけを選び、その偏りを SEC のデータで別に見積もった）は、ヘッジファンドのデータベースの生存者バイアスと同じ問題への対処である。偏りは比べる相手との比較で一部打ち消せるが、選び方そのものの偏りは残る。
\end{projbox}

\section{腕は続くか}

\begin{itemize}
\item \textbf{Kosowski・Naik・Teo（2007）}\cite{knt2007}：ブートストラップとベイズの方法で、上位のヘッジファンドのアルファが運だけでは説明できないこと、1 年の単位で成績の持続があることを示した。
\item \textbf{Jagannathan・Malakhov・Novikov（2010）}\cite{jmn2010}：データベースの偏りと推定の誤差を考慮しても、優れたファンドの成績は 3 年の単位で持続し（「熱い手」）、劣ったファンドの持続は弱いことを示した。
\item \textbf{Sun・Wang・Zheng}\cite{swz2018}：ヘッジファンドの業界全体が不調だった時期の成績（DownsideReturns）で並べると、上位の 5 分の 1 は下位の 5 分の 1 を、リスクを調整した翌年の成績で約 7\% 上回り、その予測力は 3 年まで続いた。好調だった時期の成績で並べても、その後の差は出なかった。上げ相場では、リスクを取っただけのファンドも良く見えるので、腕と運を見分けにくい（Deep Research の答え 3 は、この論文の題名を誤って書いていた。付録~\ref{ch:verify}）。
\item \textbf{Titman・Tiu（2011）}\cite{titmantiu2011}：ファクターとの連動の小さい（決定係数の低い）ファンドほど、シャープレシオとアルファが高く、手数料も高いことを示した。腕のある運用者は、ファクターのリスクを打ち消し、固有の収益を売る。
\end{itemize}

\section{規模と成績}

\begin{itemize}
\item \textbf{Berk・Green（2004）}\cite{berkgreen2004}：投資信託の理論で、腕のある運用者には資金が集まり、規模が大きくなるとアルファが下がる（規模の不経済）ので、均衡では投資家の手数料の後のアルファはゼロになり、成績の持続も見えなくなることを示した。腕は存在しても、その価値は運用者が手数料として取る。
\item \textbf{Pástor・Stambaugh・Taylor（2015）}\cite{pst2015}：投資信託で、運用会社の規模より、業界全体の規模の拡大がアルファを下げること（業界の規模の不経済）を示した。運用者の腕は時代とともに上がっているが、業界の拡大がそれを打ち消している。
\item \textbf{Getmansky（2012）}\cite{getmansky2012}：ヘッジファンドの成績と規模の関係は凹で、最適な規模があることを示した。
\item \textbf{Yin（2016）}\cite{yin2016}：ヘッジファンドの成績は規模とともに下がるが、運用者の報酬（管理報酬は規模に比例）は規模とともに増えるので、運用者は投資家にとって最適な規模を超えて資金を集める誘因を持つことを示した。
\end{itemize}

\section{手数料は誰のものか}

\begin{paperbox}[Ben-David・Birru・Rossi（2020）「成功報酬の成績」\cite{bbr2020}]
1995〜2016 年の約 22 年間のヘッジファンドで、実際に支払われた成功報酬は、利益の合計に対して、契約上の平均（約 20\%）の約 2.5 倍、つまり約 50\% に達した。投資家は、安全な金利を超えて得た収益（リスクの調整の前）の 1 ドルのうち、約 36 セントしか手にしなかった。原因は、(1) 成功報酬は利益が出た年に払われ、後の損失で返されない（契約の非対称）、(2) 投資家が成績を追って資金を出し入れする、(3) 損をしてハイウォーターマークの下に沈んだファンドが閉じる（その後の回復で報酬が相殺されない）、ことである。ファンドの生涯の成績と、受け取った成功報酬の関係は弱い。
\end{paperbox}

\begin{figure}[H]
\centering
\begin{tikzpicture}[x=1.2mm,y=1mm]
\fill[c1] (0,0) rectangle (36,10); \node[font=\footnotesize,white] at (18,5) {投資家 36 セント};
\fill[c2!80] (36,0) rectangle (100,10); \node[font=\footnotesize,white] at (68,5) {運用者（管理報酬＋成功報酬）64 セント};
\draw[ink2] (0,0) rectangle (100,10);
\node[lab,anchor=north west] at (0,-1) {安全な金利を超えた、手数料の前の収益 1 ドル（1995〜2016 年、リスクの調整の前）};
\end{tikzpicture}
\caption{ヘッジファンドの収益の分け前（Ben-David・Birru・Rossi の NBER の要旨の値から描いた）。}
\label{fig:feesplit}
\end{figure}

\ref{sec:fees} 節の手数料の例（手数料の前の収益の 36\% が運用者に渡る）は 1 年の計算だった。複数年・多数のファンド・投資家の行動を考えると、運用者の取り分はさらに大きくなる。

\section{まとめ}

\begin{keybox}[ヘッジファンドの成績について確立していること]
\begin{itemize}
\item ヘッジファンドのリターンの大部分は、オプションの形を含むファクター（Fung・Hsieh の 7 ファクターなど）で説明できる。
\item 報告されたアルファの平均は月 0.3〜0.4\% 程度だが、データの偏りを補正すると小さくなる。データベースの偏り（生存者・遡り・自己選択）は年に数\% の大きさになりうる。
\item 投資家が実際に得た収益は、資金の出し入れのタイミングのため、ファンドの成績よりさらに低い。
\item 手数料は、手数料の前の超過収益の大部分を取る。
\item 少数の優れたファンドの成績には持続があるが、規模の拡大と資金の流入がアルファを小さくする。
\end{itemize}
\end{keybox}

LLM-BLM の研究が時価総額加重に勝てなかったことは、プロの世界の平均的な成績から見て、驚くべきことではない。プロの多くも、手数料とコストの後では、市場と同じリスクのファクターの組み合わせに勝てていない。

\chapter{失敗と混雑}
\label{ch:failures}

ヘッジファンドの歴史に残る失敗は、たいてい同じ形をしている。平時には安定して儲かる取引に、大きなレバレッジをかけ、同じ取引をしている人が多く、売りたいときに売れない。この章では、代表的な失敗を並べ、その共通の構造と、研究から分かる教訓をまとめる。

\section{共通の構造：解消のらせん}

\begin{figure}[H]
\centering
\begin{tikzpicture}[node distance=8mm and 12mm]
\node[box=c2,text width=30mm] (loss) {損失};
\node[box=c4,text width=30mm,right=24mm of loss] (margin) {担保の追加の要求・\\投資家の解約・\\リスクの上限};
\node[box=c7,text width=30mm,below=12mm of margin] (sell) {ポジションの\\投げ売り};
\node[box=c1,text width=30mm,below=12mm of loss] (price) {価格の下落\\（同じ取引をしている\\人も損をする）};
\draw[arr] (loss)--(margin);
\draw[arr] (margin)--(sell);
\draw[arr] (sell)--(price);
\draw[arr] (price)--(loss);
\node[lab,align=center] at ($(loss)!0.5!(sell)$) {レバレッジ・混雑・\\流動性の不一致が\\回転を速くする};
\end{tikzpicture}
\caption{解消のらせん（Brunnermeier・Pedersen 2009 \cite{brunnermeierpedersen2009}、Pedersen 2009 \cite{pedersen2009} の考え方の模式図）。}
\label{fig:spiral}
\end{figure}

\begin{itemize}
\item \textbf{レバレッジ}：小さな価格の変化で資本の大きな割合を失い、担保の追加を求められる。
\item \textbf{混雑}：同じ取引をしている人が多いと、一人の売りがほかの人の損失を生み、売りの連鎖が起きる。Stein（2009）\cite{stein2009} は、洗練された裁定者が、何人が同じ取引をしているかを知らないこと自体が、価格を不安定にすることを論じた。
\item \textbf{流動性の不一致}：持っている資産はすぐに売れないのに、資金（借入れ、投資家の資金）はすぐに引き上げられる。
\item \textbf{モデルの前提の外れ}：過去の相関・変動率・流動性を前提にしたリスクの物差し（VaR など）は、危機の局面では役に立たない。
\end{itemize}

Pedersen（2009）\cite{pedersen2009} は、これを「皆が同時に出口に走るとき」と呼び、2007〜2008 年の出来事をこの枠組みで説明した。

\section{LTCM（1998 年）}

LTCM は 1994 年に、Salomon Brothers の債券の裁定の部門の出身者が設立し、のちにノーベル経済学賞を受ける Merton と Scholes も加わった。主な取引は、債券の相対価値（第 \ref{ch:rv} 章）で、利回りの差が縮むことに賭けた。

\begin{itemize}
\item 1998 年の初めに、資本の 25 倍前後の資産を持ち、デリバティブの想定元本はそれよりはるかに大きかった。
\item 1998 年 8 月のロシアの債務不履行の後、投資家が一斉に安全で流動性の高い資産（最新の米国債など）に逃げ、LTCM が賭けていたほぼすべての差が同時に広がった。
\item ほかの多くの銀行・ヘッジファンドが似たポジションを持っていて、LTCM の窮状が知られると、その取引の反対側に立つ人が現れた。
\item 9 月、ニューヨーク連邦準備銀行の仲介で、大手の金融機関が資本を注入し、ファンドは管理の下で清算された。
\end{itemize}

Jorion（2000）\cite{jorion2000} は、LTCM のリスク管理を分析し、短い期間の平穏なデータで推定した変動率と相関に基づく VaR が、危機の局面のリスクを大きく過小評価したこと、レバレッジが小さな価格の変化を致命的にしたこと、流動性のリスクと、ほかの参加者の行動のリスクが測られていなかったことを指摘した。Edwards（1999）\cite{edwards1999} は、この出来事が金融システムに与えた影響と、規制の論点を整理した。

\section{Amaranth（2006 年）}

Amaranth Advisors は、マルチストラテジーのファンドだったが、一人のトレーダーの天然ガスの先物の取引（冬と春の限月の価格差に賭けるスプレッドの取引）が、ファンドの資産の大部分のリスクを占めるようになっていた。2006 年 9 月、価格差が逆に動き、数週間で 60 億ドルを超える損失を出して、ファンドは閉じた。

Chincarini（2007）\cite{chincarini2007} と Till（2007）\cite{till2007} は、ポジションが市場全体の建玉のかなりの割合を占めるほど大きく、損を限るために売ろうとしても、自分の売りが価格をさらに不利に動かす状態にあったこと、リスクの物差しがこの流動性の問題を捉えていなかったことを分析した。一つのファンドの中でも、一つの取引への集中と、市場の厚さを超えるポジションが、致命的になりうる。

\section{2007 年 8 月と、2008 年}

2007 年 8 月のクオンツの急落（第 \ref{ch:statarb} 章、Khandani・Lo 2011 \cite{khandanilo2011}）は、市場全体が平穏な中で、混雑したポジションの解消の連鎖が起きた例である。

2008 年の金融危機では、ヘッジファンドの業界の全体が大きく損をした（業界の指数で年 2 割前後の損失と広く報じられた）。

\begin{itemize}
\item \textbf{Ben-David・Franzoni・Moussawi（2012）}\cite{bfm2012}：2008 年の後半に、ヘッジファンドは株の保有を大きく減らした（約 3 割）。売りは、投資家の解約と、資金の貸し手からの圧力によるもので、ほかの機関投資家より大きかった。
\item \textbf{Aragon・Strahan（2012）}\cite{aragonstrahan2012}：Lehman Brothers をプライムブローカーにしていたファンドは、資産が凍結されて売買できなくなり、破綻しやすくなった。そうしたファンドが持っていた株は、流動性が大きく下がった。ヘッジファンドが平時に流動性の供給者として果たしている役割を、逆側から示した。
\item \textbf{Mitchell・Pulvino（2012）}\cite{mitchellpulvino2012}：プライムブローカーが資金の提供を引き上げたことが、転換社債の裁定や合併裁定の価格の崩れを引き起こした（第 \ref{ch:rv} 章）。
\end{itemize}

大手のマルチストラテジーのファンドの中にも、2008 年に資産の半分前後を失ったものがあった、と報じられている。

\section{ボラティリティの売り（2018 年 2 月）と国債（2020 年 3 月）}

2018 年 2 月の VIX の急騰による、ボラティリティの売りの上場商品の崩壊（Augustin ほか 2021 \cite{volmageddon2021}）と、2020 年 3 月の国債の市場の混乱でのヘッジファンドの投げ売り（Kruttli ほか 2021 \cite{kruttli2021}）は、第 \ref{ch:rv} 章で扱った。どちらも、仕組み（日々のリバランス、担保の条件）が、価格の変化を増幅する方向の売買を機械的に生んだ。

\section{クオンツの冬（2018〜2020 年）}

割安を中心とする多くのファクターが同時に長く損をした（Blitz 2021 \cite{blitz2021}、第 \ref{ch:equityfactor} 章）。数日の急落ではなく、数年にわたる低迷で、多くのクオンツのファンドが資金を失った。リスク・プレミアムの戦略でも、何年も損が続くことがあり、そのあいだ投資家が資金を預け続けるとは限らない。

\section{2021 年：空売りの踏み上げと、隠れたレバレッジ}

\begin{itemize}
\item \textbf{GameStop}：2021 年 1 月、インターネットの掲示板で集まった個人投資家が、空売りの多い GameStop の株を買い上げ、株価は数週間で何十倍にもなった。大きな空売りを持っていたヘッジファンドの Melvin Capital は、1 月だけで資産の半分を超える損失を出し、ほかのファンドから資本の注入を受けたと報じられた。空売りには、損失に上限がないという固有のリスクがあり、空売りの残高そのものが「混雑」の指標になる。
\item \textbf{Archegos}：2021 年 3 月、ファミリー・オフィス（一族の資産の運用会社）の Archegos Capital Management が、複数の銀行と結んだトータル・リターン・スワップ（株を持たずに、株の値動きを受け取る契約）で、少数の銘柄に巨大なレバレッジのポジションを持っていた。株価が下がって担保の追加に応じられなくなると、銀行が一斉に担保の株を売り、銀行の損失は合わせて 100 億ドルを超えた（Credit Suisse だけで約 55 億ドルと報じられた）。スワップによるポジションは開示されず、各銀行は、ほかの銀行が同じ顧客に同じポジションを持たせていることを知らなかった。
\end{itemize}

これらは、学術研究の分析がまだ少ない最近の出来事で、事実関係は報道と当局の報告に基づく。

\section{混雑の測り方}

混雑は直接は見えないが、いくつかの測り方が提案されている。

\begin{itemize}
\item \textbf{保有に基づく}：機関投資家の保有の開示（米国の 13F）、空売りの残高、ファンドの資金の流れから、同じ銘柄・同じファクターに集まっている資金を測る。
\item \textbf{リターンに基づく}：Lou・Polk（2022）\cite{loupolk2022} は、モメンタムの勝ち組（あるいは負け組）の銘柄どうしの、ファクターで説明できないリターンの相関（コモメンタム）が、モメンタムの裁定者の多さを表すとした。コモメンタムが高い時期の後には、モメンタムの収益が低く、のちに反転が起きる。裁定者が多すぎると、価格は行き過ぎる。
\item \textbf{価格に基づく}：ファクターの上位と下位の銘柄の評価の差（割安の差の広がり）。混雑すると評価の差が縮む。
\end{itemize}

\section{効果の減衰}

公表の後の減衰（McLean・Pontiff 2016 \cite{mcleanpontiff2016}、第 \ref{ch:equityfactor} 章）は、「ゆっくりした混雑」である。Falck・Rej・Thesmar（2021）\cite{falck2022} は、公表されたアノマリーの標本外のシャープレシオの低下を、事前に分かる特徴で予測した。公表の年だけで低下の分散の 30\% を説明し、新しく公表されるファクターほど低下が大きい（毎年 5 ポイントずつ大きくなる）。シグナルの計算に要する操作の数や、標本の中のシャープレシオの外れ値への敏感さ（過学習に関係する変数）が、さらに 15\% を説明した。裁定に関係する変数の説明力は小さかった。

\section{教訓}

\begin{keybox}[失敗の研究から分かること]
\begin{enumerate}
\item 平時のシャープレシオは、まれな大損のリスクを示さない。最悪の場合の損失と、それに耐えられる資金の構造（レバレッジ、資金の引き上げの条件）を見る。
\item 自分の戦略が正しくても、同じ取引をしているほかの人の事情で損をする。有名な論文のシグナル、広く使われる道具（LLM の語調を含む）ほど、混雑しやすい。
\item 持っている資産の流動性と、資金の流動性を合わせる。
\item 過去のデータで推定したリスクの物差しは、危機の局面で外れる。ストレス・テストで補う。
\item 公表された戦略の効果は、時間とともに小さくなる。新しいものほど、早く小さくなる。
\end{enumerate}
\end{keybox}

\chapter{この研究への含み}
\label{ch:implications}

最後に、ここまでの基本と研究を、LLM-BLM の研究に当てはめる。

\section{LLM-BLM で分かったこと}

\begin{itemize}
\item \SP{} の決算発表文では、LLM（gpt-4o-mini）・FinBERT・ChronoGPT の読み取りは、発表をはさむ 2 日の株価の反応を強く説明した。
\item しかし、その情報は発表の当日にほぼ株価に入りきり、翌日以降のドリフトは有意に予測しなかった。
\item その点数を Black--Litterman の見通しにした月次・日次のロングのみのポートフォリオは、取引コストの後で時価総額加重を上回れなかった。
\item 注目の低い小型株で、文章の読み取りが発表の後 2〜60 取引日のドリフトを予測するかを、事前登録した実験 S1 で確かめている（この報告書を書いた時点で、結果はまだ出ていない）。
\end{itemize}

\section{研究の文献の中での位置}

\begin{table}[H]
\centering\small
\caption{LLM-BLM の結果と、関係する研究}
\label{tab:position}
\begin{tabular}{@{}P{0.32\linewidth}P{0.6\linewidth}@{}}
\toprule
LLM-BLM の結果 & 整合する研究 \\
\midrule
文章の点数は 2 日の反応を強く説明する & Tetlock ほか（2008）、Loughran・McDonald（2011）、Lopez-Lira・Tang（2023）、Chen・Kelly・Xiu（2022）（第 \ref{ch:text} 章） \\
翌日以降のドリフトは残らない & Martineau（2022）：大型株の PEAD は 2006 年以降なく、最近はマイクロキャップでも消えた。Lopez-Lira・Tang：LLM の普及で戦略の収益が低下（第 \ref{ch:event}・\ref{ch:text} 章） \\
時価総額加重に勝てない & ヘッジファンドの平均の手数料の後のアルファは小さい（第 \ref{ch:hfperf} 章）。公表されたシグナルは減衰する（第 \ref{ch:equityfactor}・\ref{ch:failures} 章） \\
小型株ならドリフトが残るかもしれない（S1 の仮説） & Hong・Lim・Stein（2000）、Chordia ほか（2009）、Akbas ほか、PEAD.txt（第 \ref{ch:event} 章）。ただしコストで消えうる（Chordia ほか 2009、Avramov ほか 2023） \\
\bottomrule
\end{tabular}
\end{table}

LLM-BLM の結果は、文献の主な流れとよく一致している。「LLM を使えば市場に勝てる」という期待に対して、「LLM は情報を正しく読むが、その情報はすでに価格に入っている」という、効率的な市場の考え方に沿った答えである。

\section{プロの工程との違い}

\begin{table}[H]
\centering\small
\caption{クオンツの運用の工程と、LLM-BLM の設計の違い}
\label{tab:gap}
\begin{tabular}{@{}P{0.14\linewidth}P{0.27\linewidth}P{0.24\linewidth}P{0.26\linewidth}@{}}
\toprule
工程 & プロの典型 & LLM-BLM & 差の影響 \\
\midrule
シグナル & 数十〜数百を組み合わせ、互いの相関を考える & LLM などの読み取りの一つ & 一つの弱いシグナルでは IR が小さい（基本法則の BR が小さい） \\
速さ & 発表の数ミリ秒〜数分で売買 & 月末・日次の終値 & 情報の大部分は売買の前に価格に入る \\
期間 & シグナルの減衰の曲線に合わせる & 月次・日次・四半期 & 速いシグナルを遅い売買で使うと、価値が消える \\
構築 & ロング・ショート、ファクター中立、TC を高く & ロングのみ、TE 3\% & 悪い読み取りを空売りに使えず、TC が低い \\
コスト & 平方根則のモデル、売買しない領域、執行の最適化 & 片道 0.3\% の一定 & 小さな資金なら妥当。大きな資金では過小 \\
リスク & ファクター型のリスクモデル、ストレス・テスト & Ledoit--Wolf の共分散 & 小規模なら十分 \\
検証 & 長いホールドアウト、試した数の記録、紙の上の運用 & 事前登録（S1）、Holm の補正 & 研究としては厳しい側 \\
\bottomrule
\end{tabular}
\end{table}

\section{研究で試せること}

ここまでの章から、学生の研究の規模で試す価値がありそうなことを、優先度の高い順に挙げる。どれも、新しい分析を始める前に事前登録の追記として書くべきものである。

\begin{enumerate}
\item \textbf{シグナルの価値を、ポートフォリオの制約と分けて測る}。ロングのみ・TE の上限の中での成績だけでなく、紙の上のロング・ショート（上位 − 下位の分位のポートフォリオ）の成績と、その移転係数（TC）を報告する。予測の価値がないのか、ロングのみの制約で価値が失われたのかが分かる（\chapref{ch:perf}、\chapref{ch:pc}）。
\item \textbf{減衰の曲線を描く}。発表の翌日の始値、翌日の終値、1 週間後、1 か月後に売買したとき、それぞれの IC と、コストの後の収益がどう変わるか。情報がどの時間の単位で価格に入るかを示せば、「なぜ月次の運用では勝てないか」を直接に説明できる（\chapref{ch:signal}）。
\item \textbf{既知のファクターへの増分を測る}。LLM の点数を、モメンタム・割安・収益性・決算の驚き（SUE）などに回帰した残差の IC を測る。LLM の点数が、数値の驚きの言い換えにすぎないのか、文章でしか分からない情報を持つのかが分かる（第 \ref{ch:equityfactor}・\ref{ch:event} 章）。
\item \textbf{反応の遅れやすい場面に絞る}。小型株（S1）のほか、アナリストの少ない銘柄、金曜日や発表の多い日、悪い知らせ、長く複雑な文章（ガイダンス、リスクの記述、前回からの文章の変化）。先行研究が遅れを示した場面である（第 \ref{ch:event}・\ref{ch:text} 章）。
\item \textbf{読み取りをリスク管理に使う}。2 日の反応の\textbf{大きさ}を予測できるなら、それは売買の向きの予測でなくても、発表の前後の量の調整や、変動率の予測に価値がある（\chapref{ch:risk}）。
\item \textbf{コストと容量を明示する}。小型株では、平方根則（\chapref{ch:cost}）で、資金の規模ごとのコストを見積もり、「どの規模の資金で、この結論が成り立つか」を書く。
\item \textbf{先読みを避け続ける}。学習の期間の止まった言語モデル（ChronoGPT）や、会社名を伏せた入力（Glasserman・Lin 2023）で、LLM が将来を知っている影響を除く（第 \ref{ch:text} 章）。
\end{enumerate}

\section{してはいけないこと}

\begin{itemize}
\item \textbf{成績が出るまで条件を変えて試し続ける}。100 通りを試せば、無意味な戦略からでも年率 1 を超えるシャープレシオが見つかる（\chapref{ch:backtest}）。試した数を記録し、事前登録にない分析は探索として区別する。
\item \textbf{レバレッジやオプションの売りで、シャープレシオを良く見せる}。平時のシャープレシオは、まれな大損を示さない（第 \ref{ch:rv}・\ref{ch:failures} 章）。
\item \textbf{今日残っている会社だけで、過去の成績を評価する}。生存者バイアスは、年に数\% の大きさになりうる（第 \ref{ch:hfperf} 章）。LLM-BLM の小型株のデータでは、この偏りを別に見積もった。
\end{itemize}

\section{結び}

ヘッジファンドの研究から分かるのは、市場に勝つことは、プロにとっても非常に難しいということである。勝っている少数の運用者は、(1) 多くの弱いシグナルを組み合わせ、(2) 速く、安く売買し、(3) ほかの人が取りたがらないリスク（流動性、急落、混雑の解消）を、それに耐えられる資金の構造で引き受け、(4) 新しいデータと方法を探し続けている。一つの新しい道具（LLM）だけで市場の平均を上回ることは、文献から見て、期待しにくい。

一方で、研究としての価値は、勝てたかどうかだけにあるのではない。LLM が決算発表文の何を読み、その情報がどの速さで価格に入り、どこに遅れが残るのかを、先読みと過学習を避けて正しく測ることは、それ自体が市場の効率についての新しい証拠になる。


\appendix
\chapter{用語の一覧}
\label{ch:glossary}

本文で使った主な用語を、英語の名前と、説明した章とともに並べる。

{\small
\begin{longtable}{@{}P{0.2\linewidth}P{0.21\linewidth}P{0.42\linewidth}P{0.08\linewidth}@{}}
\toprule
用語 & 英語 & 意味 & 章 \\
\midrule
\endfirsthead
\toprule
用語 & 英語 & 意味 & 章 \\
\midrule
\endhead
\bottomrule
\endfoot
\multicolumn{4}{@{}l}{\textbf{ヘッジファンドの仕組み}} \\
プライムブローカー & prime broker & ヘッジファンドに資金・株を貸し、売買の決済と保管をする投資銀行の部門 & \ref{ch:hf} \\
管理報酬・成功報酬 & management / performance fee & 資産に比例する報酬と、利益に比例する報酬。伝統的には 2\% と 20\% & \ref{ch:hf} \\
ハイウォーターマーク & high-water mark & 過去の最高値を超えた分にだけ成功報酬をとる決まり & \ref{ch:hf} \\
ロックアップ・ゲート & lock-up / gate & 解約できない期間と、1 回に解約できる額の上限 & \ref{ch:hf} \\
マルチマネージャー & multi-manager platform & 多数の運用チームに資金を割り当て、中央でリスクを管理する運用会社 & \ref{ch:hf} \\
グロス・ネット & gross / net exposure & 買いと売りの金額の和と差（自己資本に対する比） & \ref{ch:pc} \\
\multicolumn{4}{@{}l}{\textbf{成績の物差し}} \\
シャープレシオ & Sharpe ratio & 超過リターンの平均 ÷ その標準偏差 & \ref{ch:perf} \\
アルファ・ベータ & alpha / beta & ファクターで説明できない平均の超過リターンと、ファクターへの感応度 & \ref{ch:perf} \\
トラッキングエラー & tracking error (TE) & 比べる相手との差のリターンの標準偏差 & \ref{ch:perf} \\
情報レシオ & information ratio (IR) & 比べる相手との差の平均 ÷ TE & \ref{ch:perf} \\
情報係数 & information coefficient (IC) & 予測と実際のリターンの（順位）相関 & \ref{ch:perf}, \ref{ch:signal} \\
基本法則 & fundamental law & $\IR \approx \mathrm{TC} \times \IC \times \sqrt{\mathrm{BR}}$ & \ref{ch:perf} \\
移転係数 & transfer coefficient (TC) & 制約のため、予測を比率に移しきれない度合い（0〜1） & \ref{ch:perf}, \ref{ch:pc} \\
最大ドローダウン & maximum drawdown & 過去の最高値からの最大の下落率 & \ref{ch:perf} \\
VaR・期待ショートフォール & value at risk / expected shortfall & 一定の確率で超えない損失と、それを超えた損失の平均 & \ref{ch:risk} \\
\multicolumn{4}{@{}l}{\textbf{ファクターとシグナル}} \\
ファクター & factor & 多くの資産のリターンに共通する要因。上位と下位のロング・ショートで作る & \ref{ch:factor} \\
アノマリー & anomaly & 既知のリスクで説明できない平均リターンの規則性 & \ref{ch:factor}, \ref{ch:equityfactor} \\
モメンタム & momentum & 過去に上がったものがさらに上がる傾向（横断面・時系列） & \ref{ch:equityfactor}, \ref{ch:trend} \\
割安 & value & 簿価・利益に比べて価格の低いものが上回る傾向 & \ref{ch:equityfactor} \\
キャリー & carry & 価格が変わらなかったときに得られる収益 & \ref{ch:macro} \\
低ボラティリティ・BAB & low risk / betting against beta & 低ベータ・低変動の資産がリスクの割に高いリターンを持つ傾向 & \ref{ch:equityfactor} \\
決算発表の後のドリフト & post-earnings-announcement drift (PEAD) & 決算の驚きの方向に、発表の後も株価が動き続ける傾向 & \ref{ch:event} \\
SUE & standardized unexpected earnings & 利益の驚きを、その変動の大きさで割ったもの & \ref{ch:signal}, \ref{ch:event} \\
中立化 & neutralization & シグナルから業種・規模などの影響を回帰で除くこと & \ref{ch:signal} \\
減衰 & signal decay & シグナルの情報が時間とともに価格に入り、予測力がなくなること & \ref{ch:signal} \\
\multicolumn{4}{@{}l}{\textbf{ポートフォリオとリスク}} \\
平均分散最適化 & mean-variance optimization & 期待リターンと分散のトレードオフを最適化する方法 & \ref{ch:pc} \\
Black--Litterman & Black--Litterman model & 時価総額加重から逆算した均衡のリターンに、見通しをベイズで混ぜる方法 & \ref{ch:pc} \\
リスクパリティ & risk parity & 各資産のリスクへの寄与をそろえる方法 & \ref{ch:pc} \\
階層的リスクパリティ & hierarchical risk parity (HRP) & クラスタリングで木を作り、逆行列を使わずにリスクを配る方法 & \ref{ch:pc}, \ref{ch:pcresearch} \\
ケリー基準 & Kelly criterion & 資産の対数の期待値を最大にする賭けの割合 & \ref{ch:pc} \\
ボラティリティ・ターゲット & volatility targeting & 予測した変動率が一定になるように量を調整すること & \ref{ch:pc}, \ref{ch:pcresearch} \\
縮小推定 & shrinkage estimator & 標本の推定値を構造のある目標に向けて縮め、誤差を減らす推定 & \ref{ch:risk} \\
ファクター型のリスクモデル & factor risk model & 共分散を少数のファクターと固有の分散で近似するモデル & \ref{ch:risk} \\
GARCH・実現ボラティリティ & GARCH / realized volatility & 変動率の予測のモデルと、日中のデータで測った変動率 & \ref{ch:risk} \\
\multicolumn{4}{@{}l}{\textbf{取引と市場の仕組み}} \\
スプレッド & bid-ask spread & 最良の売値と買値の差 & \ref{ch:cost}, \ref{ch:mm} \\
マーケット・インパクト & market impact & 自分の売買が価格を不利に動かす分 & \ref{ch:cost} \\
平方根則 & square-root law & インパクトが（注文 ÷ 売買高）の平方根に比例する経験則 & \ref{ch:cost}, \ref{ch:pcresearch} \\
インプリメンテーション・ショートフォール & implementation shortfall & 判断の時点の価格で売買できた場合との成績の差 & \ref{ch:cost} \\
容量 & capacity & コストの後の利益が消えるまでに運用できる資金の大きさ & \ref{ch:cost}, \ref{ch:pcresearch} \\
板 & limit order book & 指値の注文が価格ごとに並んだもの & \ref{ch:mm} \\
逆選択 & adverse selection & 情報を持った相手と取引して損をすること & \ref{ch:mm} \\
高頻度取引 & high-frequency trading (HFT) & 秒以下の速さで注文を出し取り消す取引 & \ref{ch:mm}, \ref{ch:hft} \\
\multicolumn{4}{@{}l}{\textbf{戦略}} \\
統計的裁定・ペアトレード & statistical arbitrage / pairs trading & 似た資産の一時的なずれが戻ることに賭ける & \ref{ch:statarb} \\
共和分 & cointegration & 個々は非定常な系列の線形結合が定常になる関係 & \ref{ch:statarb} \\
トレンドフォロー & trend following, time-series momentum & 資産自身の過去のリターンの符号で売買する & \ref{ch:trend} \\
合併裁定 & merger arbitrage & 買収の発表の後、成立までの価格差をとる & \ref{ch:event} \\
ベーシス取引 & basis trade & 現物と先物の価格差をとる（国債など） & \ref{ch:rv} \\
転換社債の裁定 & convertible arbitrage & 転換社債を買い、その株を空売りする & \ref{ch:rv} \\
ボラティリティ・リスク・プレミアム & volatility risk premium & オプションが示す変動率が、実現の変動率より平均して高いこと & \ref{ch:rv} \\
\multicolumn{4}{@{}l}{\textbf{検証}} \\
先読み & look-ahead bias & その時点で知られていなかった情報を使うこと & \ref{ch:backtest}, \ref{ch:text} \\
生存者バイアス & survivorship bias & 生き残ったものだけで評価して成績が良く見えること & \ref{ch:backtest}, \ref{ch:hfperf} \\
遡って入る偏り & backfill / instant-history bias & 成績が良かったときだけ過去の成績と一緒にデータに加わる偏り & \ref{ch:hfperf} \\
デフレートしたシャープレシオ & deflated Sharpe ratio (DSR) & 試した数・歪度・尖度を考慮して補正したシャープレシオ & \ref{ch:backtest} \\
バックテストの過学習の確率 & probability of backtest overfitting (PBO) & 学習の期間で最良のものが評価の期間で中央値を下回る確率 & \ref{ch:backtest} \\
パージ・エンバーゴ & purging / embargo & 学習と評価の正解の期間の重なりを除き、隙間を空けること & \ref{ch:model}, \ref{ch:backtest} \\
事前登録 & pre-registration & 仮説と方法を結果を見る前に固定すること & \ref{ch:backtest} \\
\multicolumn{4}{@{}l}{\textbf{失敗}} \\
混雑 & crowding & 多くの運用者が同じ取引をしていること & \ref{ch:failures} \\
解消のらせん & liquidation spiral & 損失 → 担保の追加 → 投げ売り → 価格の下落 → 損失の連鎖 & \ref{ch:risk}, \ref{ch:failures} \\
\end{longtable}
}

\chapter{文献と数値の照合}
\label{ch:verify}

\section{文献の実在の照合}

報告書で引いた文献の一覧（\code{照合/文献の一覧.md}）を、\code{LLM/refcheck.py} で照合した。結果は \code{照合/照合の結果.txt} にある。

（照合の集計は、照合が終わってから作る。）


一致が弱かった 20 本は、著者名を加えて Crossref で照合し直し（上の表の右の列）、書籍と業界誌の記事を含めて、すべての実在を確かめた。照合の後に一覧に加えた Hafsi・Vittori は、Crossref の DOI の記録で確かめた（\ref{sec:othermodel} 節）。

\section{Deep Research の答えの誤りと、確かめられなかった記述}

第 II 部を書くときに見つけた、Deep Research の答え（\code{03\_方針と文献調査/ヘッジファンドの手法/deep/}）の誤りと、原典で確かめられなかった記述を並べる。「確かめた」は、論文の要旨（Crossref・arXiv・NBER のページ）または原典の PDF で確かめたことを指す。

{\small
\begin{longtable}{@{}cP{0.3\linewidth}P{0.32\linewidth}P{0.24\linewidth}@{}}
\toprule
答え & Deep Research の記述 & 照合の結果 & 報告書での扱い \\
\midrule
\endfirsthead
\toprule
答え & Deep Research の記述 & 照合の結果 & 報告書での扱い \\
\midrule
\endhead
\bottomrule
\endfoot
\multicolumn{4}{@{}l}{\textbf{誤り}} \\
1 & Khandani・Lo（2007, 2012）によれば、2007 年 8 月に「標準偏差の 25 倍」の変動が数日続いた & 「25 標準偏差」は、当時の Goldman Sachs の幹部の発言として知られるもので、論文の結論ではない & 逸話として、出所を正して書いた（第 \ref{ch:statarb} 章） \\
1 & Fama・French（2015）の 5 ファクターが「リターンの分散の 95\% を説明」 & 要旨（SSRN の版）にこの値はない & 使わない \\
1 & Menkveld（2013）：1 銘柄 1 日 1,397 回の取引、1 回 0.88 ユーロ、シャープレシオ 9.35 & 要旨にない。要旨は「シャープレシオは資本のコストの仮定に強く依存する」と書く & 使わない。要旨の値（取引への参加率 8.1\%・64.4\%、約 5 分の 4 が受け身）を書いた \\
2 & Lopez-Lira・Tang：GPT-4 が「約 90\% の精度」 & 要旨では、取引できない最初の反応についての、ポートフォリオの日の的中率 & 区別を明記して書いた（第 \ref{ch:text} 章） \\
2 & 文献の一覧の番号と URL & 番号が本文の引用と対応せず、URL の多くが解説のブログ・Wikipedia・別の論文を指す & 一覧は使わず、題名で照合し直した \\
3 & Yang ほか（2022）：データの期間の中央が 2015 年以降の研究では平均のアルファが負 & 2022 年 7 月の版の要旨と本文の検索で見つからない（平均の月 0.36\%、74 研究・1,019 推定値は一致） & 一致した値だけを書いた \\
3 & Jacobs・Müller（2020）：231 のアノマリー & 要旨は 241 & 241 と書いた \\
3 & Sun・Wang・Zheng（2016）の題名「Only Winners in Hard Times Sleep Well: \ldots Market Market Environments」 & 正しくは「Only Winners in Tough Times Repeat: Hedge Fund Performance Persistence over Different Market Conditions」 & 正しい題名で引いた \\
3 & Jorion（2000）・Pástor ほか（2020）の URL & それぞれ Amaranth の論文、別の論文（量的・裁量の運用の規模の不経済）を指す & 題名で照合し直した \\
3 & Ben-David ほか（2023）NBER & NBER の WP は 2020 年（w27454） & 2020 年として引いた \\
\multicolumn{4}{@{}l}{\textbf{原典と一致した数値}} \\
1 & Greenwood・Sammon：指数の効果は 1990 年代 7.6\% → 最近 0.8\% & NBER の要旨と一致（1980 年代は 3.4\%） & 書いた（第 \ref{ch:event} 章） \\
1 & Womack（1996）：買い +2.4\%、売り $-9.1\%$ & 要旨と一致 & 書いた \\
1 & Mitchell・Pulvino（2001）：合併裁定の超過リターン年 4\%、プットの売りに似る & 要旨と一致 & 書いた \\
2 & Korajczyk・Sadka（2004）：容量 約 50 億ドル & 要旨（1999 年 12 月の時価総額の水準）と一致 & 書いた（第 \ref{ch:pcresearch} 章） \\
2 & Hafsi・Vittori の DQN による執行 & 題名・著者が Crossref で一致（2026 年の会議録） & 内容は答えの要約として書いた \\
3 & Ben-David ほか：実効の成功報酬 約 50\%、投資家の取り分 36 セント & NBER の要旨と一致 & 書いた（第 \ref{ch:hfperf} 章） \\
3 & McLean・Pontiff（2016）：標本外で 26\%、公表の後に 58\% 低下 & 分野でよく知られた値（要旨） & 書いた \\
3 & Yang ほか（2022）：74 研究、1,019 推定値、平均 月 0.36\% & 原典の PDF と一致 & 書いた \\
\multicolumn{4}{@{}l}{\textbf{確かめられなかった、または学術論文でない出典}} \\
1 & ペアトレードの年率 8\%・最大ドローダウン 12\%、機械学習のシグナルの減衰 3〜6 か月 & 出典は業界の統計のサイト（Gitnux） & 使わない \\
1 & 個人のファクター投資の超過リターン 年 3.27\% & 出典は修士論文 & 使わない \\
1 & 転換社債の裁定 年 9.40\%（1994〜2004 年） & 出典は大学の論文の抜粋で、著者が不明 & 使わない \\
1 & マクロのファンドのシャープレシオ 1.0・0.8・0.9 & 出典を確かめられない & 使わない \\
1 & Carr・Wu（2009）Table 5 の $-0.594$、売りの情報レシオ 3 超 & 原典で確かめていない & 数値は使わず、結論だけを書いた \\
1 & Bernard・Thomas（1989）Table 2：60 日で 4.2\% & 原典で確かめていない & 数値は使わない \\
2 & BloombergGPT の学習費用 約 267 万ドル、FinGPT 約 65 ドル & 原典で確かめていない & 使わない \\
3 & 2008 年のヘッジファンドの損失「総資産の 7\%、自己資本の 140\%」 & 出典を確かめられない & 使わない \\
3 & LTCM の想定の変動率 8.1\% → 13.6\%・19.2\% & 出典を確かめられない & 使わない \\
3 & ボルマゲドン：「ベータ約 4.0」「VIX 先物 約 93,000 枚」 & 出典は報道・解説 & 使わない \\
\end{longtable}
}

\section{照合の限界}

\begin{itemize}
\item \code{refcheck.py} の題名の照合は、文献の実在と、題名・著者・年の大筋を確かめるもので、本文の数値や主張が論文の内容と合っているかは確かめない。数値は、上の表のように個別に要旨や原典で確かめた。
\item 照合の一致の度合いが高くても、別の論文に一致していることがある（例：Harvey・Liu・Zhu（2016）の題名「\ldots and the Cross-Section of Expected Returns」は、題名が短く、別の論文に一致した）。そうしたものは自分で確かめた。
\item 年は、作業論文の版（SSRN・NBER）の年で一致することがある。報告書の文献の一覧には、主に学術誌に載った年を書いた。
\end{itemize}

\section{別のモデルによる点検}
\label{sec:othermodel}

書き上げた本文の全章を Gemini（gemini-3.1-pro-preview）に 1 回だけ送り、事実の誤り（著者・年・結果の取り違え、数値、定義、歴史）を探させた（\code{照合/gemini\_査読.md}）。指摘は 1 件だけで、「Hafsi・Vittori の論文は実在しない生成 AI の捏造だ」というものだった。しかし、Crossref の DOI（10.1109/acdsa67686.2026.11467851）の記録で、IEEE の 2026 年 2 月の会議録の論文（著者 Yadh Hafsi・Edoardo Vittori、9 ページ）として実在を確かめた。Gemini の知識より新しい文献を「架空」と判定した誤りと考えられる。LLM による照合は、新しい文献の実在の判定には使えない。


% 作成スクリプト/文献.py が 照合/文献の一覧.md から作る。手で直さない。
\begin{thebibliography}{999}
\small
\setlength{\itemsep}{0.15em}
\bibitem{afj2013} Agarwal, V., Fos, V., Jiang, W. (2013). ``Inferring Reporting-Related Biases in Hedge Fund Databases from Hedge Fund Equity Holdings.'' \textit{Management Science}.
\bibitem{aflm2011} Agarwal, V., Fung, W. H., Loon, Y. C., Naik, N. Y. (2011). ``Risk and Return in Convertible Arbitrage: Evidence from the Convertible Bond Market.'' \textit{Journal of Empirical Finance}.
\bibitem{agarwalnaik2004} Agarwal, V., Naik, N. Y. (2004). ``Risks and Portfolio Decisions Involving Hedge Funds.'' \textit{Review of Financial Studies}.
\bibitem{ace2013} Aiken, A. L., Clifford, C. P., Ellis, J. (2013). ``Out of the Dark: Hedge Fund Reporting Biases and Commercial Databases.'' \textit{Review of Financial Studies}.
\bibitem{akbas2017} Akbas, F., Markov, S., Subasi, M., Weisbrod, E. (2017). ``Determinants and Consequences of Information Processing Delay: Evidence from the Thomson Reuters Institutional Brokers' Estimate System.'' \textit{Journal of Financial Economics}.
\bibitem{almgrenchriss2000} Almgren, R., Chriss, N. (2000). ``Optimal Execution of Portfolio Transactions.'' \textit{Journal of Risk}.
\bibitem{athl2005} Almgren, R., Thum, C., Hauptmann, E., Li, H. (2005). ``Direct Estimation of Equity Market Impact.'' \textit{Risk}.
\bibitem{angbekaert2002} Ang, A., Bekaert, G. (2002). ``International Asset Allocation with Regime Shifts.'' \textit{Review of Financial Studies}.
\bibitem{ahxz2006} Ang, A., Hodrick, R. J., Xing, Y., Zhang, X. (2006). ``The Cross-Section of Volatility and Expected Returns.'' \textit{Journal of Finance}.
\bibitem{abo2022} Aquilina, M., Budish, E., O'Neill, P. (2022). ``Quantifying the High-Frequency Trading 'Arms Race'.'' \textit{Quarterly Journal of Economics}.
\bibitem{araci2019} Araci, D. (2019). ``FinBERT: Financial Sentiment Analysis with Pre-trained Language Models.'' \textit{arXiv:1908.10063}.
\bibitem{aragonstrahan2012} Aragon, G. O., Strahan, P. E. (2012). ``Hedge Funds as Liquidity Providers: Evidence from the Lehman Bankruptcy.'' \textit{Journal of Financial Economics}.
\bibitem{qmj2019} Asness, C. S., Frazzini, A., Pedersen, L. H. (2019). ``Quality Minus Junk.'' \textit{Review of Accounting Studies}.
\bibitem{amp2013} Asness, C. S., Moskowitz, T. J., Pedersen, L. H. (2013). ``Value and Momentum Everywhere.'' \textit{Journal of Finance}.
\bibitem{akl2001} Asness, C., Krail, R., Liew, J. (2001). ``Do Hedge Funds Hedge?.'' \textit{Journal of Portfolio Management}.
\bibitem{volmageddon2021} Augustin, P., Cheng, I.-H., Van den Bergen, L. (2021). ``Volmageddon and the Failure of Short Volatility Products.'' \textit{Financial Analysts Journal}.
\bibitem{avellanedalee2010} Avellaneda, M., Lee, J.-H. (2010). ``Statistical Arbitrage in the US Equities Market.'' \textit{Quantitative Finance}.
\bibitem{avellanedastoikov2008} Avellaneda, M., Stoikov, S. (2008). ``High-Frequency Trading in a Limit Order Book.'' \textit{Quantitative Finance}.
\bibitem{acm2023} Avramov, D., Cheng, S., Metzker, L. (2023). ``Machine Learning vs. Economic Restrictions: Evidence from Stock Return Predictability.'' \textit{Management Science}.
\bibitem{charlatanism2014} Bailey, D. H., Borwein, J. M., López de Prado, M., Zhu, Q. J. (2014). ``Pseudo-Mathematics and Financial Charlatanism: The Effects of Backtest Overfitting on Out-of-Sample Performance.'' \textit{Notices of the American Mathematical Society}.
\bibitem{pbo2017} Bailey, D. H., Borwein, J. M., López de Prado, M., Zhu, Q. J. (2017). ``The Probability of Backtest Overfitting.'' \textit{Journal of Computational Finance}.
\bibitem{dsr2014} Bailey, D. H., López de Prado, M. (2014). ``The Deflated Sharpe Ratio: Correcting for Selection Bias, Backtest Overfitting, and Non-Normality.'' \textit{Journal of Portfolio Management}.
\bibitem{bakersavasoglu2002} Baker, M., Savaşoglu, S. (2002). ``Limited Arbitrage in Mergers and Acquisitions.'' \textit{Journal of Financial Economics}.
\bibitem{bakshikapadia2003} Bakshi, G., Kapadia, N. (2003). ``Delta-Hedged Gains and the Negative Market Volatility Risk Premium.'' \textit{Review of Financial Studies}.
\bibitem{ballbrown1968} Ball, R., Brown, P. (1968). ``An Empirical Evaluation of Accounting Income Numbers.'' \textit{Journal of Accounting Research}.
\bibitem{banz1981} Banz, R. W. (1981). ``The Relationship between Return and Market Value of Common Stocks.'' \textit{Journal of Financial Economics}.
\bibitem{bbhk2019} Baron, M., Brogaard, J., Hagströmer, B., Kirilenko, A. (2019). ``Risk and Return in High-Frequency Trading.'' \textit{Journal of Financial and Quantitative Analysis}.
\bibitem{barroso2015} Barroso, P., Santa-Clara, P. (2015). ``Momentum Has Its Moments.'' \textit{Journal of Financial Economics}.
\bibitem{bbr2020} Ben-David, I., Birru, J., Rossi, A. (2020). ``The Performance of Hedge Fund Performance Fees.'' \textit{NBER Working Paper}.
\bibitem{bfm2012} Ben-David, I., Franzoni, F., Moussawi, R. (2012). ``Hedge Fund Stock Trading in the Financial Crisis of 2007–2009.'' \textit{Review of Financial Studies}.
\bibitem{bh1995} Benjamini, Y., Hochberg, Y. (1995). ``Controlling the False Discovery Rate: A Practical and Powerful Approach to Multiple Testing.'' \textit{Journal of the Royal Statistical Society Series B}.
\bibitem{berkgreen2004} Berk, J. B., Green, R. C. (2004). ``Mutual Fund Flows and Performance in Rational Markets.'' \textit{Journal of Political Economy}.
\bibitem{bernardthomas1989} Bernard, V. L., Thomas, J. K. (1989). ``Post-Earnings-Announcement Drift: Delayed Price Response or Risk Premium?.'' \textit{Journal of Accounting Research}.
\bibitem{bernardthomas1990} Bernard, V. L., Thomas, J. K. (1990). ``Evidence That Stock Prices Do Not Fully Reflect the Implications of Current Earnings for Future Earnings.'' \textit{Journal of Accounting and Economics}.
\bibitem{blacklitterman1992} Black, F., Litterman, R. (1992). ``Global Portfolio Optimization.'' \textit{Financial Analysts Journal}.
\bibitem{blitz2021} Blitz, D. (2021). ``The Quant Crisis of 2018–2020: Cornered by Big Growth.'' \textit{Journal of Portfolio Management}.
\bibitem{bollenpool2009} Bollen, N. P. B., Pool, V. K. (2009). ``Do Hedge Fund Managers Misreport Returns? Evidence from the Pooled Distribution.'' \textit{Journal of Finance}.
\bibitem{bollerslev1986} Bollerslev, T. (1986). ``Generalized Autoregressive Conditional Heteroskedasticity.'' \textit{Journal of Econometrics}.
\bibitem{bouchaud2018} Bouchaud, J.-P., Bonart, J., Donier, J., Gould, M. (2018). ``Trades, Quotes and Prices: Financial Markets Under the Microscope.'' \textit{Cambridge University Press}.
\bibitem{bjpt2008} Brav, A., Jiang, W., Partnoy, F., Thomas, R. (2008). ``Hedge Fund Activism, Corporate Governance, and Firm Performance.'' \textit{Journal of Finance}.
\bibitem{bhr2014} Brogaard, J., Hendershott, T., Riordan, R. (2014). ``High-Frequency Trading and Price Discovery.'' \textit{Review of Financial Studies}.
\bibitem{bgi1999} Brown, S. J., Goetzmann, W. N., Ibbotson, R. G. (1999). ``Offshore Hedge Funds: Survival and Performance, 1989–95.'' \textit{Journal of Business}.
\bibitem{bnp2008} Brunnermeier, M. K., Nagel, S., Pedersen, L. H. (2008). ``Carry Trades and Currency Crashes.'' \textit{NBER Macroeconomics Annual}.
\bibitem{brunnermeierpedersen2009} Brunnermeier, M. K., Pedersen, L. H. (2009). ``Market Liquidity and Funding Liquidity.'' \textit{Review of Financial Studies}.
\bibitem{bcs2015} Budish, E., Cramton, P., Shim, J. (2015). ``The High-Frequency Trading Arms Race: Frequent Batch Auctions as a Market Design Response.'' \textit{Quarterly Journal of Economics}.
\bibitem{deephedging2019} Buehler, H., Gonon, L., Teichmann, J., Wood, B. (2019). ``Deep Hedging.'' \textit{Quantitative Finance}.
\bibitem{campbellthompson2008} Campbell, J. Y., Thompson, S. B. (2008). ``Predicting Excess Stock Returns Out of Sample: Can Anything Beat the Historical Average?.'' \textit{Review of Financial Studies}.
\bibitem{carhart1997} Carhart, M. M. (1997). ``On Persistence in Mutual Fund Performance.'' \textit{Journal of Finance}.
\bibitem{carrwu2009} Carr, P., Wu, L. (2009). ``Variance Risk Premiums.'' \textit{Review of Financial Studies}.
\bibitem{cederburg2020} Cederburg, S., O'Doherty, M. S., Wang, F., Yan, X. S. (2020). ``On the Performance of Volatility-Managed Portfolios.'' \textit{Journal of Financial Economics}.
\bibitem{chenzimmermann2022} Chen, A. Y., Zimmermann, T. (2022). ``Open Source Cross-Sectional Asset Pricing.'' \textit{Critical Finance Review}.
\bibitem{cpz2024} Chen, L., Pelger, M., Zhu, J. (2024). ``Deep Learning in Asset Pricing.'' \textit{Management Science}.
\bibitem{ckx2022} Chen, Y., Kelly, B. T., Xiu, D. (2022). ``Expected Returns and Large Language Models.'' \textit{SSRN Working Paper}.
\bibitem{chincarini2007} Chincarini, L. B. (2007). ``The Amaranth Debacle: Failure of Risk Measures or Failure of Risk Management?.'' \textit{Journal of Alternative Investments}.
\bibitem{cgsss2009} Chordia, T., Goyal, A., Sadka, G., Sadka, R., Shivakumar, L. (2009). ``Liquidity and the Post-Earnings-Announcement Drift.'' \textit{Financial Analysts Journal}.
\bibitem{cst2014} Chordia, T., Subrahmanyam, A., Tong, Q. (2014). ``Have Capital Market Anomalies Attenuated in the Recent Era of High Liquidity and Trading Activity?.'' \textit{Journal of Accounting and Economics}.
\bibitem{cdt2002} Clarke, R., de Silva, H., Thorley, S. (2002). ``Portfolio Constraints and the Fundamental Law of Active Management.'' \textit{Financial Analysts Journal}.
\bibitem{cochrane2011} Cochrane, J. H. (2011). ``Presidential Address: Discount Rates.'' \textit{Journal of Finance}.
\bibitem{cmn2020} Cohen, L., Malloy, C., Nguyen, Q. (2020). ``Lazy Prices.'' \textit{Journal of Finance}.
\bibitem{cks2014} Cont, R., Kukanov, A., Stoikov, S. (2014). ``The Price Impact of Order Book Events.'' \textit{Journal of Financial Econometrics}.
\bibitem{corsi2009} Corsi, F. (2009). ``A Simple Approximate Long-Memory Model of Realized Volatility.'' \textit{Journal of Financial Econometrics}.
\bibitem{covalshumway2001} Coval, J. D., Shumway, T. (2001). ``Expected Option Returns.'' \textit{Journal of Finance}.
\bibitem{danielmoskowitz2016} Daniel, K., Moskowitz, T. J. (2016). ``Momentum Crashes.'' \textit{Journal of Financial Economics}.
\bibitem{dellavigna2009} DellaVigna, S., Pollet, J. M. (2009). ``Investor Inattention and Friday Earnings Announcements.'' \textit{Journal of Finance}.
\bibitem{dgu2009} DeMiguel, V., Garlappi, L., Uppal, R. (2009). ``Optimal Versus Naive Diversification: How Inefficient Is the 1/N Portfolio Strategy?.'' \textit{Review of Financial Studies}.
\bibitem{dmnr2020} DeMiguel, V., Martín-Utrera, A., Nogales, F. J., Uppal, R. (2020). ``A Transaction-Cost Perspective on the Multitude of Firm Characteristics.'' \textit{Review of Financial Studies}.
\bibitem{dichevyu2011} Dichev, I. D., Yu, G. (2011). ``Higher Risk, Lower Returns: What Hedge Fund Investors Really Earn.'' \textit{Journal of Financial Economics}.
\bibitem{dieboldmariano1995} Diebold, F. X., Mariano, R. S. (1995). ``Comparing Predictive Accuracy.'' \textit{Journal of Business and Economic Statistics}.
\bibitem{dofaff2010} Do, B., Faff, R. (2010). ``Does Simple Pairs Trading Still Work?.'' \textit{Financial Analysts Journal}.
\bibitem{dmv2009} Driessen, J., Maenhout, P. J., Vilkov, G. (2009). ``The Price of Correlation Risk: Evidence from Equity Options.'' \textit{Journal of Finance}.
\bibitem{dly2007} Duarte, J., Longstaff, F. A., Yu, F. (2007). ``Risk and Return in Fixed-Income Arbitrage: Nickels in Front of a Steamroller?.'' \textit{Review of Financial Studies}.
\bibitem{edwards1999} Edwards, F. R. (1999). ``Hedge Funds and the Collapse of Long-Term Capital Management.'' \textit{Journal of Economic Perspectives}.
\bibitem{engle1982} Engle, R. F. (1982). ``Autoregressive Conditional Heteroscedasticity with Estimates of the Variance of United Kingdom Inflation.'' \textit{Econometrica}.
\bibitem{englegranger1987} Engle, R. F., Granger, C. W. J. (1987). ``Co-integration and Error Correction: Representation, Estimation, and Testing.'' \textit{Econometrica}.
\bibitem{falck2022} Falck, A., Rej, A., Thesmar, D. (2021). ``Why and How Systematic Strategies Decay.'' \textit{arXiv:2105.01380}.
\bibitem{fama1970} Fama, E. F. (1970). ``Efficient Capital Markets: A Review of Theory and Empirical Work.'' \textit{Journal of Finance}.
\bibitem{fama1984} Fama, E. F. (1984). ``Forward and Spot Exchange Rates.'' \textit{Journal of Monetary Economics}.
\bibitem{ff1993} Fama, E. F., French, K. R. (1993). ``Common Risk Factors in the Returns on Stocks and Bonds.'' \textit{Journal of Financial Economics}.
\bibitem{ff2015} Fama, E. F., French, K. R. (2015). ``A Five-Factor Asset Pricing Model.'' \textit{Journal of Financial Economics}.
\bibitem{famamacbeth1973} Fama, E. F., MacBeth, J. D. (1973). ``Risk, Return, and Equilibrium: Empirical Tests.'' \textit{Journal of Political Economy}.
\bibitem{fjl2022} Frankel, R., Jennings, J., Lee, J. (2022). ``Disclosure Sentiment: Machine Learning vs. Dictionary Methods.'' \textit{Management Science}.
\bibitem{fim2018} Frazzini, A., Israel, R., Moskowitz, T. J. (2018). ``Trading Costs.'' \textit{SSRN Working Paper}.
\bibitem{bab2014} Frazzini, A., Pedersen, L. H. (2014). ``Betting Against Beta.'' \textit{Journal of Financial Economics}.
\bibitem{fnw2020} Freyberger, J., Neuhierl, A., Weber, M. (2020). ``Dissecting Characteristics Nonparametrically.'' \textit{Review of Financial Studies}.
\bibitem{funghsieh1997} Fung, W., Hsieh, D. A. (1997). ``Empirical Characteristics of Dynamic Trading Strategies: The Case of Hedge Funds.'' \textit{Review of Financial Studies}.
\bibitem{funghsieh2000} Fung, W., Hsieh, D. A. (2000). ``Performance Characteristics of Hedge Funds and Commodity Funds: Natural vs. Spurious Biases.'' \textit{Journal of Financial and Quantitative Analysis}.
\bibitem{funghsieh2001} Fung, W., Hsieh, D. A. (2001). ``The Risk in Hedge Fund Strategies: Theory and Evidence from Trend Followers.'' \textit{Review of Financial Studies}.
\bibitem{funghsieh2004} Fung, W., Hsieh, D. A. (2004). ``Hedge Fund Benchmarks: A Risk-Based Approach.'' \textit{Financial Analysts Journal}.
\bibitem{fhnr2008} Fung, W., Hsieh, D. A., Naik, N. Y., Ramadorai, T. (2008). ``Hedge Funds: Performance, Risk, and Capital Formation.'' \textit{Journal of Finance}.
\bibitem{ggr2006} Gatev, E., Goetzmann, W. N., Rouwenhorst, K. G. (2006). ``Pairs Trading: Performance of a Relative-Value Arbitrage Rule.'' \textit{Review of Financial Studies}.
\bibitem{getmansky2012} Getmansky, M. (2012). ``The Life Cycle of Hedge Funds: Fund Flows, Size, Competition, and Performance.'' \textit{Quarterly Journal of Finance}.
\bibitem{glm2004} Getmansky, M., Lo, A. W., Makarov, I. (2004). ``An Econometric Model of Serial Correlation and Illiquidity in Hedge Fund Returns.'' \textit{Journal of Financial Economics}.
\bibitem{grs1989} Gibbons, M. R., Ross, S. A., Shanken, J. (1989). ``A Test of the Efficiency of a Given Portfolio.'' \textit{Econometrica}.
\bibitem{glasserman2023} Glasserman, P., Lin, C. (2023). ``Assessing Look-Ahead Bias in Stock Return Predictions Generated by GPT Sentiment Analysis.'' \textit{arXiv:2309.17322}.
\bibitem{glostenmilgrom1985} Glosten, L. R., Milgrom, P. R. (1985). ``Bid, Ask and Transaction Prices in a Specialist Market with Heterogeneously Informed Traders.'' \textit{Journal of Financial Economics}.
\bibitem{greenwoodsammon} Greenwood, R., Sammon, M. (2022). ``The Disappearing Index Effect.'' \textit{NBER Working Paper}.
\bibitem{grinold1989} Grinold, R. C. (1989). ``The Fundamental Law of Active Management.'' \textit{Journal of Portfolio Management}.
\bibitem{grinoldkahn2000} Grinold, R. C., Kahn, R. N. (2000). ``Active Portfolio Management: A Quantitative Approach for Producing Superior Returns and Controlling Risk.'' \textit{McGraw-Hill}.
\bibitem{grossman1980} Grossman, S. J., Stiglitz, J. E. (1980). ``On the Impossibility of Informationally Efficient Markets.'' \textit{American Economic Review}.
\bibitem{gkx2020} Gu, S., Kelly, B., Xiu, D. (2020). ``Empirical Asset Pricing via Machine Learning.'' \textit{Review of Financial Studies}.
\bibitem{garleanupedersen2013} Gârleanu, N., Pedersen, L. H. (2013). ``Dynamic Trading with Predictable Returns and Transaction Costs.'' \textit{Journal of Finance}.
\bibitem{hafsivittori} Hafsi, Y., Vittori, E. (2026). ``Optimal Execution with Reinforcement Learning in a Multi-Agent Market Simulator.'' \textit{2026 International Conference on Artificial Intelligence, Computer, Data Sciences and Applications (ACDSA)}.
\bibitem{hamilton1989} Hamilton, J. D. (1989). ``A New Approach to the Economic Analysis of Nonstationary Time Series and the Business Cycle.'' \textit{Econometrica}.
\bibitem{hansen2005} Hansen, P. R. (2005). ``A Test for Superior Predictive Ability.'' \textit{Journal of Business and Economic Statistics}.
\bibitem{hhs2012} Hansen, P. R., Huang, Z., Shek, H. H. (2012). ``Realized GARCH: A Joint Model for Returns and Realized Measures of Volatility.'' \textit{Journal of Applied Econometrics}.
\bibitem{harrisgurel1986} Harris, L., Gurel, E. (1986). ``Price and Volume Effects Associated with Changes in the S\&P 500 List: New Evidence for the Existence of Price Pressures.'' \textit{Journal of Finance}.
\bibitem{hlz2016} Harvey, C. R., Liu, Y., Zhu, H. (2016). ``... and the Cross-Section of Expected Returns.'' \textit{Review of Financial Studies}.
\bibitem{hasanhodziclo2007} Hasanhodzic, J., Lo, A. W. (2007). ``Can Hedge-Fund Returns Be Replicated?: The Linear Case.'' \textit{Journal of Investment Management}.
\bibitem{helitterman1999} He, G., Litterman, R. (1999). ``The Intuition Behind Black-Litterman Model Portfolios.'' \textit{Goldman Sachs Investment Management Research}.
\bibitem{chronogpt2025} He, S., Lv, L., Manela, A., Wu, J. (2025). ``Chronologically Consistent Large Language Models.'' \textit{arXiv:2502.21206}.
\bibitem{hjm2011} Hendershott, T., Jones, C. M., Menkveld, A. J. (2011). ``Does Algorithmic Trading Improve Liquidity?.'' \textit{Journal of Finance}.
\bibitem{hlt2009} Hirshleifer, D., Lim, S. S., Teoh, S. H. (2009). ``Driven to Distraction: Extraneous Events and Underreaction to Earnings News.'' \textit{Journal of Finance}.
\bibitem{hostoll1981} Ho, T., Stoll, H. R. (1981). ``Optimal Dealer Pricing under Transactions and Return Uncertainty.'' \textit{Journal of Financial Economics}.
\bibitem{holm1979} Holm, S. (1979). ``A Simple Sequentially Rejective Multiple Test Procedure.'' \textit{Scandinavian Journal of Statistics}.
\bibitem{hls2000} Hong, H., Lim, T., Stein, J. C. (2000). ``Bad News Travels Slowly: Size, Analyst Coverage, and the Profitability of Momentum Strategies.'' \textit{Journal of Finance}.
\bibitem{hongstein1999} Hong, H., Stein, J. C. (1999). ``A Unified Theory of Underreaction, Momentum Trading, and Overreaction in Asset Markets.'' \textit{Journal of Finance}.
\bibitem{hxz2015} Hou, K., Xue, C., Zhang, L. (2015). ``Digesting Anomalies: An Investment Approach.'' \textit{Review of Financial Studies}.
\bibitem{hxz2020} Hou, K., Xue, C., Zhang, L. (2020). ``Replicating Anomalies.'' \textit{Review of Financial Studies}.
\bibitem{hwy2023} Huang, A. H., Wang, H., Yang, Y. (2023). ``FinBERT: A Large Language Model for Extracting Information from Financial Text.'' \textit{Contemporary Accounting Research}.
\bibitem{hlwz2020} Huang, D., Li, J., Wang, L., Zhou, G. (2020). ``Time Series Momentum: Is It There?.'' \textit{Journal of Financial Economics}.
\bibitem{hop2017} Hurst, B., Ooi, Y. H., Pedersen, L. H. (2017). ``A Century of Evidence on Trend-Following Investing.'' \textit{Journal of Portfolio Management}.
\bibitem{ilmanen2011} Ilmanen, A. (2011). ``Expected Returns: An Investor's Guide to Harvesting Market Rewards.'' \textit{Wiley}.
\bibitem{ikm2020} Israel, R., Kelly, B. T., Moskowitz, T. J. (2020). ``Can Machines 'Learn' Finance?.'' \textit{Journal of Investment Management}.
\bibitem{israelmoskowitz2013} Israel, R., Moskowitz, T. J. (2013). ``The Role of Shorting, Firm Size, and Time on Market Anomalies.'' \textit{Journal of Financial Economics}.
\bibitem{jacobsmuller2020} Jacobs, H., Müller, S. (2020). ``Anomalies across the Globe: Once Public, No Longer Existent?.'' \textit{Journal of Financial Economics}.
\bibitem{jagannathanma2003} Jagannathan, R., Ma, T. (2003). ``Risk Reduction in Large Portfolios: Why Imposing the Wrong Constraints Helps.'' \textit{Journal of Finance}.
\bibitem{jmn2010} Jagannathan, R., Malakhov, A., Novikov, D. (2010). ``Do Hot Hands Exist among Hedge Fund Managers? An Empirical Evaluation.'' \textit{Journal of Finance}.
\bibitem{jegadeesh1990} Jegadeesh, N. (1990). ``Evidence of Predictable Behavior of Security Returns.'' \textit{Journal of Finance}.
\bibitem{jt1993} Jegadeesh, N., Titman, S. (1993). ``Returns to Buying Winners and Selling Losers: Implications for Stock Market Efficiency.'' \textit{Journal of Finance}.
\bibitem{jkp2023} Jensen, T. I., Kelly, B., Pedersen, L. H. (2023). ``Is There a Replication Crisis in Finance?.'' \textit{Journal of Finance}.
\bibitem{jkx2023} Jiang, J., Kelly, B., Xiu, D. (2023). ``(Re-)Imag(in)ing Price Trends.'' \textit{Journal of Finance}.
\bibitem{jlw2012} Jiang, W., Li, K., Wang, W. (2012). ``Hedge Funds and Chapter 11.'' \textit{Journal of Finance}.
\bibitem{jxl2017} Jiang, Z., Xu, D., Liang, J. (2017). ``A Deep Reinforcement Learning Framework for the Financial Portfolio Management Problem.'' \textit{arXiv:1706.10059}.
\bibitem{jorion2000} Jorion, P. (2000). ``Risk Management Lessons from Long-Term Capital Management.'' \textit{European Financial Management}.
\bibitem{kkx2019} Ke, Z. T., Kelly, B. T., Xiu, D. (2019). ``Predicting Returns with Text Data.'' \textit{NBER Working Paper}.
\bibitem{kps2019} Kelly, B. T., Pruitt, S., Su, Y. (2019). ``Characteristics Are Covariances: A Unified Model of Risk and Return.'' \textit{Journal of Financial Economics}.
\bibitem{kmz2024} Kelly, B., Malamud, S., Zhou, K. (2024). ``The Virtue of Complexity in Return Prediction.'' \textit{Journal of Finance}.
\bibitem{kelly1956} Kelly, J. L. (1956). ``A New Interpretation of Information Rate.'' \textit{Bell System Technical Journal}.
\bibitem{khandanilo2011} Khandani, A. E., Lo, A. W. (2011). ``What Happened to the Quants in August 2007? Evidence from Factors and Transactions Data.'' \textit{Journal of Financial Markets}.
\bibitem{kmn2024} Kim, A. G., Muhn, M., Nikolaev, V. V. (2024). ``Financial Statement Analysis with Large Language Models.'' \textit{arXiv:2407.17866}.
\bibitem{kimtsewald2016} Kim, A. Y., Tse, Y., Wald, J. K. (2016). ``Time Series Momentum and Volatility Scaling.'' \textit{Journal of Financial Markets}.
\bibitem{kkst2017} Kirilenko, A., Kyle, A. S., Samadi, M., Tuzun, T. (2017). ``The Flash Crash: High-Frequency Trading in an Electronic Market.'' \textit{Journal of Finance}.
\bibitem{kmpv2018} Koijen, R. S. J., Moskowitz, T. J., Pedersen, L. H., Vrugt, E. B. (2018). ``Carry.'' \textit{Journal of Financial Economics}.
\bibitem{korajczyksadka2004} Korajczyk, R. A., Sadka, R. (2004). ``Are Momentum Profits Robust to Trading Costs?.'' \textit{Journal of Finance}.
\bibitem{knt2007} Kosowski, R., Naik, N. Y., Teo, M. (2007). ``Do Hedge Funds Deliver Alpha? A Bayesian and Bootstrap Analysis.'' \textit{Journal of Financial Economics}.
\bibitem{kns2020} Kozak, S., Nagel, S., Santosh, S. (2020). ``Shrinking the Cross-Section.'' \textit{Journal of Financial Economics}.
\bibitem{kruttli2021} Kruttli, M. S., Monin, P. J., Petrasek, L., Watugala, S. W. (2021). ``Hedge Fund Treasury Trading and Funding Fragility: Evidence from the COVID-19 Crisis.'' \textit{Finance and Economics Discussion Series, Federal Reserve Board}.
\bibitem{kyle1985} Kyle, A. S. (1985). ``Continuous Auctions and Insider Trading.'' \textit{Econometrica}.
\bibitem{lsv1994} Lakonishok, J., Shleifer, A., Vishny, R. W. (1994). ``Contrarian Investment, Extrapolation, and Risk.'' \textit{Journal of Finance}.
\bibitem{laloux1999} Laloux, L., Cizeau, P., Bouchaud, J.-P., Potters, M. (1999). ``Noise Dressing of Financial Correlation Matrices.'' \textit{Physical Review Letters}.
\bibitem{ledoitwolf2004} Ledoit, O., Wolf, M. (2004). ``Honey, I Shrunk the Sample Covariance Matrix.'' \textit{Journal of Portfolio Management}.
\bibitem{ledoitwolf2020} Ledoit, O., Wolf, M. (2020). ``Analytical Nonlinear Shrinkage of Large-Dimensional Covariance Matrices.'' \textit{Annals of Statistics}.
\bibitem{lehmann1990} Lehmann, B. N. (1990). ``Fads, Martingales, and Market Efficiency.'' \textit{Quarterly Journal of Economics}.
\bibitem{lwz2022} Leippold, M., Wang, Q., Zhou, W. (2022). ``Machine Learning in the Chinese Stock Market.'' \textit{Journal of Financial Economics}.
\bibitem{liang2000} Liang, B. (2000). ``Hedge Funds: The Living and the Dead.'' \textit{Journal of Financial and Quantitative Analysis}.
\bibitem{lintner1965} Lintner, J. (1965). ``The Valuation of Risk Assets and the Selection of Risky Investments in Stock Portfolios and Capital Budgets.'' \textit{Review of Economics and Statistics}.
\bibitem{finrl2020} Liu, X.-Y., Yang, H., Chen, Q., Zhang, R., Yang, L., Xiao, B., Wang, C. D. (2020). ``FinRL: A Deep Reinforcement Learning Library for Automated Stock Trading in Quantitative Finance.'' \textit{arXiv:2011.09607}.
\bibitem{lo2002} Lo, A. W. (2002). ``The Statistics of Sharpe Ratios.'' \textit{Financial Analysts Journal}.
\bibitem{lopezliratang2023} Lopez-Lira, A., Tang, Y. (2023). ``Can ChatGPT Forecast Stock Price Movements? Return Predictability and Large Language Models.'' \textit{arXiv:2304.07619}.
\bibitem{loupolk2022} Lou, D., Polk, C. (2022). ``Comomentum: Inferring Arbitrage Activity from Return Correlations.'' \textit{Review of Financial Studies}.
\bibitem{lps2019} Lou, D., Polk, C., Skouras, S. (2019). ``A Tug of War: Overnight versus Intraday Expected Returns.'' \textit{Journal of Financial Economics}.
\bibitem{lm2011} Loughran, T., McDonald, B. (2011). ``When Is a Liability Not a Liability? Textual Analysis, Dictionaries, and 10-Ks.'' \textit{Journal of Finance}.
\bibitem{lrv2011} Lustig, H., Roussanov, N., Verdelhan, A. (2011). ``Common Risk Factors in Currency Markets.'' \textit{Review of Financial Studies}.
\bibitem{hrp2016} López de Prado, M. (2016). ``Building Diversified Portfolios that Outperform Out of Sample.'' \textit{Journal of Portfolio Management}.
\bibitem{lopezdeprado2018} López de Prado, M. (2018). ``Advances in Financial Machine Learning.'' \textit{Wiley}.
\bibitem{mrt2010} Maillard, S., Roncalli, T., Teïletche, J. (2010). ``The Properties of Equally Weighted Risk Contribution Portfolios.'' \textit{Journal of Portfolio Management}.
\bibitem{malkielsaha2005} Malkiel, B. G., Saha, A. (2005). ``Hedge Funds: Risk and Return.'' \textit{Financial Analysts Journal}.
\bibitem{markowitz1952} Markowitz, H. (1952). ``Portfolio Selection.'' \textit{Journal of Finance}.
\bibitem{martineau2022} Martineau, C. (2022). ``Rest in Peace Post-Earnings Announcement Drift.'' \textit{Critical Finance Review}.
\bibitem{mcleanpontiff2016} McLean, R. D., Pontiff, J. (2016). ``Does Academic Research Destroy Stock Return Predictability?.'' \textit{Journal of Finance}.
\bibitem{msss2012} Menkhoff, L., Sarno, L., Schmeling, M., Schrimpf, A. (2012). ``Carry Trades and Global Foreign Exchange Volatility.'' \textit{Journal of Finance}.
\bibitem{menkveld2013} Menkveld, A. J. (2013). ``High Frequency Trading and the New Market Makers.'' \textit{Journal of Financial Markets}.
\bibitem{peadtxt} Meursault, V., Liang, P. J., Routledge, B. R., Scanlon, M. M. (2023). ``PEAD.txt: Post-Earnings-Announcement Drift Using Text.'' \textit{Journal of Financial and Quantitative Analysis}.
\bibitem{michaud1989} Michaud, R. O. (1989). ``The Markowitz Optimization Enigma: Is 'Optimized' Optimal?.'' \textit{Financial Analysts Journal}.
\bibitem{mpp2007} Mitchell, M., Pedersen, L. H., Pulvino, T. (2007). ``Slow Moving Capital.'' \textit{American Economic Review}.
\bibitem{mitchellpulvino2001} Mitchell, M., Pulvino, T. (2001). ``Characteristics of Risk and Return in Risk Arbitrage.'' \textit{Journal of Finance}.
\bibitem{mitchellpulvino2012} Mitchell, M., Pulvino, T. (2012). ``Arbitrage Crashes and the Speed of Capital.'' \textit{Journal of Financial Economics}.
\bibitem{moreiramuir2017} Moreira, A., Muir, T. (2017). ``Volatility-Managed Portfolios.'' \textit{Journal of Finance}.
\bibitem{mop2012} Moskowitz, T. J., Ooi, Y. H., Pedersen, L. H. (2012). ``Time Series Momentum.'' \textit{Journal of Financial Economics}.
\bibitem{nagel2012} Nagel, S. (2012). ``Evaporating Liquidity.'' \textit{Review of Financial Studies}.
\bibitem{nfk2006} Nevmyvaka, Y., Feng, Y., Kearns, M. (2006). ``Reinforcement Learning for Optimized Trade Execution.'' \textit{Proceedings of the 23rd International Conference on Machine Learning}.
\bibitem{nlj2021} Ning, B., Lin, F. H. T., Jaimungal, S. (2021). ``Double Deep Q-Learning for Optimal Execution.'' \textit{Applied Mathematical Finance}.
\bibitem{novymarx2013} Novy-Marx, R. (2013). ``The Other Side of Value: The Gross Profitability Premium.'' \textit{Journal of Financial Economics}.
\bibitem{novymarxvelikov2016} Novy-Marx, R., Velikov, M. (2016). ``A Taxonomy of Anomalies and Their Trading Costs.'' \textit{Review of Financial Studies}.
\bibitem{pedersen2009} Pedersen, L. H. (2009). ``When Everyone Runs for the Exit.'' \textit{International Journal of Central Banking}.
\bibitem{pedersen2015} Pedersen, L. H. (2015). ``Efficiently Inefficient: How Smart Money Invests and Market Prices Are Determined.'' \textit{Princeton University Press}.
\bibitem{perold1988} Perold, A. F. (1988). ``The Implementation Shortfall: Paper versus Reality.'' \textit{Journal of Portfolio Management}.
\bibitem{plerou1999} Plerou, V., Gopikrishnan, P., Rosenow, B., Amaral, L. A. N., Stanley, H. E. (1999). ``Universal and Nonuniversal Properties of Cross Correlations in Financial Time Series.'' \textit{Physical Review Letters}.
\bibitem{pst2015} Pástor, Ľ., Stambaugh, R. F., Taylor, L. A. (2015). ``Scale and Skill in Active Management.'' \textit{Journal of Financial Economics}.
\bibitem{ross1976} Ross, S. A. (1976). ``The Arbitrage Theory of Capital Asset Pricing.'' \textit{Journal of Economic Theory}.
\bibitem{sharpe1964} Sharpe, W. F. (1964). ``Capital Asset Prices: A Theory of Market Equilibrium under Conditions of Risk.'' \textit{Journal of Finance}.
\bibitem{sharpe1966} Sharpe, W. F. (1966). ``Mutual Fund Performance.'' \textit{Journal of Business}.
\bibitem{shleifer1986} Shleifer, A. (1986). ``Do Demand Curves for Stocks Slope Down?.'' \textit{Journal of Finance}.
\bibitem{shleifer1997} Shleifer, A., Vishny, R. W. (1997). ``The Limits of Arbitrage.'' \textit{Journal of Finance}.
\bibitem{sloan1996} Sloan, R. G. (1996). ``Do Stock Prices Fully Reflect Information in Accruals and Cash Flows about Future Earnings?.'' \textit{Accounting Review}.
\bibitem{stein2009} Stein, J. C. (2009). ``Presidential Address: Sophisticated Investors and Market Efficiency.'' \textit{Journal of Finance}.
\bibitem{swz2018} Sun, Z., Wang, A. W., Zheng, L. (2018). ``Only Winners in Tough Times Repeat: Hedge Fund Performance Persistence over Different Market Conditions.'' \textit{Journal of Financial and Quantitative Analysis}.
\bibitem{tetlock2007} Tetlock, P. C. (2007). ``Giving Content to Investor Sentiment: The Role of Media in the Stock Market.'' \textit{Journal of Finance}.
\bibitem{tsm2008} Tetlock, P. C., Saar-Tsechansky, M., Macskassy, S. (2008). ``More Than Words: Quantifying Language to Measure Firms' Fundamentals.'' \textit{Journal of Finance}.
\bibitem{till2007} Till, H. (2007). ``The Amaranth Collapse: What Happened and What Have We Learned Thus Far?.'' \textit{EDHEC Risk and Asset Management Research Centre}.
\bibitem{titmantiu2011} Titman, S., Tiu, C. (2011). ``Do the Best Hedge Funds Hedge?.'' \textit{Review of Financial Studies}.
\bibitem{toth2011} Tóth, B., Lempérière, Y., Deremble, C., de Lataillade, J., Kockelkoren, J., Bouchaud, J.-P. (2011). ``Anomalous Price Impact and the Critical Nature of Liquidity in Financial Markets.'' \textit{Physical Review X}.
\bibitem{white2000} White, H. (2000). ``A Reality Check for Data Snooping.'' \textit{Econometrica}.
\bibitem{womack1996} Womack, K. L. (1996). ``Do Brokerage Analysts' Recommendations Have Investment Value?.'' \textit{Journal of Finance}.
\bibitem{bloomberggpt2023} Wu, S., Irsoy, O., Lu, S., Dabravolski, V., Dredze, M., Gehrmann, S., Kambadur, P., Rosenberg, D., Mann, G. (2023). ``BloombergGPT: A Large Language Model for Finance.'' \textit{arXiv:2303.17564}.
\bibitem{yhin2022} Yang, F., Havranek, T., Irsova, Z., Novak, J. (2022). ``Hedge Fund Performance: A Quantitative Survey.'' \textit{IES Working Paper, Charles University}.
\bibitem{fingpt2023} Yang, H., Liu, X.-Y., Wang, C. D. (2023). ``FinGPT: Open-Source Financial Large Language Models.'' \textit{arXiv:2306.06031}.
\bibitem{yin2016} Yin, C. (2016). ``The Optimal Size of Hedge Funds: Conflict between Investors and Fund Managers.'' \textit{Journal of Finance}.
\end{thebibliography}


\end{document}
